\chapter{Complexity --- Cook--Levin Hardness}
%%% generated by the blueprint pipeline from TCSlib.Complexity.CookLevin.Hardness
%%%%%%%%%%%%%%%%%%%%%%%%%%%%%%%%%%%%%%%%%%%%%%%%%%%%%%%%%%%%%%%%%%%%%%%%%%%%%%%
\section{Overview}

This module completes the Cook--Levin theorem of Arora--Barak (Theorem 2.10,
via Lemma 2.11 and Lemma 2.14): the languages $\mathrm{SAT}$ and $\mathrm{3SAT}$
of serialized satisfiable CNF formulas are $\mathrm{NP}$-complete. It shows that
hardness transfers along polynomial-time reductions, develops the tableau over an
oblivious verifier --- snapshot block encodings, the six ordered clause-group
families, and a large toolkit of native multi-tape Turing-machine gadgets for
counting, recording, searching, and emitting --- and assembles a polynomial-time
machine that emits the serialized tableau of any $\mathrm{NP}$ verifier. Together
with equisatisfiability of the tableau and the exact certificate relation, this
yields the $\mathrm{NP}$-hardness and $\mathrm{NP}$-completeness of
$\mathrm{SAT}$ and $\mathrm{3SAT}$.

%%%%%%%%%%%%%%%%%%%%%%%%%%%%%%%%%%%%%%%%%%%%%%%%%%%%%%%%%%%%%%%%%%%%%%%%%%%%%%%
\section{Declarations}

\begin{theorem}[Hardness transfers along reductions]
\lean{Complexity.NPHard.polyTimeReducible}
\label{Complexity.NPHard.polyTimeReducible}
\leanok
\uses{Complexity.NPHard, Complexity.PolyTimeReducible.trans}
\difficulty{1}
Let $L$ and $L'$ be languages over the Boolean alphabet. If every language in $\mathrm{NP}$ polynomial-time many-one reduces to $L$, and $L$ polynomial-time reduces to $L'$, then every language in $\mathrm{NP}$ polynomial-time reduces to $L'$.
\end{theorem}

\begin{definition}[Index of the last emission round]
\lean{Complexity.clLastRound}
\label{Complexity.clLastRound}
\leanok
For an input length $n$, a work-tape count $k$, and a time horizon $T$, the last round index of the Cook--Levin clause-emission schedule is $n + (k+3)T + k + 1$, one less than the total number of clause groups.
\end{definition}

\begin{definition}[The six ordered clause-group families]
\lean{Complexity.clGroups}
\label{Complexity.clGroups}
\leanok
Given an input length $n$, a tape count $k$, a horizon $T$, and six clause-family generators, the resulting list arranges the clause groups in a fixed order: $n$ pinning groups, one initial group, $T$ state groups, $T+1$ input groups, $k(T+1)$ work groups ordered by time and then by tape, and $T$ acceptance groups. An empty clause group still occupies its own slot.
\end{definition}

\begin{lemma}[Count of clause groups]
\lean{Complexity.clGroups_length}
\label{Complexity.clGroups_length}
\leanok
\uses{Complexity.clGroups, Complexity.clLastRound}
\difficulty{3}
For every input length $n$, tape count $k$, horizon $T$, and any choice of the six clause-family generators, the ordered list of clause groups has length exactly $n + (k+3)T + k + 2$, one more than the last round index; groups are counted, not individual clauses.
\end{lemma}

\begin{definition}[Serialized fragment of a clause group]
\lean{Complexity.clFragment}
\label{Complexity.clFragment}
\leanok
\uses{Std.Sat.CNF.serializeClause}
For a CNF formula over natural-number variables, its fragment is the bit string obtained by emitting, for each clause in order, a leading clause-marker bit followed by the serialization of that clause; no formula terminator is appended.
\end{definition}

\begin{lemma}[Fragments respect concatenation]
\lean{Complexity.clFragment_append}
\label{Complexity.clFragment_append}
\leanok
\uses{Complexity.clFragment}
\difficulty{1}
For CNF formulas $\varphi$ and $\psi$, the serialized fragment of $\varphi$ followed by $\psi$ equals the fragment of $\varphi$ concatenated with the fragment of $\psi$, with all clauses kept in their original order.
\end{lemma}

\begin{definition}[Round chunk of the serialized formula]
\lean{Complexity.clChunk}
\label{Complexity.clChunk}
\leanok
\uses{Complexity.clFragment}
Given a last round index $R$, a family of clause groups indexed by round, and a round $i$, the chunk at $i$ is the serialized fragment of the group at $i$, followed by the single formula-terminator bit exactly when $i = R$; an empty final group still emits the terminator.
\end{definition}

\begin{lemma}[Proper prefixes carry no terminator]
\lean{Complexity.clChunks_before}
\label{Complexity.clChunks_before}
\leanok
\uses{Complexity.clChunk, Complexity.clFragment, Complexity.clFragment_append}
\difficulty{4}
Let $R$ be the last round index and let a family assign a clause group to each round. For every $j \le R$, the concatenation of the chunks of rounds $0,\dots,j-1$ equals the serialized fragment of the concatenation of the first $j$ groups, so no formula terminator appears in a proper prefix of the rounds.
\end{lemma}

\begin{lemma}[Chunks serialize the concatenated formula]
\lean{Complexity.clChunks_serialize}
\label{Complexity.clChunks_serialize}
\leanok
\uses{Complexity.clChunk, Complexity.clChunks_before, Complexity.clFragment, Std.Sat.CNF.serialize}
\difficulty{3}
Let $R$ be the last round index and let a family assign a clause group to each round. The concatenation of the chunks of rounds $0,\dots,R$ equals the full serialization of the concatenation of all $R+1$ groups, formula terminator included. The identity concerns output words only and makes no satisfiability claim.
\end{lemma}

\begin{lemma}[Indexed lookup reconstructs a list]
\lean{Complexity.clGroups_index}
\label{Complexity.clGroups_index}
\leanok
\difficulty{3}
For every list $L$ of CNF formulas, the list obtained by looking up the entries of $L$ at indices $0,\dots,|L|-1$, with the empty formula as default for out-of-range lookups, is $L$ itself; the default is never used at a legal index.
\end{lemma}

\begin{lemma}[Serialization of the six-family list]
\lean{Complexity.clGroups_serialize}
\label{Complexity.clGroups_serialize}
\leanok
\uses{Complexity.clChunk, Complexity.clChunks_serialize, Complexity.clGroups, Complexity.clGroups_index, Complexity.clGroups_length, Complexity.clLastRound, Std.Sat.CNF.serialize}
\difficulty{3}
For all parameters $n$, $k$, $T$ and all six clause-family generators, concatenating the round chunks of the six-family group list over rounds $0,\dots,R$, where $R = n+(k+3)T+k+1$ is its computed last round, equals the serialization of the flattened six-family formula, without rearranging any clause.
\end{lemma}

\begin{definition}[Packing of snapshot wires]
\lean{Complexity.clPack}
\label{Complexity.clPack}
\leanok
Given a count $m$ of input variables, a block width $B$, a time $t$, and a bit position $i$, the packed global variable index is $m + tB + i$: snapshot bits occupy consecutive blocks of width $B$ placed after all $m$ input variables.
\end{definition}

\begin{lemma}[Oblivious quadratic-overhead verifier]
\lean{Complexity.clObliviousVerifier}
\label{Complexity.clObliviousVerifier}
\leanok
\uses{Complexity.DTIME, Complexity.P, Complexity.mem_P_iff, Complexity.oblivious_of_mem_DTIME, Complexity.timeConstructible_poly, Turing.FinTM, Turing.FinTM.ComputesInTime.mono, Turing.FinTM.DecidesInTime, Turing.FinTM.Oblivious, Turing.FinTM.not_computesInTime_zero}
\difficulty{4}
Let $V$ be a language in $\mathrm{P}$. Then there exist a finite machine $M$ and constants $c, A, d > 0$ such that $M$ is oblivious, meaning its head positions depend only on the input length and the time, and $M$ decides $V$ within time $c\,(A(n+1)^d + 1)^2$ on inputs of length $n$. Only the time bound is enlarged; the language, and hence the certificate relation, is unchanged.
\end{lemma}

\begin{lemma}[Polynomial budget for the emission loop]
\lean{Complexity.clLoop_polyBound}
\label{Complexity.clLoop_polyBound}
\leanok
\uses{Complexity.PolyBound}
\difficulty{4}
Let $P$ and $R$ be polynomially bounded functions on the naturals and let $c$ be a constant. Then the function $n \mapsto c\,(P(n)+1)(R(n)+2)$ is polynomially bounded.
\end{lemma}

\begin{definition}[Fixed CNF table for a Boolean predicate]
\lean{Complexity.clTemplate}
\label{Complexity.clTemplate}
\leanok
\uses{Complexity.exists_cnf_boolFun}
For an arity $\ell$ and a Boolean predicate $p$ on $\ell$-bit vectors, a CNF formula computing $p$ is chosen once and for all; the choice depends only on $\ell$ and $p$, before any input or index wiring is supplied.
\end{definition}

\begin{lemma}[Size and correctness of the chosen table]
\lean{Complexity.clTemplate_spec}
\label{Complexity.clTemplate_spec}
\leanok
\uses{Complexity.clTemplate, Complexity.exists_cnf_boolFun, Std.Sat.CNF.WidthAtMost, Std.Sat.CNF.numVars}
\difficulty{1}
For every arity $\ell$ and predicate $p$ on $\ell$-bit vectors, the chosen CNF table uses at most $\ell$ variables, has at most $2^\ell$ clauses, has width at most $\ell$, and evaluates under any assignment $a$ to exactly $p$ applied to the first $\ell$ bits of $a$.
\end{lemma}

\begin{lemma}[Evaluation of a relabeled table]
\lean{Complexity.clTemplate_eval}
\label{Complexity.clTemplate_eval}
\leanok
\uses{Complexity.clTemplate, Complexity.clTemplate_spec}
\difficulty{1}
For every arity $\ell$, predicate $p$ on $\ell$-bit vectors, wiring map $w$ on variable indices, and assignment $a$, the relabeling of the chosen CNF table along $w$ evaluates to $p$ applied to the bits $a(w(0)),\dots,a(w(\ell-1))$; repeated wire indices are permitted and no injectivity is assumed.
\end{lemma}

\begin{definition}[Binary code for a finite alphabet]
\lean{Complexity.CLFieldCode}
\label{Complexity.CLFieldCode}
\leanok
A binary code for a type consists of a bit width $B$, an encoding map from the type into $B$-bit vectors, and a total decoding map that is a left inverse of the encoding. The code is fixed once for the machine, before any reduction input.
\end{definition}

\begin{definition}[A fixed binary code always exists]
\lean{Complexity.clFieldCode}
\label{Complexity.clFieldCode}
\leanok
\uses{Complexity.CLFieldCode}
Every nonempty finite type carries a binary code whose width is the cardinality of the type: since a finite set injects into the Boolean vectors indexed by itself, an embedding exists, and any left inverse of it serves as the total decoder. Taking the cardinality as bit width avoids any logarithm convention.
\end{definition}

\begin{definition}[Two-bit encoding of an optional symbol]
\lean{Complexity.clSymbolCode}
\label{Complexity.clSymbolCode}
\leanok
An optional Boolean tape symbol is encoded in two bits: the first bit records whether a symbol is present, and the second records its value, defaulting to false for a blank.
\end{definition}

\begin{definition}[Two-bit symbol decoder]
\lean{Complexity.clSymbolDecode}
\label{Complexity.clSymbolDecode}
\leanok
A pair of bits decodes to an optional Boolean symbol: if the presence bit is set, the result is the value bit; otherwise the result is the blank symbol. The decoder is total on all four bit patterns.
\end{definition}

\begin{lemma}[Symbol decoding inverts encoding]
\lean{Complexity.clSymbolDecode_code}
\label{Complexity.clSymbolDecode_code}
\leanok
\uses{Complexity.clSymbolCode, Complexity.clSymbolDecode}
\difficulty{2}
For every optional Boolean symbol $s$, decoding the two-bit encoding of $s$ returns $s$, including the blank case and the false-valued case.
\end{lemma}

\begin{definition}[Product encoding of a snapshot]
\lean{Complexity.clBlockEncode}
\label{Complexity.clBlockEncode}
\leanok
\uses{Complexity.CLFieldCode, Complexity.Snapshot, Complexity.clSymbolCode, Turing.FinTM}
Given a machine with $k$ work tapes and a binary code of width $B$ for its optional states, a snapshot, consisting of a state, an input symbol, and one symbol per work tape, is encoded as a vector of $B + 2 + 2k$ bits: the state bits, then two input-symbol bits, then two bits per work tape, in literal consecutive slices.
\end{definition}

\begin{definition}[State slice of a block]
\lean{Complexity.clStateSlice}
\label{Complexity.clStateSlice}
\leanok
\uses{Complexity.CLFieldCode, Turing.FinTM}
Within a snapshot block of $B + 2 + 2k$ bits, where $B$ is the state-code width and $k$ the number of work tapes, the state slice reads the first $B$ bits.
\end{definition}

\begin{definition}[Input-symbol slice of a block]
\lean{Complexity.clInputSlice}
\label{Complexity.clInputSlice}
\leanok
\uses{Complexity.CLFieldCode, Turing.FinTM}
Within a snapshot block of $B + 2 + 2k$ bits, where $B$ is the state-code width and $k$ the number of work tapes, the input slice reads the two bits that follow the $B$ state bits.
\end{definition}

\begin{definition}[Work-symbol slice of a block]
\lean{Complexity.clWorkSlice}
\label{Complexity.clWorkSlice}
\leanok
\uses{Complexity.CLFieldCode, Turing.FinTM}
Within a snapshot block of $B + 2 + 2k$ bits, where $B$ is the state-code width and $k$ the number of work tapes, the work slice for tape $\tau$ reads the two bits at offset $B + 2 + 2\tau$.
\end{definition}

\begin{lemma}[State slice of an encoded snapshot]
\lean{Complexity.clBlock_state}
\label{Complexity.clBlock_state}
\leanok
\uses{Complexity.CLFieldCode, Complexity.Snapshot, Complexity.clBlockEncode, Complexity.clStateSlice, Turing.FinTM}
\difficulty{1}
For every snapshot $s$ of a machine with a fixed optional-state binary code, the state slice of the encoded block of $s$ equals the code's encoding of the state of $s$.
\end{lemma}

\begin{lemma}[Input slice of an encoded snapshot]
\lean{Complexity.clBlock_input}
\label{Complexity.clBlock_input}
\leanok
\uses{Complexity.CLFieldCode, Complexity.Snapshot, Complexity.clBlockEncode, Complexity.clInputSlice, Complexity.clSymbolCode, Turing.FinTM}
\difficulty{1}
For every snapshot $s$ of a machine with a fixed optional-state binary code, the input slice of the encoded block of $s$ equals the two-bit encoding of the input symbol of $s$.
\end{lemma}

\begin{lemma}[Work slice of an encoded snapshot]
\lean{Complexity.clBlock_work}
\label{Complexity.clBlock_work}
\leanok
\uses{Complexity.CLFieldCode, Complexity.Snapshot, Complexity.clBlockEncode, Complexity.clSymbolCode, Complexity.clWorkSlice, Turing.FinTM}
\difficulty{3}
For every snapshot $s$ and every work tape $\tau$ of a machine with a fixed optional-state binary code, the work slice at $\tau$ of the encoded block of $s$ equals the two-bit encoding of the symbol of $s$ on tape $\tau$.
\end{lemma}

\begin{lemma}[Blocks agreeing on all slices are equal]
\lean{Complexity.clBlock_ext}
\label{Complexity.clBlock_ext}
\leanok
\uses{Complexity.CLFieldCode, Complexity.clInputSlice, Complexity.clStateSlice, Complexity.clWorkSlice, Turing.FinTM}
\difficulty{4}
If two snapshot blocks of $B + 2 + 2k$ bits have equal state slices, equal input slices, and equal work slices on every tape, then the blocks are equal bit for bit. Equality of the decoded fields alone would not suffice; the claim pins every individual bit.
\end{lemma}

\begin{lemma}[Injectivity of snapshot encoding]
\lean{Complexity.clBlock_injective}
\label{Complexity.clBlock_injective}
\leanok
\uses{Complexity.CLFieldCode, Complexity.clBlockEncode, Complexity.clBlock_input, Complexity.clBlock_state, Complexity.clBlock_work, Complexity.clInputSlice, Complexity.clStateSlice, Complexity.clSymbolDecode_code, Complexity.clWorkSlice, Turing.FinTM}
\difficulty{4}
For a fixed machine and a fixed binary code of its optional states, the snapshot block encoding is injective: distinct snapshots produce distinct bit blocks, because each individual field code is injective.
\end{lemma}

\begin{definition}[Totalized block decoding]
\lean{Complexity.clBlockDecode}
\label{Complexity.clBlockDecode}
\leanok
\uses{Complexity.CLFieldCode, Complexity.Snapshot, Complexity.clBlockEncode, Turing.FinTM}
A snapshot bit block decodes to the unique snapshot that encodes to it, when one exists; every other bit pattern decodes to one fixed snapshot with halted state and all-blank symbols, rather than repairing fields independently.
\end{definition}

\begin{lemma}[Block decoding inverts encoding]
\lean{Complexity.clBlockDecode_encode}
\label{Complexity.clBlockDecode_encode}
\leanok
\uses{Complexity.CLFieldCode, Complexity.Snapshot, Complexity.clBlockDecode, Complexity.clBlockEncode, Complexity.clBlock_injective, Turing.FinTM}
\difficulty{3}
For every snapshot $s$ of a machine with a fixed optional-state binary code, decoding the encoded block of $s$ returns $s$; the junk fallback is never reached on genuine encodings.
\end{lemma}

\begin{definition}[Width of a snapshot block]
\lean{Complexity.clWidth}
\label{Complexity.clWidth}
\leanok
\uses{Complexity.CLFieldCode, Turing.FinTM}
For a machine with $k$ work tapes and a state code of width $B$, each snapshot block occupies exactly $B + 2 + 2k$ bits; the width is a constant of the fixed machine.
\end{definition}

\begin{definition}[Source block of a window]
\lean{Complexity.clSourceBits}
\label{Complexity.clSourceBits}
\leanok
Within a reconstruction window of $2B+1$ bits, where $B$ is the block width, the source block consists of the first $B$ bits.
\end{definition}

\begin{definition}[Target block of a window]
\lean{Complexity.clTargetBits}
\label{Complexity.clTargetBits}
\leanok
Within a reconstruction window of $2B+1$ bits, where $B$ is the block width, the target block consists of the $B$ bits starting at position $B$.
\end{definition}

\begin{definition}[Selected input bit of a window]
\lean{Complexity.clWindowBit}
\label{Complexity.clWindowBit}
\leanok
Within a reconstruction window of $2B+1$ bits, where $B$ is the block width, the final bit, at position $2B$, carries a selected interior input variable.
\end{definition}

\begin{definition}[Kinds of reconstruction templates]
\lean{Complexity.CLTemplateKind}
\label{Complexity.CLTemplateKind}
\leanok
The finite catalogue of template kinds for the Cook--Levin tableau over a machine with $k$ tapes: an initial kind carrying an input-presence flag, a state kind, an input kind carrying a presence flag, a work kind for each tape carrying an earlier-visit flag, and an acceptance kind. Boundary and absent-previous-visit variants are separate fixed tables.
\end{definition}

\begin{definition}[The six reconstruction predicates]
\lean{Complexity.clPredicate}
\label{Complexity.clPredicate}
\leanok
\uses{Complexity.CLFieldCode, Complexity.CLTemplateKind, Complexity.clBlockDecode, Complexity.clBlockEncode, Complexity.clInputSlice, Complexity.clSourceBits, Complexity.clStateSlice, Complexity.clSymbolCode, Complexity.clTargetBits, Complexity.clWidth, Complexity.clWindowBit, Complexity.clWorkSlice, Complexity.emitted, Complexity.stepState, Complexity.writtenOrKept, Turing.FinTM}
For a machine with a state code and each template kind, a Boolean predicate on windows of $2B+1$ bits, where $B$ is the block width: the initial kind demands that the target block encode the start state reading the selected input bit, or a blank when no input is present; the state kind, that the target state slice encode the stepped state of the decoded source block; the input kind, that the target input slice encode the selected bit or a blank; the work kind for tape $\tau$, that the target work slice at $\tau$ encode the symbol written or kept at the last earlier visit, or a blank on a first visit; and the acceptance kind, that the decoded source snapshot does not emit the output bit false. All kinds share the same constant arity, and unused window bits are harmless.
\end{definition}

\begin{definition}[Window wiring into global variables]
\lean{Complexity.clWire}
\label{Complexity.clWire}
\leanok
\uses{Complexity.clPack}
Given $m$ input variables, a block width $B$, a source time, a target time, and a selected input index $p$, the local window index $j$ is wired to the $j$-th bit of the source block when $j < B$, to the $(j-B)$-th bit of the target block when $B \le j < 2B$, and to $p$ otherwise. Repeated global indices are intentional and supported by relabeling.
\end{definition}

\begin{definition}[One wired clause group]
\lean{Complexity.clGroup}
\label{Complexity.clGroup}
\leanok
\uses{Complexity.CLFieldCode, Complexity.CLTemplateKind, Complexity.clPredicate, Complexity.clTemplate, Complexity.clWidth, Complexity.clWire, Turing.FinTM}
For a template kind together with an input-variable count $m$, a source time, a target time, and a selected input index, the clause group is the fixed CNF table of that kind's reconstruction predicate, relabeled along the window wiring into global variable indices.
\end{definition}

\begin{definition}[Input clause group at a time]
\lean{Complexity.clInputGroup}
\label{Complexity.clInputGroup}
\leanok
\uses{Complexity.CLFieldCode, Complexity.clGroup, Complexity.inputPosAt, Turing.FinTM}
For an input-variable count $m$ and a time $t$, if the oblivious input head at time $t$ sits at an interior position $p$ with $0 < p \le m$, the group is the present-input table wired to input variable $p-1$; otherwise it is the blank-input table with a safe dummy index of zero.
\end{definition}

\begin{definition}[Work clause group at a time and tape]
\lean{Complexity.clWorkGroup}
\label{Complexity.clWorkGroup}
\leanok
\uses{Complexity.CLFieldCode, Complexity.clGroup, Complexity.prevVisit, Turing.FinTM}
For an input-variable count $m$, a time $t$, and a work tape $\tau$: when some strictly earlier time $s$ was the last visit to the cell scanned on $\tau$ at time $t$, the group is the earlier-visit table wired from the snapshot block at $s$; when the cell was never visited before, it is the blank first-visit table.
\end{definition}

\begin{definition}[Group list of the pure tableau]
\lean{Complexity.clTableauGroups}
\label{Complexity.clTableauGroups}
\leanok
\uses{Complexity.CLFieldCode, Complexity.clGroup, Complexity.clGroups, Complexity.clInputGroup, Complexity.clWorkGroup, Turing.FinTM}
For a verifier machine with a state code, certificate parameters $C$ and $e$, an instance $x$, and a horizon $T$, the ordered list of clause groups comprises: one unit clause pinning each of the $|x|$ input variables to the corresponding bit of $x$, the initial group, state groups for each time below $T$, input groups for each time up to $T$, work groups for all times up to $T$ and all tapes, and acceptance groups for each time below $T$. The total variable prefix has length $m = |x| + C(|x|+1)^e$, using the exact original certificate polynomial.
\end{definition}

\begin{definition}[The pure Cook--Levin tableau formula]
\lean{Complexity.clTableau}
\label{Complexity.clTableau}
\leanok
\uses{Complexity.CLFieldCode, Complexity.clTableauGroups, Turing.FinTM}
For a verifier machine with a state code, certificate parameters $C$ and $e$, an instance $x$, and a horizon $T$, the tableau formula is the concatenation of all clause groups of the six families in their fixed order.
\end{definition}

\begin{lemma}[Chunked serialization of the tableau]
\lean{Complexity.clTableau_chunks}
\label{Complexity.clTableau_chunks}
\leanok
\uses{Complexity.CLFieldCode, Complexity.clChunk, Complexity.clGroups_serialize, Complexity.clLastRound, Complexity.clTableau, Complexity.clTableauGroups, Std.Sat.CNF.serialize, Turing.FinTM}
\difficulty{1}
For every verifier machine with a state code, certificate parameters, instance, and horizon, concatenating the round chunks of the tableau's ordered group list over all rounds yields exactly the serialization of the tableau formula, including the empty-template and final-terminator cases.
\end{lemma}

\begin{lemma}[The previous visit is the greatest earlier match]
\lean{Complexity.clPrev_spec}
\label{Complexity.clPrev_spec}
\leanok
\uses{Complexity.prevVisit, Complexity.workPosAt, Turing.FinTM}
\difficulty{3}
Suppose the last-visit function for input count $m$, time $t$, and tape $\tau$ returns a time $s$. Then $s < t$, the work-head position on $\tau$ at time $s$ equals the position at time $t$, and every time $r < t$ scanning the same position satisfies $r \le s$.
\end{lemma}

\begin{lemma}[Saturating cursor orbit]
\lean{Complexity.clCursor_orbit}
\label{Complexity.clCursor_orbit}
\leanok
\difficulty{4}
Let $R$ be a bound, let a packing map assign a word to each cursor value, and suppose a step function sends the packed word of $i$ to the packed word of $\min(i+1, R+1)$ for every $i \le R+1$. Then for all $i$, iterating the step function $i$ times from the packed word of $0$ yields the packed word of $\min(i, R+1)$: the cursor reaches each legal round and then stalls at the one-past-last value.
\end{lemma}

\begin{lemma}[Polynomial emitter from a native body]
\lean{Complexity.clEmitter_of_body}
\label{Complexity.clEmitter_of_body}
\leanok
\uses{Complexity.PolyBound, Complexity.PolyTimeComputable, Complexity.clChunk, Complexity.clChunks_serialize, Complexity.clCursor_orbit, Complexity.clLoop_polyBound, Std.Sat.CNF.serialize, Turing.Cfg.ofWords, Turing.FinTM, Turing.FinTM.ComputesFunInTime, Turing.FinTM.ComputesInTime.mono, Turing.FinTM.exists_emitLoopTM, Turing.MultiTapeTM.initCfg, Turing.MultiTapeTM.runFrom, Turing.stateWord}
\difficulty{5}
Suppose a native machine with a distinguished anchor state, a packing of cursor values into words, abstract step and emission functions, a round bound $R$, and a budget $P$ are given, where $R$ and $P$ are polynomially bounded and some helper machine computes the binary representation of $R$ within time $P$. Assume: stepping a packed cursor advances it saturatingly up to one past the last round; emission at cursor $i \le R$ produces exactly the chunk of round $i$, and nothing beyond the last round; from any input the machine reaches the packed start word, first returning to the anchor, within time $P$; and from each packed cursor the machine returns to the anchor in positive time at most $P$, having appended exactly the prescribed chunk to its output. Then the function sending an input to the serialization of the concatenation of all clause groups over all rounds is computable in polynomial time. No preparation producer or native body is constructed here.
\end{lemma}

\begin{lemma}[Normalized verifier for an NP language]
\lean{Complexity.clNPVerifier}
\label{Complexity.clNPVerifier}
\leanok
\uses{Complexity.NP, Complexity.clObliviousVerifier, Turing.FinTM, Turing.FinTM.DecidesInTime, Turing.FinTM.Oblivious}
\difficulty{1}
Let $L$ be a language in $\mathrm{NP}$. Then there exist certificate parameters $C$ and $e$, a verifier language $V$, a finite oblivious machine $M$, and constants $c, A, d > 0$ such that a string $x$ belongs to $L$ exactly when some certificate $u$ of length $C(|x|+1)^e$ makes the concatenation of $x$ and $u$ belong to $V$, and $M$ decides $V$ within time $c\,(A(n+1)^d+1)^2$ on inputs of length $n$. Neither certificate parameter is enlarged by the normalization.
\end{lemma}

\begin{lemma}[Linear-time contracts are polynomial time]
\lean{Complexity.clNative_linear}
\label{Complexity.clNative_linear}
\leanok
\uses{Complexity.PolyTimeComputable, Turing.FinTM, Turing.FinTM.ComputesFunInTime}
\difficulty{2}
If some finite machine computes a string function $f$ within time $a(n+1)$ on inputs of length $n$, for a constant $a$, then $f$ is polynomial-time computable.
\end{lemma}

\begin{lemma}[Mapping the payload of an encoded pair]
\lean{Complexity.clNative_map}
\label{Complexity.clNative_map}
\leanok
\uses{Complexity.PolyTimeComputable, Turing.FinTM.ComputesInTime.mono, Turing.FinTM.computesFunInTime_pairMapSnd, Turing.pairDecode, Turing.pairEncode}
\difficulty{4}
If $g$ is a polynomial-time computable string function, then so is the function that decodes its input as an encoded pair $(a,b)$ and returns the encoded pair $(a, g(b))$, returning the empty string when the input is not an encoded pair. The first field is retained with its exact bytes.
\end{lemma}

\begin{lemma}[Pairing two computed fields]
\lean{Complexity.clNative_pair}
\label{Complexity.clNative_pair}
\leanok
\uses{Complexity.PolyTimeComputable, Complexity.PolyTimeComputable.comp, Complexity.clNative_linear, Complexity.clNative_map, Turing.FinTM.computesFunInTime_const, Turing.FinTM.computesFunInTime_pairConcat, Turing.FinTM.computesFunInTime_pairDup, Turing.FinTM.computesFunInTime_pairFst, Turing.FinTM.computesFunInTime_pairSnd, Turing.pairDecode_pairEncode, Turing.pairEncode}
\difficulty{5}
If $f$ and $g$ are polynomial-time computable string functions, then so is the function sending a string $x$ to the encoded pair of $f(x)$ and $g(x)$, where both component computations may consult the original input.
\end{lemma}

\begin{lemma}[Appending two computed fields]
\lean{Complexity.clNative_append}
\label{Complexity.clNative_append}
\leanok
\uses{Complexity.PolyTimeComputable, Complexity.PolyTimeComputable.comp, Complexity.clNative_linear, Complexity.clNative_pair, Turing.FinTM.computesFunInTime_pairConcat, Turing.pairDecode_pairEncode}
\difficulty{2}
If $f$ and $g$ are polynomial-time computable string functions, then so is the function sending a string $x$ to the concatenation of $f(x)$ and $g(x)$.
\end{lemma}

\begin{lemma}[Unary value of the certificate polynomial]
\lean{Complexity.clNative_unary}
\label{Complexity.clNative_unary}
\leanok
\uses{Complexity.PolyTimeComputable, Turing.FinTM.computesFunInTime_polyUnary}
\difficulty{1}
For all constants $C$ and $e$, the function sending a string $x$ to the word of exactly $C(|x|+1)^e$ true bits is polynomial-time computable, including the zero-coefficient and zero-degree cases; the exact value is never replaced by a majorant.
\end{lemma}

\begin{definition}[Constant-fill scanning machine]
\lean{Complexity.clFillTM}
\label{Complexity.clFillTM}
\leanok
\uses{Turing.FinTM}
For a fixed bit $b$, a machine with no work tapes scans its input left to right and writes one copy of $b$ to the output for each input bit, halting silently at the right boundary; the empty input is covered by the boundary transition.
\end{definition}

\begin{lemma}[Run of the constant-fill machine]
\lean{Complexity.clFill_run}
\label{Complexity.clFill_run}
\leanok
\uses{Complexity.clFillTM, Turing.Action.apply, Turing.Cfg, Turing.Cfg.inputSymbol, Turing.MultiTapeTM.runFrom, Turing.MultiTapeTM.runFrom_succ_eq_step, Turing.MultiTapeTM.step, Turing.inputSymbolInner, Turing.moveInputPos_pos_of_ne_right}
\difficulty{4}
For a fixed bit $b$, an input $x$, a position $i \le |x|$, and a current output word $w$, running the constant-fill machine for $|x| - i + 1$ steps from the configuration scanning position $i$ with output $w$ halts at the right boundary with output $w$ extended by $|x| - i$ copies of $b$.
\end{lemma}

\begin{lemma}[Polynomial-time constant fill]
\lean{Complexity.clNative_fill}
\label{Complexity.clNative_fill}
\leanok
\uses{Complexity.PolyTimeComputable, Complexity.clFillTM, Complexity.clFill_run, Turing.Cfg, Turing.FinTM.computesInTime_iff, Turing.MultiTapeTM.initCfg}
\difficulty{3}
For every fixed bit $b$, the function sending an input $x$ to the word of $|x|$ copies of $b$ is polynomial-time computable, witnessed by an actual native scan of the input rather than by substituting the instance for the simulator's physical input.
\end{lemma}

\begin{definition}[Arithmetic preparation header]
\lean{Complexity.clPrepHeader}
\label{Complexity.clPrepHeader}
\leanok
\uses{Turing.pairEncode}
For parameters $C, e, c, A, d$ and an instance $x$, set $Q = C(|x|+1)^e$, $m = |x| + Q$, and $T = c\,(A(m+1)^d+1)^2$. The header is the nested pair encoding of six fields: the instance $x$, the binary representations of $Q$, of $m$, and of $T$, the all-false reference input of length $m$, and a unary clock of $T$ true bits. Trajectory and last-visit records are not part of it.
\end{definition}

\begin{lemma}[The header is native polynomial time]
\lean{Complexity.clPrepHeader_native}
\label{Complexity.clPrepHeader_native}
\leanok
\uses{Complexity.PolyTimeComputable, Complexity.PolyTimeComputable.comp, Complexity.clNative_append, Complexity.clNative_fill, Complexity.clNative_linear, Complexity.clNative_pair, Complexity.clNative_unary, Complexity.clPrepHeader, Complexity.polyTimeComputable_id, Turing.FinTM.computesFunInTime_lengthBits}
\difficulty{4}
For all parameters $C, e, c, A, d$, the arithmetic preparation header, comprising the retained instance, the exact binary certificate length, reference length, and horizon, the all-false reference input, and the unary horizon clock, is computable in polynomial time by native machines with sequential access; no positivity assumption is placed on any parameter.
\end{lemma}

\begin{definition}[Embedded reference configuration]
\lean{Complexity.clRefCfg}
\label{Complexity.clRefCfg}
\leanok
\uses{Turing.Cfg, Turing.FinTM, Turing.FinTM.bufferTape, Turing.FinTM.tapeBlocks}
Given a source machine, a host-state embedding, a source configuration over a virtual input word, a position tag bit, a host input position, an administrative bank of $l$ extra tapes with their heads, and an output word, the host configuration places the virtual input on a buffer tape next to the source work tapes and the bank. The source state is embedded into a live host state even when halted, the source output is deliberately absent, and the physical output is the separately supplied word.
\end{definition}

\begin{definition}[One embedded reference action]
\lean{Complexity.clRefAction}
\label{Complexity.clRefAction}
\leanok
\uses{Turing.Action, Turing.FinTM, Turing.FinTM.tapeBlocks, Turing.FinTM.virtualMove, Turing.FinTM.virtualNextTag}
For an embedded source state, a position tag bit, the host work-tape reads, and arbitrary simultaneous operations on the $l$ administrative tapes, the host action simulates one source step through the input buffer: a live source state applies its transition's work-tape writes and virtual input move before its successor state is internalized, so erasures and terminal writes take effect; a halted source keeps all source components fixed. In both cases source emission and physical termination are suppressed, and the bank operations are applied.
\end{definition}

\begin{lemma}[One-step reference correspondence]
\lean{Complexity.clRef_apply}
\label{Complexity.clRef_apply}
\leanok
\uses{Complexity.clRefAction, Complexity.clRefCfg, Turing.Action.apply, Turing.Cfg, Turing.Cfg.inputSymbol, Turing.Cfg.workTapeSymbols, Turing.FinTM, Turing.FinTM.VirtualTag, Turing.FinTM.bufferTape_inputSymbol, Turing.FinTM.tapeBlocks_buffer, Turing.FinTM.virtualMove, Turing.FinTM.virtualMove_correct, Turing.FinTM.virtualNextTag, Turing.MultiTapeTM.step, Turing.MultiTapeTM.step_of_halt, Turing.moveInputPos_zero}
\difficulty{5}
Let $c$ be a source configuration over a virtual input, with a valid position tag bit $b$. Applying the embedded reference action to the embedded host configuration of $c$ yields exactly the embedded configuration of the stepped source configuration, for some valid new tag bit: each administrative tape is updated by its requested write and head move, while the host input position and the physical output are unchanged.
\end{lemma}

\begin{definition}[Clocked reference runner]
\lean{Complexity.clRefClockTM}
\label{Complexity.clRefClockTM}
\leanok
\uses{Complexity.clRefAction, Turing.FinTM, Turing.FinTM.controlAction}
For a source machine, the host machine carries one unary clock tape, one input buffer tape, and one tape per source work tape, and runs the embedded reference simulation while consuming one clock cell per source step. Exhaustion of the clock dispatches to a distinguished live return state; halting of the source alone does not dispatch. Its initial-configuration contract assumes a prepared configuration with the header installed.
\end{definition}

\begin{definition}[Configuration of the clocked runner]
\lean{Complexity.clRefClockCfg}
\label{Complexity.clRefClockCfg}
\leanok
\uses{Complexity.clRefCfg, Complexity.clRefClockTM, Turing.Cfg, Turing.FinTM, Turing.FinTM.bufferTape}
For a source configuration $c$, a position tag bit, a clock budget $T$, and an elapsed count $t$, the clocked-runner configuration embeds $c$ over the virtual input buffer, fills the clock tape with $T$ true cells and places its head at $t$, and has empty output, so source time is recorded separately from any future administrative cost.
\end{definition}

\begin{definition}[Little-endian binary increment]
\lean{Complexity.clCountInc}
\label{Complexity.clCountInc}
\leanok
The increment of a little-endian bit list is as follows: a leading false bit turns into true; a leading true bit becomes false and the carry propagates into the rest; the empty list becomes the single bit true, extending the word on overflow.
\end{definition}

\begin{definition}[Carry length of an increment]
\lean{Complexity.clCountCarry}
\label{Complexity.clCountCarry}
\leanok
For a little-endian bit list, the carry count is the number of initial true bits, which is exactly the number of cells cleared by one increment.
\end{definition}

\begin{lemma}[Increment matches successor]
\lean{Complexity.clCountInc_bits}
\label{Complexity.clCountInc_bits}
\leanok
\uses{Complexity.clCountInc}
\difficulty{4}
For every natural number $n$, applying the list increment to the little-endian binary representation of $n$ yields the binary representation of $n+1$, including the overflow case.
\end{lemma}

\begin{lemma}[Increment grows by at most one bit]
\lean{Complexity.clCountInc_length}
\label{Complexity.clCountInc_length}
\leanok
\uses{Complexity.clCountCarry, Complexity.clCountInc}
\difficulty{3}
For every bit list $w$, the incremented list has length at most $|w| + 1$, and the carry count of $w$ is at most the length of the incremented list, so all cleared cells lie within the incremented word.
\end{lemma}

\begin{definition}[Carry action of the counter]
\lean{Complexity.clCountBump}
\label{Complexity.clCountBump}
\leanok
\uses{Turing.Action}
For a prescribed input-head move and a scanned work symbol, the carry action is as follows: on reading a true bit it writes false, moves the work head right, and keeps carrying; on any other symbol it writes true, moves the work head left, and switches to the rewinding state.
\end{definition}

\begin{definition}[Silent in-place increment machine]
\lean{Complexity.clCountTM}
\label{Complexity.clCountTM}
\leanok
\uses{Complexity.clCountBump, Turing.FinTM, Turing.FinTM.controlAction}
A four-state machine with one work tape increments a little-endian binary counter in place: one state propagates the carry, one rewinds the head to the origin, and an absorbing live state marks completion. The machine is silent: it never advances its native input, never emits output, and never starts a second increment after returning.
\end{definition}

\begin{definition}[Tape holding a finite word]
\lean{Complexity.clCountTape}
\label{Complexity.clCountTape}
\leanok
For a bit list $w$, the tape assigns to each cell $z \ge 0$ the $z$-th bit of $w$ when it exists and a blank otherwise; every negative cell is blank.
\end{definition}

\begin{definition}[Canonical counter configurations]
\lean{Complexity.clCountCfg}
\label{Complexity.clCountCfg}
\leanok
\uses{Complexity.clCountTape, Turing.Cfg}
For an input word, a control state, an input position, a work-head cell, a stored word, and an output, the canonical configuration of the increment machine holds the stored word on its single work tape; these configurations carry the carry, rewind, and return invariants.
\end{definition}

\begin{lemma}[Reading past a prefix]
\lean{Complexity.clCountTape_read}
\label{Complexity.clCountTape_read}
\leanok
\uses{Complexity.clCountTape}
\difficulty{3}
For bit lists $u$ and $v$, the tape holding $u$ followed by $v$ reads, at cell $|u|$, the first bit of $v$, or a blank when $v$ is empty.
\end{lemma}

\begin{lemma}[Writing at the suffix head]
\lean{Complexity.clCountTape_write}
\label{Complexity.clCountTape_write}
\leanok
\uses{Complexity.clCountTape, Complexity.clCountTape_read}
\difficulty{4}
For bit lists $u$ and $v$ and a bit $b$, updating the tape that holds $u$ followed by $v$ at cell $|u|$ with the symbol $b$ yields the tape holding $u$ followed by $b$ and then the tail of $v$: the write either replaces the first suffix bit or extends the word, all other cells are unchanged, and negative cells stay blank.
\end{lemma}

\begin{lemma}[One carry transition of the counter]
\lean{Complexity.clCount_carry_step}
\label{Complexity.clCount_carry_step}
\leanok
\uses{Complexity.clCountBump, Complexity.clCountCfg, Complexity.clCountTM, Complexity.clCountTape, Complexity.clCountTape_read, Complexity.clCountTape_write, Turing.Action.apply, Turing.MultiTapeTM.step, Turing.moveInputPos_zero}
\difficulty{4}
For an input word, an input position, and a stored word split as a prefix $u$ followed by a suffix $v$, one step of the increment machine in carry control at cell $|u|$ yields: when $v$ starts with a true bit, the cell rewritten to false with the head moved right, still in carry control; and otherwise the cell rewritten to true with the head moved left, in rewind control. Only the scanned cell changes.
\end{lemma}

\begin{lemma}[The carry phase flips the leading true bits]
\lean{Complexity.clCount_carry}
\label{Complexity.clCount_carry}
\leanok
\uses{Complexity.clCountCarry, Complexity.clCountCfg, Complexity.clCountInc, Complexity.clCountTM, Complexity.clCount_carry_step, Turing.MultiTapeTM.runFrom, Turing.MultiTapeTM.runFrom_succ_eq_step, Turing.MultiTapeTM.runFrom_zero}
\difficulty{4}
For every input word, input position, prefix $u$, and suffix $v$, running the increment machine in carry control at cell $|u|$ for one more step than the carry count of $v$ replaces $v$ by its increment and leaves the machine in rewind control one cell left of the last written position. Exactly the initial true bits of $v$ are cleared before the final true bit is written.
\end{lemma}

\begin{lemma}[The rewind phase returns to the origin]
\lean{Complexity.clCount_rewind}
\label{Complexity.clCount_rewind}
\leanok
\uses{Complexity.clCountCfg, Complexity.clCountTM, Complexity.clCountTape, Turing.Action.apply, Turing.Cfg.workTapeSymbols, Turing.MultiTapeTM.runFrom, Turing.MultiTapeTM.runFrom_succ_eq_step, Turing.MultiTapeTM.runFrom_zero, Turing.MultiTapeTM.step, Turing.moveInputPos_zero}
\difficulty{4}
For every input word, input position, stored word $w$, and count $j \le |w|$, running the increment machine in rewind control at cell $j - 1$ for $j + 1$ steps crosses the written cells, detects the untouched blank at cell $-1$, and ends at cell $0$ in the absorbing return control, with the stored word unchanged.
\end{lemma}

\begin{lemma}[The carry never exceeds the word length]
\lean{Complexity.clCountCarry_le}
\label{Complexity.clCountCarry_le}
\leanok
\uses{Complexity.clCountCarry}
\difficulty{3}
For every bit list $w$, the carry count of $w$, the number of initial true bits, is at most the length of $w$.
\end{lemma}

\begin{lemma}[Complete timed in-place increment]
\lean{Complexity.clCount_run}
\label{Complexity.clCount_run}
\leanok
\uses{Complexity.clCountCarry, Complexity.clCountCfg, Complexity.clCountInc, Complexity.clCountInc_length, Complexity.clCountTM, Complexity.clCount_carry, Complexity.clCount_rewind, Turing.MultiTapeTM.runFrom, Turing.MultiTapeTM.runFrom_add}
\difficulty{3}
For every input word, input position, and stored word $w$, running the increment machine from carry control at cell $0$ for exactly twice the carry count of $w$ plus two steps ends in the return control at cell $0$ with the stored word replaced by its increment; the native input head never moves from its supplied position.
\end{lemma}

\begin{lemma}[The returned counter is absorbing]
\lean{Complexity.clCount_idle}
\label{Complexity.clCount_idle}
\leanok
\uses{Complexity.clCountTM, Turing.Cfg, Turing.FinTM.controlAction_apply, Turing.MultiTapeTM.step, Turing.moveInputPos_zero}
\difficulty{3}
For every configuration of the increment machine whose control is the return state, one step leaves the entire configuration unchanged, so any upper bound on the running time can be cut to the actual first return.
\end{lemma}

\begin{lemma}[Returned counter configurations stay fixed]
\lean{Complexity.clCount_idle_run}
\label{Complexity.clCount_idle_run}
\leanok
\uses{Complexity.clCountTM, Complexity.clCount_idle, Turing.Cfg, Turing.MultiTapeTM.runFrom, Turing.MultiTapeTM.runFrom_succ_eq_step}
\difficulty{3}
For every configuration of the increment machine in the return control and every duration $t$, running for $t$ steps leaves the whole configuration, not merely its control, unchanged.
\end{lemma}

\begin{lemma}[Strictly positive first return of the counter]
\lean{Complexity.clCount_first}
\label{Complexity.clCount_first}
\leanok
\uses{Complexity.clCountCarry, Complexity.clCountCarry_le, Complexity.clCountCfg, Complexity.clCountInc, Complexity.clCountTM, Complexity.clCount_idle_run, Complexity.clCount_run, Turing.MultiTapeTM.runFrom, Turing.MultiTapeTM.runFrom_add, Turing.MultiTapeTM.runFrom_zero}
\difficulty{5}
For every input word, input position, and stored word $w$, there is a time $t$ with $0 < t \le 2|w| + 2$ such that the increment machine, started in carry control at cell $0$, is never in the return control before $t$ and at time $t$ sits in the return control at cell $0$ with the stored word replaced by its increment.
\end{lemma}

\begin{lemma}[The counter tape is the canonical buffer]
\lean{Complexity.clCountTape_eq}
\label{Complexity.clCountTape_eq}
\leanok
\uses{Complexity.clCountTape, Turing.FinTM.bufferTape}
\difficulty{2}
For every bit list $w$, the tape holding $w$ on the nonnegative cells with blanks elsewhere coincides with the library's canonical buffer tape for $w$, including every negative cell and both empty-word boundaries.
\end{lemma}

\begin{definition}[Counter increment beside a retained reference]
\lean{Complexity.clRefCountTM}
\label{Complexity.clRefCountTM}
\leanok
\uses{Complexity.clCountTM, Turing.FinTM, Turing.FinTM.leftAction}
For a source machine, the host is a machine that runs the four-state binary increment on its leftmost tape while retaining the entire virtual reference configuration, buffer tape, source tapes, and heads, in a disjoint bank; the finite control stores the source state and the boundary tag throughout the call.
\end{definition}

\begin{lemma}[First return of the retained-reference increment]
\lean{Complexity.clRefCount_first}
\label{Complexity.clRefCount_first}
\leanok
\uses{Complexity.clCountCfg, Complexity.clCountInc_bits, Complexity.clCountTM, Complexity.clCountTape_eq, Complexity.clCount_first, Complexity.clRefCfg, Complexity.clRefCountTM, Turing.Cfg, Turing.FinTM, Turing.FinTM.bufferTape, Turing.FinTM.leftCfg, Turing.FinTM.leftCfg_run, Turing.FinTM.tapeBlocks, Turing.MultiTapeTM.runFrom}
\difficulty{5}
For every source configuration, tag bit, and counter value $n$, there is a strictly positive time $t$, at most twice the binary length of $n$ plus two, such that the counter-beside-reference host first reaches its return control exactly at $t$, at which point the counter tape holds the binary representation of $n+1$ with its head restored to zero, while the source configuration, tag, and every bank tape and head are untouched, and the represented source clock has not advanced.
\end{lemma}

\begin{lemma}[Binary width bound]
\lean{Complexity.clCount_width}
\label{Complexity.clCount_width}
\leanok
\difficulty{4}
For all naturals $n$ and $w$ with $n < 2^w$, the binary representation of $n$ has at most $w$ bits, so one administrative increment is charged by two $w$-bit scans plus two transitions, with no unary representation of positions.
\end{lemma}

\begin{definition}[Component action of the counter bank]
\lean{Complexity.clBankPart}
\label{Complexity.clBankPart}
\leanok
\uses{Complexity.clCountTM, Turing.Action, Turing.FinTM.controlAction}
Given a vector of optional component controls and the scanned symbols of $l$ counter tapes, the component action at coordinate $i$ is the increment machine's transition for control $i$ on symbol $i$ when that component is live, and a stationary no-write action when the component is stopped.
\end{definition}

\begin{definition}[Bank of simultaneous binary counters]
\lean{Complexity.clBankTM}
\label{Complexity.clBankTM}
\leanok
\uses{Complexity.clBankPart, Turing.FinTM}
For a fixed number $l$ of disjoint binary-counter tapes, one machine performs simultaneous in-place increments: its finite control holds the vector of component controls, each step applies every component's counter action to its own tape, and a returned component stays idle while others are still carrying or rewinding.
\end{definition}

\begin{definition}[Product configuration of the bank]
\lean{Complexity.clBankCfg}
\label{Complexity.clBankCfg}
\leanok
\uses{Complexity.clBankTM, Turing.Cfg}
Given an input position and one single-tape counter configuration per coordinate, the bank configuration has the vector of component controls as its state, each component's tape and head on the corresponding bank tape, the supplied physical input position, and empty output: every increment is silent.
\end{definition}

\begin{lemma}[Projection of the bank action]
\lean{Complexity.clBank_part}
\label{Complexity.clBank_part}
\leanok
\uses{Complexity.clBankCfg, Complexity.clBankPart, Complexity.clCountTM, Turing.Cfg, Turing.Cfg.inputSymbol, Turing.Cfg.workTapeSymbols, Turing.FinTM.controlAction}
\difficulty{3}
For every input position, every family of component configurations, and every coordinate $i$, the bank's component action at $i$ computed from the product configuration equals the one-tape counter transition of component $i$ on its own scanned symbol, or the stationary action when component $i$ has stopped.
\end{lemma}

\begin{lemma}[One bank step is one step per component]
\lean{Complexity.clBank_step}
\label{Complexity.clBank_step}
\leanok
\uses{Complexity.clBankCfg, Complexity.clBankPart, Complexity.clBankTM, Complexity.clBank_part, Complexity.clCountTM, Turing.Action.apply, Turing.Cfg, Turing.Cfg.inputSymbol, Turing.Cfg.workTapeSymbols, Turing.FinTM.controlAction, Turing.MultiTapeTM.step, Turing.moveInputPos_zero}
\difficulty{4}
For every input position and family of component configurations, one step of the counter bank equals the product of one single-tape counter step per component: tape contents and independently positioned heads are each updated by their own component, the physical input stays fixed, and all component outputs are suppressed.
\end{lemma}

\begin{lemma}[Bank runs project to component runs]
\lean{Complexity.clBank_run}
\label{Complexity.clBank_run}
\leanok
\uses{Complexity.clBankCfg, Complexity.clBankTM, Complexity.clBank_step, Complexity.clCountTM, Turing.Cfg, Turing.MultiTapeTM.runFrom, Turing.MultiTapeTM.runFrom_succ_eq_step'}
\difficulty{3}
For every input position, family of component configurations, and duration $t$, running the counter bank for $t$ steps yields the product of the $t$-step runs of all components, so every counter trajectory is an exact projection of the one native bank run.
\end{lemma}

\begin{definition}[Starting control of a bank component]
\lean{Complexity.clBankStart}
\label{Complexity.clBankStart}
\leanok
A selected counter component starts in the carry control, and an unselected one starts directly in its absorbing return control.
\end{definition}

\begin{lemma}[One component meets the common deadline]
\lean{Complexity.clBank_component}
\label{Complexity.clBank_component}
\leanok
\uses{Complexity.clBankStart, Complexity.clCountCarry, Complexity.clCountCarry_le, Complexity.clCountCfg, Complexity.clCountInc_bits, Complexity.clCountTM, Complexity.clCount_idle_run, Complexity.clCount_run, Turing.MultiTapeTM.runFrom, Turing.MultiTapeTM.runFrom_add}
\difficulty{4}
For every input word, input position, selection flag, counter value $n$, and width bound $W$ with the binary representation of $n$ at most $W$ bits long, running one counter component for $2W + 2$ steps from its selection-dependent start control ends in the return control at head zero, holding the binary representation of $n+1$ if the component was selected and of $n$ otherwise.
\end{lemma}

\begin{lemma}[The bank increments the selected coordinates]
\lean{Complexity.clBank_finish}
\label{Complexity.clBank_finish}
\leanok
\uses{Complexity.clBankCfg, Complexity.clBankStart, Complexity.clBankTM, Complexity.clBank_component, Complexity.clBank_run, Complexity.clCountCfg, Turing.MultiTapeTM.runFrom}
\difficulty{3}
For every input word, input position, selection vector, counter values, and common width bound $W$ covering all the values' binary representations, running the counter bank for $2W + 2$ steps increments exactly the selected coordinates, leaves the unselected ones unchanged, and restores every head to zero; the cost is independent of the number of counters.
\end{lemma}

\begin{lemma}[The all-returned bank is absorbing]
\lean{Complexity.clBank_idle}
\label{Complexity.clBank_idle}
\leanok
\uses{Complexity.clBankPart, Complexity.clBankTM, Complexity.clCountTM, Turing.Action.apply, Turing.Cfg, Turing.FinTM.controlAction, Turing.MultiTapeTM.step, Turing.moveInputPos_zero}
\difficulty{3}
Every configuration of the counter bank whose control has all components in the return state is left entirely unchanged by one step of the bank.
\end{lemma}

\begin{lemma}[First all-returned time of the bank]
\lean{Complexity.clBank_first}
\label{Complexity.clBank_first}
\leanok
\uses{Complexity.clBankCfg, Complexity.clBankStart, Complexity.clBankTM, Complexity.clBank_finish, Complexity.clBank_idle, Complexity.clCountCfg, Turing.MultiTapeTM.runFrom, Turing.MultiTapeTM.runFrom_add}
\difficulty{5}
For every input word, input position, selection vector, counter values, and common width bound $W$, there is a time $t \le 2W + 2$ such that the bank control is not all-returned at any earlier time, and at $t$ the bank holds exactly the requested incremented values with every head restored to zero. The duration may be zero when nothing is selected; positivity is supplied by the surrounding dispatch controller.
\end{lemma}

\begin{definition}[Actual source-head displacements]
\lean{Complexity.clMoves}
\label{Complexity.clMoves}
\leanok
\uses{Turing.FinTM, Turing.FinTM.virtualMove}
For a source state, a position tag bit, a scanned virtual input symbol, and scanned work symbols, the displacement vector records, per coordinate, the clamped virtual input displacement and each work tape's signed move from the source transition; a halted source contributes zero on every coordinate. The input component is the actual movement, not the possibly blocked requested move.
\end{definition}

\begin{definition}[Reference position vector]
\lean{Complexity.clPositions}
\label{Complexity.clPositions}
\leanok
\uses{Turing.Cfg, Turing.FinTM}
For a source configuration, the reference position vector is the buffer's shifted input coordinate, the input position minus one, followed by the signed work-head coordinates; adding one to the first coordinate recovers the input position.
\end{definition}

\begin{lemma}[Displacements track the stepped positions]
\lean{Complexity.clMoves_correct}
\label{Complexity.clMoves_correct}
\leanok
\uses{Complexity.clMoves, Complexity.clPositions, Turing.Action.apply, Turing.Cfg, Turing.Cfg.inputSymbol, Turing.Cfg.workTapeSymbols, Turing.FinTM, Turing.FinTM.VirtualTag, Turing.FinTM.virtualMove_correct, Turing.MultiTapeTM.step, Turing.MultiTapeTM.step_of_halt}
\difficulty{4}
Let $c$ be a source configuration with a valid position tag bit. For every coordinate, the reference position of the stepped configuration equals the reference position of $c$ plus the recorded displacement, including blocked input moves and the halting transition.
\end{lemma}

\begin{definition}[Counter selection from displacements]
\lean{Complexity.clSelect}
\label{Complexity.clSelect}
\leanok
Given a displacement vector on $l$ coordinates, the selection on $2l$ counters is as follows: for each coordinate, the counter in the first bank is marked exactly on a positive move and the counter in the second bank exactly on a negative move; a zero move selects neither.
\end{definition}

\begin{definition}[Advancing the movement counts]
\lean{Complexity.clAdvance}
\label{Complexity.clAdvance}
\leanok
\uses{Complexity.clSelect}
Given a displacement vector and current movement counts on the doubled coordinates, each selected count increases by one and every other count is unchanged.
\end{definition}

\begin{definition}[Signed value of a counter pair]
\lean{Complexity.clSigned}
\label{Complexity.clSigned}
\leanok
Given movement counts on $2l$ coordinates, the signed position at coordinate $i$ is the count in the first bank at $i$ minus the count in the second bank at $i$.
\end{definition}

\begin{lemma}[Signed counts advance by the displacement]
\lean{Complexity.clSigned_advance}
\label{Complexity.clSigned_advance}
\leanok
\uses{Complexity.clAdvance, Complexity.clSelect, Complexity.clSigned}
\difficulty{2}
For every displacement vector, movement counts, and coordinate $i$, the signed value of the advanced counts at $i$ equals the signed value of the original counts plus the displacement at $i$.
\end{lemma}

\begin{lemma}[Counts grow by at most one per step]
\lean{Complexity.clAdvance_bound}
\label{Complexity.clAdvance_bound}
\leanok
\uses{Complexity.clAdvance}
\difficulty{2}
If every movement count is at most $t$, then after advancing along any displacement vector every count is at most $t + 1$.
\end{lemma}

\begin{lemma}[Comparing signed positions without subtraction]
\lean{Complexity.clSigned_eq}
\label{Complexity.clSigned_eq}
\leanok
\uses{Complexity.clSigned}
\difficulty{2}
For movement-count vectors $n$ and $m$ and a coordinate $i$, the signed positions at $i$ agree exactly when the two cross-sums of first-bank and second-bank counts are equal; the criterion compares values rather than the non-unique two-counter codes.
\end{lemma}

\begin{definition}[Source-step and counter tracking machine]
\lean{Complexity.clTrackTM}
\label{Complexity.clTrackTM}
\leanok
\uses{Complexity.clBankStart, Complexity.clBankTM, Complexity.clMoves, Complexity.clRefAction, Complexity.clSelect, Turing.FinTM, Turing.FinTM.controlAction, Turing.FinTM.leftAction}
For a source machine, the tracking host carries two counter banks, one per movement direction, beside the reference bank of buffer and source tapes. In its anchor phase it executes one faithful embedded source action and selects the counters matching the actual clamped displacements; in its bank phase it runs the selected counter increments without advancing the source again, then dispatches back to the anchor. The anchor is a live control phase, distinct from source halting.
\end{definition}

\begin{definition}[Canonical tracking seam]
\lean{Complexity.clTrackCfg}
\label{Complexity.clTrackCfg}
\leanok
\uses{Complexity.clRefCfg, Complexity.clTrackTM, Turing.Cfg, Turing.FinTM, Turing.FinTM.bufferTape}
For a source configuration, a tag bit, and movement counts, the canonical tracking configuration holds every movement counter in canonical binary form with its head at zero, retains the full virtual reference bank, and has empty physical output.
\end{definition}

\begin{lemma}[One source action enters the bank phase]
\lean{Complexity.clTrack_source}
\label{Complexity.clTrack_source}
\leanok
\uses{Complexity.clBankStart, Complexity.clMoves, Complexity.clRefCfg, Complexity.clRef_apply, Complexity.clSelect, Complexity.clTrackCfg, Complexity.clTrackTM, Turing.Action.apply, Turing.Cfg, Turing.Cfg.inputSymbol, Turing.Cfg.workTapeSymbols, Turing.FinTM, Turing.FinTM.VirtualTag, Turing.FinTM.bufferTape, Turing.FinTM.bufferTape_inputSymbol, Turing.MultiTapeTM.step}
\difficulty{4}
Let $c$ be a source configuration with a valid position tag bit and let movement counts be given in canonical form. One step of the tracking host from the canonical seam applies the full source action, produces a valid tag for the stepped configuration, and enters the bank phase with exactly the counters of the actual displacements selected; no counter has yet changed, and the physical output is still empty.
\end{lemma}

\begin{lemma}[Guarded left-bank lifting until first exit]
\lean{Complexity.clLeft_until}
\label{Complexity.clLeft_until}
\leanok
\uses{Turing.Action.apply, Turing.Cfg, Turing.Cfg.inputSymbol, Turing.Cfg.workTapeSymbols, Turing.FinTM.leftAction, Turing.FinTM.leftCfg, Turing.FinTM.leftCfg_apply, Turing.MultiTapeTM, Turing.MultiTapeTM.runFrom, Turing.MultiTapeTM.runFrom_succ_eq_step', Turing.MultiTapeTM.step}
\difficulty{5}
Suppose a host machine agrees, on every embedded non-exit control, with the disjoint-bank lift of a smaller machine acting on the left tapes. Then for every start configuration of the smaller machine, every retained right-bank contents and heads, and every horizon $t$ such that the smaller machine does not reach the exit control before $t$, the host run agrees with the lifted run of the smaller machine at every time up to $t$. No agreement is required after the first exit.
\end{lemma}

\begin{lemma}[The bank frame is the reference frame]
\lean{Complexity.clTrack_frame}
\label{Complexity.clTrack_frame}
\leanok
\uses{Complexity.clBankCfg, Complexity.clBankTM, Complexity.clCountCfg, Complexity.clCountTape_eq, Complexity.clRefCfg, Turing.Cfg, Turing.FinTM, Turing.FinTM.bufferTape, Turing.FinTM.leftCfg, Turing.FinTM.tapeBlocks}
\difficulty{3}
For every source configuration, tag bit, component controls, and counter values, the left-bank lift of the counter-bank product configuration over the retained source tapes coincides literally with the embedded reference configuration whose administrative tapes hold the counters in binary, including the virtual input buffer, every source tape, and all heads.
\end{lemma}

\begin{lemma}[Silent dispatch back to the anchor]
\lean{Complexity.clTrack_dispatch}
\label{Complexity.clTrack_dispatch}
\leanok
\uses{Complexity.clRefCfg, Complexity.clTrackCfg, Complexity.clTrackTM, Turing.Cfg, Turing.FinTM, Turing.FinTM.bufferTape, Turing.FinTM.controlAction_apply, Turing.MultiTapeTM.step, Turing.moveInputPos_zero}
\difficulty{2}
For every source configuration, tag bit, and counter values, one step of the tracking host from the all-returned bank phase yields the anchor control, with no change to the source time, any tape, any head, or the output.
\end{lemma}

\begin{lemma}[One full tracking round]
\lean{Complexity.clTrack_round}
\label{Complexity.clTrack_round}
\leanok
\uses{Complexity.clAdvance, Complexity.clBankCfg, Complexity.clBankStart, Complexity.clBankTM, Complexity.clBank_first, Complexity.clBank_run, Complexity.clCountCfg, Complexity.clLeft_until, Complexity.clMoves, Complexity.clSelect, Complexity.clTrackCfg, Complexity.clTrackTM, Complexity.clTrack_dispatch, Complexity.clTrack_frame, Complexity.clTrack_source, Turing.Cfg, Turing.Cfg.inputSymbol, Turing.Cfg.workTapeSymbols, Turing.FinTM, Turing.FinTM.VirtualTag, Turing.FinTM.bufferTape, Turing.FinTM.leftCfg, Turing.MultiTapeTM.runFrom, Turing.MultiTapeTM.runFrom_succ_eq_step, Turing.MultiTapeTM.runFrom_succ_eq_step', Turing.MultiTapeTM.step}
\difficulty{5}
Let $c$ be a source configuration with a valid position tag bit, and let movement counts be given whose binary representations fit a common width $W$. Then there is a strictly positive time $\tau \le 2W + 4$ such that: every intermediate time strictly between $0$ and $\tau$ is in the bank phase, not at the anchor; and at $\tau$ the tracking host is back at the canonical seam for the stepped source configuration, with a valid new tag and with exactly the clamped displacement counts advanced. Source output is suppressed and every counter head is restored.
\end{lemma}

\begin{lemma}[Appending one bit to a buffer]
\lean{Complexity.clBuffer_append_bit}
\label{Complexity.clBuffer_append_bit}
\leanok
\uses{Complexity.clCountTape_eq, Complexity.clCountTape_write, Turing.FinTM.bufferTape}
\difficulty{1}
For every word $w$ and bit $b$, updating the canonical buffer tape of $w$ at cell $|w|$, its right blank, with $b$ yields the canonical buffer tape of $w$ extended by $b$; no other cell changes, including negative cells and the empty-buffer boundary.
\end{lemma}

\begin{definition}[Two fixed tape slots]
\lean{Complexity.clTwo}
\label{Complexity.clTwo}
\leanok
For two values, the assignment places the first on tape zero and the second on tape one, which gives two disjoint fixed slots with no run-time indexing.
\end{definition}

\begin{definition}[Two-tape record copier]
\lean{Complexity.clCopyTM}
\label{Complexity.clCopyTM}
\leanok
\uses{Complexity.clTwo, Turing.FinTM, Turing.FinTM.controlAction}
A five-state machine with two work tapes copies a self-delimiting field: it doubles each source bit onto the record tape, writes one aligned false-true separator, rewinds the source head to cell zero, and stops in an absorbing completed control with the record head at its new right blank. The physical output stays silent and the source tape is read-only.
\end{definition}

\begin{definition}[Copier configuration]
\lean{Complexity.clCopyCfg}
\label{Complexity.clCopyCfg}
\leanok
\uses{Complexity.clTwo, Turing.Cfg, Turing.FinTM.bufferTape}
For an input word, a control state, an input position, a source-head cell, a source word, and an accumulated record, the copier configuration holds the source word and the record on its two buffer tapes, with the record head at the record's right end.
\end{definition}

\begin{lemma}[One copier write appends one bit]
\lean{Complexity.clCopy_write}
\label{Complexity.clCopy_write}
\leanok
\uses{Complexity.clBuffer_append_bit, Complexity.clCopyCfg, Complexity.clTwo, Turing.Action, Turing.Action.apply, Turing.moveInputPos_zero}
\difficulty{3}
For every copier configuration and every action that writes a bit $b$ at the record head while moving the source head by $d$, applying the action yields the configuration with exactly $b$ appended to the record and the source head moved by $d$; nothing else changes.
\end{lemma}

\begin{lemma}[Two transitions double one source bit]
\lean{Complexity.clCopy_pair}
\label{Complexity.clCopy_pair}
\leanok
\uses{Complexity.clCopyCfg, Complexity.clCopyTM, Complexity.clCopy_write, Complexity.clTwo, Turing.Action.apply, Turing.Cfg.workTapeSymbols, Turing.FinTM.bufferTape, Turing.FinTM.bufferTape_nat, Turing.MultiTapeTM.runFrom, Turing.MultiTapeTM.runFrom_succ_eq_step', Turing.MultiTapeTM.runFrom_zero, Turing.MultiTapeTM.step}
\difficulty{4}
For every input word, input position, source prefix $u$, remaining source bits starting with $b$, and record, two steps of the copier from its reading control at cell $|u|$ append the two bits $b, b$ to the record and advance the source head to cell $|u| + 1$; the source tape is unchanged.
\end{lemma}

\begin{lemma}[The forward pass copies the whole source]
\lean{Complexity.clCopy_forward}
\label{Complexity.clCopy_forward}
\leanok
\uses{Complexity.clCopyCfg, Complexity.clCopyTM, Complexity.clCopy_pair, Turing.MultiTapeTM.runFrom, Turing.MultiTapeTM.runFrom_add}
\difficulty{4}
For every input word, input position, source word split into a prefix $u$ and suffix $w$, and record, running the copier for $2|w|$ steps from the reading control at cell $|u|$ appends the doubling of every bit of $w$ to the record and leaves the source head at the right end; the field separator has not yet been written.
\end{lemma}

\begin{lemma}[The separator writes close the field]
\lean{Complexity.clCopy_separator}
\label{Complexity.clCopy_separator}
\leanok
\uses{Complexity.clCopyCfg, Complexity.clCopyTM, Complexity.clCopy_write, Complexity.clTwo, Turing.Action.apply, Turing.Cfg.workTapeSymbols, Turing.FinTM.bufferTape, Turing.MultiTapeTM.runFrom, Turing.MultiTapeTM.runFrom_succ_eq_step', Turing.MultiTapeTM.runFrom_zero, Turing.MultiTapeTM.step}
\difficulty{4}
For every input word, input position, source word $w$, and record, two steps of the copier from the reading control at the source's right blank yield the rewinding control, with the aligned separator bits false, true appended to the record and the source head moved one cell left.
\end{lemma}

\begin{lemma}[The copier rewinds the source head]
\lean{Complexity.clCopy_rewind}
\label{Complexity.clCopy_rewind}
\leanok
\uses{Complexity.clCopyCfg, Complexity.clCopyTM, Complexity.clTwo, Turing.Action.apply, Turing.Cfg.workTapeSymbols, Turing.FinTM.bufferTape, Turing.FinTM.bufferTape_nat, Turing.MultiTapeTM.runFrom, Turing.MultiTapeTM.runFrom_succ_eq_step, Turing.MultiTapeTM.runFrom_zero, Turing.MultiTapeTM.step, Turing.moveInputPos_zero}
\difficulty{4}
For every input word, input position, source word $w$, record, and count $j \le |w|$, running the copier in rewinding control at source cell $j - 1$ for $j + 1$ steps returns the source head to cell zero and enters the completed control, preserving the entire accumulated record; each crossed bit costs one transition and leaving the left blank costs one more.
\end{lemma}

\begin{lemma}[Exact full record-copy contract]
\lean{Complexity.clCopy_run}
\label{Complexity.clCopy_run}
\leanok
\uses{Complexity.clCopyCfg, Complexity.clCopyTM, Complexity.clCopy_forward, Complexity.clCopy_rewind, Complexity.clCopy_separator, Turing.MultiTapeTM.runFrom, Turing.MultiTapeTM.runFrom_add, Turing.pairEncode}
\difficulty{3}
For every input word, input position, source word $w$, and record, running the copier for $3|w| + 3$ steps from the start yields the completed control, in which the record is extended by exactly the self-delimiting pair encoding of $w$ with empty payload and the source head is restored to zero; the empty field still takes three positive transitions.
\end{lemma}

\begin{lemma}[The completed copier is absorbing]
\lean{Complexity.clCopy_idle}
\label{Complexity.clCopy_idle}
\leanok
\uses{Complexity.clCopyTM, Turing.Action.apply, Turing.Cfg, Turing.FinTM.controlAction, Turing.FinTM.controlAction_apply, Turing.MultiTapeTM.step, Turing.moveInputPos_zero}
\difficulty{3}
Every copier configuration in the completed control is left entirely unchanged by one step of the copier.
\end{lemma}

\begin{lemma}[First positive completion of the copier]
\lean{Complexity.clCopy_first}
\label{Complexity.clCopy_first}
\leanok
\uses{Complexity.clCopyCfg, Complexity.clCopyTM, Complexity.clCopy_idle, Complexity.clCopy_run, Turing.MultiTapeTM.runFrom, Turing.MultiTapeTM.runFrom_add, Turing.MultiTapeTM.runFrom_zero, Turing.pairEncode}
\difficulty{5}
For every input word, input position, source word $w$, and record, there is a time $t$ with $0 < t \le 3|w| + 3$ such that the copier is never in the completed control before $t$, and at $t$ it holds the record extended by the self-delimiting encoding of $w$, with the source untouched and both heads as specified; a host can therefore dispatch on completion without knowing the field length.
\end{lemma}

\begin{definition}[Relocating an action to fixed host slots]
\lean{Complexity.clSlotAction}
\label{Complexity.clSlotAction}
\leanok
\uses{Turing.Action}
Given a partial selection from $l$ host tapes back to $k$ local tapes and a state embedding, an action on $k$ tapes is relocated to the host: each selected host tape receives the corresponding local write and move, every unselected host tape is stationary, and the input move, output emission, and embedded successor state are carried over.
\end{definition}

\begin{definition}[Relocated phase configuration]
\lean{Complexity.clSlotCfg}
\label{Complexity.clSlotCfg}
\leanok
\uses{Turing.Cfg}
Given a partial selection from host tapes to local tapes, a state embedding, and retained contents and heads for the unselected host tapes, a local configuration is relocated to the host: selected tapes and heads come from the local configuration, unselected ones keep their retained values, and the input position and physical output are those of the local configuration.
\end{definition}

\begin{lemma}[Relocation commutes with one action]
\lean{Complexity.clSlot_apply}
\label{Complexity.clSlot_apply}
\leanok
\uses{Complexity.clSlotAction, Complexity.clSlotCfg, Turing.Action, Turing.Action.apply, Turing.Cfg}
\difficulty{3}
For every partial selection, state embedding, retained tapes and heads, action, and local configuration, applying the relocated action to the relocated configuration equals relocating the result of applying the action locally; writes, head motion, and final emission are preserved, and inactive host tapes and heads are fixed.
\end{lemma}

\begin{lemma}[Guarded relocation of a whole phase]
\lean{Complexity.clSlot_run}
\label{Complexity.clSlot_run}
\leanok
\uses{Complexity.clSlotAction, Complexity.clSlotCfg, Complexity.clSlot_apply, Turing.Cfg, Turing.Cfg.workTapeSymbols, Turing.MultiTapeTM, Turing.MultiTapeTM.runFrom, Turing.MultiTapeTM.runFrom_succ_eq_step', Turing.MultiTapeTM.step, Turing.MultiTapeTM.step_of_halt}
\difficulty{5}
Let a host machine agree, on every embedded control satisfying a guard, with the relocation of a smaller machine's transition through a tape injection whose selection is a left inverse. Then for every retained host contents and heads, every local start configuration, and every horizon $t$ such that all local controls strictly before $t$ satisfy the guard, the host run up to time $t$ equals the relocation of the local run; the guard is not needed at the endpoint itself.
\end{lemma}

\begin{definition}[Tape slots of one row-copy call]
\lean{Complexity.clRowIndex}
\label{Complexity.clRowIndex}
\leanok
\uses{Complexity.clTwo}
For a selected counter coordinate $i$ among $l$ counters, the two tapes of a copy call are embedded as: the source slot is counter tape $i$, and the record slot is the final extra tape. The index is fixed in finite control for the whole field.
\end{definition}

\begin{definition}[Partial inverse of the row-copy slots]
\lean{Complexity.clRowSelect}
\label{Complexity.clRowSelect}
\leanok
For a selected counter coordinate $i$, the selection maps counter tape $i$ back to the copy call's source slot and the final record tape back to its record slot, and leaves every other counter tape unselected.
\end{definition}

\begin{lemma}[Row-copy selection inverts the injection]
\lean{Complexity.clRow_inverse}
\label{Complexity.clRow_inverse}
\leanok
\uses{Complexity.clRowIndex, Complexity.clRowSelect, Complexity.clTwo}
\difficulty{2}
For every selected counter coordinate $i$ and each of the two local copy tapes, the row-copy selection applied to the embedded slot returns that local tape, so both actual copy tapes are selected back to their local indices.
\end{lemma}

\begin{definition}[Row controller copying every counter]
\lean{Complexity.clRowTM}
\label{Complexity.clRowTM}
\leanok
\uses{Complexity.clCopyTM, Complexity.clRowIndex, Complexity.clRowSelect, Complexity.clSlotAction, Turing.FinTM, Turing.FinTM.controlAction}
For $l$ counter tapes plus one record tape, a controller runs the two-tape field copier on each counter in fixed tape order: its control holds the current field index and the copier control, one silent dispatch is charged after each completed field, and after the last field the controller idles in a live end-of-row control.
\end{definition}

\begin{definition}[Row-copy seam configuration]
\lean{Complexity.clRowCfg}
\label{Complexity.clRowCfg}
\leanok
\uses{Complexity.clRowTM, Turing.Cfg, Turing.FinTM.bufferTape}
For an input word, an input position, a field index, a copier control, stored counter words, and an accumulated record, the row seam holds every counter as a buffer at head zero and the append-only record with its head at the record's exact current length; the physical output is empty.
\end{definition}

\begin{lemma}[The row seam is a relocated copier]
\lean{Complexity.clRow_frame}
\label{Complexity.clRow_frame}
\leanok
\uses{Complexity.clCopyCfg, Complexity.clRowCfg, Complexity.clRowSelect, Complexity.clSlotCfg, Complexity.clTwo, Turing.FinTM.bufferTape}
\difficulty{4}
For every input word, input position, selected counter $i$, copier control, counter words, and record, relocating the two-tape copier configuration for counter $i$ into the row controller's slots, with all other counters retained, yields exactly the row seam at field $i$; every unselected counter retains its whole word, not just its current cell.
\end{lemma}

\begin{lemma}[One field is stored and the index advances]
\lean{Complexity.clRow_field}
\label{Complexity.clRow_field}
\leanok
\uses{Complexity.clCopyCfg, Complexity.clCopyTM, Complexity.clCopy_first, Complexity.clRowCfg, Complexity.clRowIndex, Complexity.clRowSelect, Complexity.clRowTM, Complexity.clRow_frame, Complexity.clRow_inverse, Complexity.clSlot_run, Turing.Action.apply, Turing.FinTM.bufferTape, Turing.FinTM.controlAction_apply, Turing.MultiTapeTM.runFrom, Turing.MultiTapeTM.runFrom_succ_eq_step', Turing.moveInputPos_zero, Turing.pairEncode}
\difficulty{5}
For every input word, input position, selected counter $i$, counter words $w$, and record, there is a strictly positive time $t \le 3|w(i)| + 4$ after which the row controller has appended the self-delimiting encoding of $w(i)$ to the record and advanced its field index from $i$ to $i+1$, with the entire counter bank, the input head, and the empty physical output restored.
\end{lemma}

\begin{definition}[Stored prefix of a row]
\lean{Complexity.clRowPrefix}
\label{Complexity.clRowPrefix}
\leanok
\uses{Turing.pairEncode}
For every family of counter words on $l$ coordinates and every count $j$, the stored row prefix is the concatenation of the self-delimiting encodings of the first $j$ words in increasing tape order; indices beyond $l$ contribute nothing, and actual controllers only use prefixes up to the field count.
\end{definition}

\begin{lemma}[The controller stores each row prefix]
\lean{Complexity.clRow_prefix_run}
\label{Complexity.clRow_prefix_run}
\leanok
\uses{Complexity.clRowCfg, Complexity.clRowPrefix, Complexity.clRowTM, Complexity.clRow_field, Turing.MultiTapeTM.runFrom, Turing.MultiTapeTM.runFrom_add}
\difficulty{4}
For every input word, input position, counter words all of length at most $W$, record, and count $j \le l$, there is a time $t \le j(3W + 4)$ after which the row controller, started at field zero, has appended exactly the stored prefix of the first $j$ fields to the record and sits at field index $j$ with every counter and head restored.
\end{lemma}

\begin{lemma}[A complete row is stored]
\lean{Complexity.clRow_stored}
\label{Complexity.clRow_stored}
\leanok
\uses{Complexity.clRowCfg, Complexity.clRowPrefix, Complexity.clRowTM, Complexity.clRow_prefix_run, Turing.MultiTapeTM.runFrom}
\difficulty{1}
For every input word, input position, counter words of length at most $W$, and record, there is a time $t \le l(3W+4)$ after which the row controller has appended the complete stored row, all $l$ self-delimiting fields in tape order, to the record, with every source counter unchanged and its head at zero.
\end{lemma}

\begin{lemma}[The completed row controller is absorbing]
\lean{Complexity.clRow_idle}
\label{Complexity.clRow_idle}
\leanok
\uses{Complexity.clRowTM, Turing.Cfg, Turing.FinTM.controlAction_apply, Turing.MultiTapeTM.step, Turing.moveInputPos_zero}
\difficulty{3}
Every configuration of the row controller whose control is the end-of-row state with copier control zero is left entirely unchanged by one step.
\end{lemma}

\begin{lemma}[First end-of-row return]
\lean{Complexity.clRow_first}
\label{Complexity.clRow_first}
\leanok
\uses{Complexity.clRowCfg, Complexity.clRowPrefix, Complexity.clRowTM, Complexity.clRow_idle, Complexity.clRow_stored, Turing.MultiTapeTM.runFrom, Turing.MultiTapeTM.runFrom_add, Turing.MultiTapeTM.runFrom_zero}
\difficulty{5}
For a positive number $l$ of counters, every input word, input position, counter words of length at most $W$, and record, there is a time $t$ with $0 < t \le l(3W+4)$ such that the row controller never shows the end-of-row control before $t$, and at $t$ the record has been extended by the complete stored row with all counter heads restored; positivity comes from the fixed positive field count.
\end{lemma}

\begin{lemma}[Length ledger of a stored row]
\lean{Complexity.clRowPrefix_length}
\label{Complexity.clRowPrefix_length}
\leanok
\uses{Complexity.clRowPrefix, Turing.pairEncode}
\difficulty{4}
For counter words of length at most $W$ and every count $j$, the stored row prefix of the first $j$ fields has length at most $j(2W + 2)$: each field costs its doubled bits plus the two separator bits, independently of the native copying runtime.
\end{lemma}

\begin{lemma}[Widths fit the elapsed time]
\lean{Complexity.clElapsed_width}
\label{Complexity.clElapsed_width}
\leanok
\uses{Complexity.clCount_width}
\difficulty{1}
For all naturals $n \le t$, the binary representation of $n$ has at most $t$ bits. The bound coarsely charges native scans; all stored movement counts remain in canonical binary form.
\end{lemma}

\begin{definition}[Deterministic virtual-input tag]
\lean{Complexity.clTag}
\label{Complexity.clTag}
\leanok
For an input position in a word of length $n$, extended by the two boundary blanks, the deterministic tag is true exactly when the position is nonzero; it is used only in the pure record specification.
\end{definition}

\begin{lemma}[The deterministic tag is valid]
\lean{Complexity.clTag_valid}
\label{Complexity.clTag_valid}
\leanok
\uses{Complexity.clTag, Turing.FinTM.VirtualTag}
\difficulty{2}
For every input position, the deterministic nonzero-position tag correctly identifies both boundary blanks and is a valid virtual-input tag for that position.
\end{lemma}

\begin{lemma}[Move symbols embed into the integers]
\lean{Complexity.clMove_inj}
\label{Complexity.clMove_inj}
\leanok
\difficulty{2}
For any two of the three native move symbols, if their integer displacements are equal then the symbols themselves are equal.
\end{lemma}

\begin{lemma}[All valid tags give the same displacement]
\lean{Complexity.clMoves_tag}
\label{Complexity.clMoves_tag}
\leanok
\uses{Complexity.clMove_inj, Complexity.clMoves, Complexity.clMoves_correct, Complexity.clTag, Complexity.clTag_valid, Turing.Cfg, Turing.Cfg.inputSymbol, Turing.Cfg.workTapeSymbols, Turing.FinTM, Turing.FinTM.VirtualTag}
\difficulty{3}
Let $c$ be a source configuration and let $b$ be any valid position tag for its input position. The displacement vector computed from $b$ equals the displacement vector computed from the deterministic nonzero-position tag; at an interior cell either tag yields the same clamped movement.
\end{lemma}

\begin{definition}[Canonical movement counts over time]
\lean{Complexity.clCounts}
\label{Complexity.clCounts}
\leanok
\uses{Complexity.clAdvance, Complexity.clMoves, Complexity.clTag, Turing.Cfg.inputSymbol, Turing.Cfg.workTapeSymbols, Turing.FinTM, Turing.MultiTapeTM.initCfg, Turing.MultiTapeTM.runFrom}
For a source machine and a virtual input $y$, the movement counts at time $0$ are all zero, and the counts at time $t+1$ advance the counts at time $t$ by the actual clamped displacements of the source step at time $t$, computed with the deterministic tag. Counts record actual moves, not requested outward moves, and freeze after the source halts.
\end{definition}

\begin{lemma}[One round advances the specified counts]
\lean{Complexity.clCounts_succ}
\label{Complexity.clCounts_succ}
\leanok
\uses{Complexity.clAdvance, Complexity.clCounts, Complexity.clMoves, Complexity.clMoves_tag, Turing.Cfg.inputSymbol, Turing.Cfg.workTapeSymbols, Turing.FinTM, Turing.FinTM.VirtualTag, Turing.MultiTapeTM.initCfg, Turing.MultiTapeTM.runFrom}
\difficulty{1}
For every source machine, virtual input $y$, time $t$, and valid tag $b$ for the position reached at time $t$, advancing the counts at time $t$ by the displacements computed from $b$ yields exactly the counts at time $t+1$.
\end{lemma}

\begin{lemma}[Counts are bounded by elapsed time]
\lean{Complexity.clCounts_bound}
\label{Complexity.clCounts_bound}
\leanok
\uses{Complexity.clAdvance_bound, Complexity.clCounts, Turing.FinTM}
\difficulty{3}
For every source machine, virtual input $y$, time $t$, and coordinate, the stored movement count at time $t$ is at most $t$.
\end{lemma}

\begin{lemma}[Signed counts are the exact positions]
\lean{Complexity.clCounts_positions}
\label{Complexity.clCounts_positions}
\leanok
\uses{Complexity.clAdvance, Complexity.clCounts, Complexity.clMoves_correct, Complexity.clPositions, Complexity.clSigned, Complexity.clSigned_advance, Complexity.clTag_valid, Turing.Cfg.init, Turing.FinTM, Turing.MultiTapeTM.initCfg, Turing.MultiTapeTM.runFrom, Turing.MultiTapeTM.runFrom_succ_eq_step'}
\difficulty{4}
For every source machine, virtual input $y$, time $t$, and coordinate, the signed value of the stored count pair equals the reference position of the configuration reached at time $t$: the shifted clamped input position on the first coordinate and the signed work-head positions on the rest.
\end{lemma}

\begin{definition}[Field count of a stored row]
\lean{Complexity.clRecFields}
\label{Complexity.clRecFields}
\leanok
\uses{Turing.FinTM}
For a source machine with $k$ work tapes, each stored trajectory row has $(1+k) + (1+k)$ binary movement-count fields: one positive and one negative counter per head, input head included.
\end{definition}

\begin{definition}[Serialized trajectory prefix]
\lean{Complexity.clRecords}
\label{Complexity.clRecords}
\leanok
\uses{Complexity.clCounts, Complexity.clRecFields, Complexity.clRowPrefix, Turing.FinTM}
For a source machine and a virtual input $y$, the serialized trajectory prefix at count $t$ is the concatenation of the stored rows of the movement counts at times $0, \dots, t-1$, each row consisting of all its self-delimiting binary fields in tape order; the complete inclusive trajectory at horizon $T$ uses prefix length $T+1$.
\end{definition}

\begin{definition}[Control phases of the recorder]
\lean{Complexity.clRecState}
\label{Complexity.clRecState}
\leanok
\uses{Complexity.clRecFields, Complexity.clRowTM, Complexity.clTrackTM, Turing.FinTM}
The recorder's finite control is a disjoint sum of five phases: row copying, carrying the source state, tag, and row-controller control; the clock test; the source-step release; the running counter update, carrying a tracking control; and the completed live return.
\end{definition}

\begin{definition}[Tracking tapes inside the recorder]
\lean{Complexity.clRecTrackIndex}
\label{Complexity.clRecTrackIndex}
\leanok
\uses{Complexity.clRecFields, Complexity.clTrackTM, Turing.FinTM}
The tracking machine's tapes, two counter banks, the virtual input buffer, and the source tapes, are embedded into the recorder's tapes by skipping the record slot and the unary clock slot.
\end{definition}

\begin{definition}[Partial inverse for the tracking tapes]
\lean{Complexity.clRecTrackSelect}
\label{Complexity.clRecTrackSelect}
\leanok
\uses{Complexity.clRecFields, Complexity.clTrackTM, Turing.FinTM}
The recorder's selection maps each counter-bank tape and each source-side tape back to its tracking-machine coordinate and leaves the record and clock slots unselected, so both administrative record tapes are preserved during tracking.
\end{definition}

\begin{definition}[Row-copy tapes inside the recorder]
\lean{Complexity.clRecRowIndex}
\label{Complexity.clRecRowIndex}
\leanok
\uses{Complexity.clRecFields, Complexity.clRowTM, Turing.FinTM}
The row controller's tapes, the movement counters plus the record tape, are embedded into the recorder's tapes with the counters first and the record slot immediately after; the clock, buffer, and source tapes are skipped.
\end{definition}

\begin{definition}[Partial inverse for the row-copy tapes]
\lean{Complexity.clRecRowSelect}
\label{Complexity.clRecRowSelect}
\leanok
\uses{Complexity.clRecFields, Complexity.clRowTM, Turing.FinTM}
The recorder's selection maps each movement counter and the record slot back to its row-controller coordinate and leaves the clock, virtual-input buffer, and every source tape unselected.
\end{definition}

\begin{lemma}[Tracking selection inverts its injection]
\lean{Complexity.clRecTrack_inverse}
\label{Complexity.clRecTrack_inverse}
\leanok
\uses{Complexity.clRecTrackIndex, Complexity.clRecTrackSelect, Complexity.clTrackTM, Turing.FinTM}
\difficulty{2}
For every tracking-machine tape coordinate, the recorder's tracking selection applied to the embedded coordinate returns that coordinate.
\end{lemma}

\begin{lemma}[Row selection inverts its injection]
\lean{Complexity.clRecRow_inverse}
\label{Complexity.clRecRow_inverse}
\leanok
\uses{Complexity.clRecFields, Complexity.clRecRowIndex, Complexity.clRecRowSelect, Complexity.clRowTM, Turing.FinTM}
\difficulty{2}
For every row-controller tape coordinate, the recorder's row selection applied to the embedded coordinate returns that coordinate.
\end{lemma}

\begin{definition}[Slot of the unary clock]
\lean{Complexity.clRecClockIndex}
\label{Complexity.clRecClockIndex}
\leanok
\uses{Complexity.clRecFields, Turing.FinTM}
The unary clock slot is the recorder tape placed after the movement counters and the record slot and before the virtual-input buffer and the source tapes, disjoint from both the source and row-copy banks.
\end{definition}

\begin{definition}[Native inclusive trajectory recorder]
\lean{Complexity.clRecTM}
\label{Complexity.clRecTM}
\leanok
\uses{Complexity.clRecClockIndex, Complexity.clRecFields, Complexity.clRecRowIndex, Complexity.clRecRowSelect, Complexity.clRecState, Complexity.clRecTrackIndex, Complexity.clRecTrackSelect, Complexity.clRowTM, Complexity.clSlotAction, Complexity.clTrackTM, Turing.FinTM, Turing.FinTM.controlAction}
For a source machine, the recorder cycles through four phases: it copies the current row of movement counters onto the record, tests the unary clock, releases exactly one tracked source step with its counter updates, and returns to row copying; when the clock is exhausted it enters a completed live control. The current row is stored before the clock test, so both the initial and the final rows are present, and every counter and copy transition lies outside source time.
\end{definition}

\begin{definition}[Whole recorder configuration]
\lean{Complexity.clRecCfg}
\label{Complexity.clRecCfg}
\leanok
\uses{Complexity.clRecFields, Complexity.clRecState, Complexity.clRecTM, Turing.Cfg, Turing.FinTM, Turing.FinTM.bufferTape, Turing.FinTM.tapeBlocks}
For a recorder control, a source configuration, movement counts, a stored record, a clock budget $T$, and a logical time $t$, the recorder configuration holds every movement counter in canonical binary at head zero, the actual record prefix with its head at the record length, the unary clock with its head at $t$, and the complete virtual reference bank of buffer and source tapes; only the record and clock heads may be away from zero.
\end{definition}

\begin{lemma}[Relocated row copying matches the recorder seam]
\lean{Complexity.clRec_row_frame}
\label{Complexity.clRec_row_frame}
\leanok
\uses{Complexity.clRecCfg, Complexity.clRecFields, Complexity.clRecRowSelect, Complexity.clRecState, Complexity.clRowCfg, Complexity.clSlotCfg, Turing.Cfg, Turing.FinTM}
\difficulty{3}
For every source configuration, tag, movement counts, record, clock budget, logical time, field index, and copier control, relocating the row-copy seam into the recorder's tapes, retaining every inactive slot from the recorder configuration itself, yields exactly the recorder configuration in the row-copying phase; the inactive bank includes every source cell and head and the unchanged clock.
\end{lemma}

\begin{lemma}[Row-frame inactivity]
\lean{Complexity.clRec_row_inactive}
\label{Complexity.clRec_row_inactive}
\leanok
\uses{Complexity.clRecCfg, Complexity.clRecFields, Complexity.clRecRowSelect, Complexity.clRecState, Complexity.clRowTM, Complexity.clSlotCfg, Turing.Cfg, Turing.FinTM}
\difficulty{3}
For every source configuration, tag, movement counts, clock budget, and logical time, the retained inactive tapes and heads used when relocating a row-copy configuration into the recorder are the same for any two record contents and row controls, so consecutive fields and rows can append to the record without changing source time.
\end{lemma}

\begin{lemma}[The recorder stores one complete row]
\lean{Complexity.clRec_copy}
\label{Complexity.clRec_copy}
\leanok
\uses{Complexity.clRecCfg, Complexity.clRecFields, Complexity.clRecRowIndex, Complexity.clRecRowSelect, Complexity.clRecRow_inverse, Complexity.clRecState, Complexity.clRecTM, Complexity.clRec_row_frame, Complexity.clRec_row_inactive, Complexity.clRowCfg, Complexity.clRowPrefix, Complexity.clRowTM, Complexity.clRow_first, Complexity.clSlot_run, Turing.Cfg, Turing.FinTM, Turing.FinTM.controlAction_apply, Turing.MultiTapeTM.runFrom, Turing.MultiTapeTM.runFrom_succ_eq_step', Turing.MultiTapeTM.step, Turing.moveInputPos_zero}
\difficulty{5}
For every source configuration, tag, movement counts with binary representations of width at most $W$, record, clock budget, and logical time, there is a time $s \le F(3W+4) + 1$, where $F$ is the field count, after which the recorder has appended the complete stored row of the counters to the record and dispatched into its clock-test phase; all source storage and the logical clock are unchanged throughout.
\end{lemma}

\begin{lemma}[Left and right bank slots are disjoint]
\lean{Complexity.clLeft_ne_right}
\label{Complexity.clLeft_ne_right}
\leanok
\difficulty{2}
For every left coordinate among the first $a$ tapes and every right coordinate among the following $b$ tapes of a disjoint bank, the two embedded tape indices are distinct.
\end{lemma}

\begin{lemma}[Right and left bank slots are disjoint]
\lean{Complexity.clRight_ne_left}
\label{Complexity.clRight_ne_left}
\leanok
\uses{Complexity.clLeft_ne_right}
\difficulty{1}
For every right coordinate and left coordinate of a disjoint tape bank, the embedded right index is distinct from the embedded left index.
\end{lemma}

\begin{lemma}[One clock cell authorizes one source step]
\lean{Complexity.clRec_tick}
\label{Complexity.clRec_tick}
\leanok
\uses{Complexity.clLeft_ne_right, Complexity.clRecCfg, Complexity.clRecClockIndex, Complexity.clRecFields, Complexity.clRecTM, Complexity.clRight_ne_left, Turing.Action.apply, Turing.Cfg, Turing.Cfg.workTapeSymbols, Turing.FinTM, Turing.FinTM.tapeBlocks, Turing.MultiTapeTM.step, Turing.moveInputPos_zero}
\difficulty{4}
For every source configuration, tag, movement counts, record, clock budget $T$, and logical time $t < T$, one recorder step from the clock-test phase reads the live unary clock cell, advances only the clock head from $t$ to $t+1$, and enters the source-release phase; no source head moves and no record bit is appended.
\end{lemma}

\begin{lemma}[The exhausted clock dispatches to completion]
\lean{Complexity.clRec_stop}
\label{Complexity.clRec_stop}
\leanok
\uses{Complexity.clRecCfg, Complexity.clRecClockIndex, Complexity.clRecFields, Complexity.clRecTM, Turing.Action.apply, Turing.Cfg, Turing.Cfg.workTapeSymbols, Turing.FinTM, Turing.FinTM.controlAction_apply, Turing.MultiTapeTM.step, Turing.moveInputPos_zero}
\difficulty{3}
For every source configuration, tag, movement counts, record, and clock budget $T$, one recorder step from the clock-test phase at logical time $T$ reads the exhausted clock blank and yields the completed live control, silently; no further source transition can occur.
\end{lemma}

\begin{lemma}[Relocated tracking matches the recorder seam]
\lean{Complexity.clRec_track_frame}
\label{Complexity.clRec_track_frame}
\leanok
\uses{Complexity.clRecCfg, Complexity.clRecFields, Complexity.clRecState, Complexity.clRecTrackSelect, Complexity.clRefCfg, Complexity.clSlotCfg, Complexity.clTrackCfg, Turing.Cfg, Turing.FinTM, Turing.FinTM.tapeBlocks}
\difficulty{3}
For every source configuration, tag, movement counts, record, clock budget, and logical time, relocating the canonical tracking seam into the recorder's tapes, retaining the record and clock slots from the recorder configuration itself, yields exactly the recorder configuration in the counter-update phase.
\end{lemma}

\begin{lemma}[Tracking-frame inactivity]
\lean{Complexity.clRec_track_inactive}
\label{Complexity.clRec_track_inactive}
\leanok
\uses{Complexity.clRecCfg, Complexity.clRecFields, Complexity.clRecState, Complexity.clRecTrackSelect, Complexity.clSlotCfg, Complexity.clTrackTM, Turing.Cfg, Turing.FinTM}
\difficulty{3}
The retained inactive tapes and heads used when relocating a tracking configuration into the recorder depend only on the stored record and clock, so they agree for any two source configurations, tags, and movement counts; source writes and counter updates therefore do not disturb the frame.
\end{lemma}

\begin{lemma}[Releasing a relocated subroutine]
\lean{Complexity.clSlot_release}
\label{Complexity.clSlot_release}
\leanok
\uses{Complexity.clSlotAction, Complexity.clSlotCfg, Complexity.clSlot_apply, Turing.Action.apply, Turing.Cfg, Turing.Cfg.inputSymbol, Turing.Cfg.workTapeSymbols, Turing.MultiTapeTM, Turing.MultiTapeTM.step}
\difficulty{4}
Suppose a host control acts, on all reads, as the relocation of a smaller machine's transition at a fixed local state $q$, through a tape injection whose selection is a left inverse. Then one host step from the relocated configuration of any local configuration in state $q$, with its control replaced by that release control, equals the relocation of one local step. Executing the first action before observing any return guard handles equal entry and exit controls.
\end{lemma}

\begin{lemma}[One recorded source round]
\lean{Complexity.clRec_advance}
\label{Complexity.clRec_advance}
\leanok
\uses{Complexity.clAdvance, Complexity.clMoves, Complexity.clRecCfg, Complexity.clRecFields, Complexity.clRecState, Complexity.clRecTM, Complexity.clRecTrackIndex, Complexity.clRecTrackSelect, Complexity.clRecTrack_inverse, Complexity.clRec_track_frame, Complexity.clRec_track_inactive, Complexity.clSlotCfg, Complexity.clSlot_release, Complexity.clSlot_run, Complexity.clTrackCfg, Complexity.clTrackTM, Complexity.clTrack_round, Turing.Cfg, Turing.Cfg.inputSymbol, Turing.Cfg.workTapeSymbols, Turing.FinTM, Turing.FinTM.VirtualTag, Turing.FinTM.controlAction_apply, Turing.MultiTapeTM.runFrom, Turing.MultiTapeTM.runFrom_succ_eq_step, Turing.MultiTapeTM.runFrom_succ_eq_step', Turing.MultiTapeTM.step, Turing.moveInputPos_zero}
\difficulty{5}
Let $c$ be a source configuration with a valid position tag, and let the movement counts have binary representations of width at most $W$. From the source-release phase, there is a time $s \le 2W + 5$ after which the recorder has executed exactly one source transition together with its full signed counter update and dispatched to the next row-copy phase, with a valid tag for the stepped configuration; the record and clock contents and heads are framed throughout.
\end{lemma}

\begin{lemma}[Recorder invariant through the horizon]
\lean{Complexity.clRec_prefix}
\label{Complexity.clRec_prefix}
\leanok
\uses{Complexity.clCounts, Complexity.clCounts_bound, Complexity.clCounts_succ, Complexity.clElapsed_width, Complexity.clRecCfg, Complexity.clRecFields, Complexity.clRecTM, Complexity.clRec_advance, Complexity.clRec_copy, Complexity.clRec_tick, Complexity.clRecords, Turing.Cfg.init, Turing.FinTM, Turing.FinTM.VirtualTag, Turing.MultiTapeTM.initCfg, Turing.MultiTapeTM.runFrom, Turing.MultiTapeTM.runFrom_add, Turing.MultiTapeTM.runFrom_succ_eq_step, Turing.MultiTapeTM.runFrom_succ_eq_step'}
\difficulty{5}
For every source machine, physical input, virtual input $y$, horizon $T$, and logical time $j \le T$, there is a time $s \le j\,(F(3T+4) + 2T + 7)$, where $F$ is the field count, after which the recorder, started from the prepared blank seam, holds on its record tape exactly the stored rows of times $0,\dots,j-1$, carries the movement counts and source configuration of time $j$ with a valid tag, and is ready to store the current row. The invariant concerns actual tape contents, not merely a running schedule.
\end{lemma}

\begin{lemma}[The recorder stores the full trajectory]
\lean{Complexity.clRec_complete}
\label{Complexity.clRec_complete}
\leanok
\uses{Complexity.clCounts, Complexity.clCounts_bound, Complexity.clElapsed_width, Complexity.clRecCfg, Complexity.clRecFields, Complexity.clRecTM, Complexity.clRec_copy, Complexity.clRec_prefix, Complexity.clRec_stop, Complexity.clRecords, Turing.FinTM, Turing.MultiTapeTM.initCfg, Turing.MultiTapeTM.runFrom, Turing.MultiTapeTM.runFrom_add, Turing.MultiTapeTM.runFrom_succ_eq_step'}
\difficulty{4}
For every source machine, physical input, virtual input $y$, and horizon $T$, there is a time $s \le (T+1)(F(3T+4) + 2T + 7)$, where $F$ is the field count, after which the recorder, started from the prepared blank seam, is in its completed live control with the stored rows of all times $0,\dots,T$ on the record tape, the movement counts and the exact horizon configuration of the source, and an empty physical output.
\end{lemma}

\begin{lemma}[Prepared initialization of the recorder]
\lean{Complexity.clRec_prepared}
\label{Complexity.clRec_prepared}
\leanok
\uses{Complexity.clRecCfg, Complexity.clRecFields, Complexity.clRecTM, Turing.Cfg.init, Turing.Cfg.ofWords, Turing.FinTM, Turing.FinTM.tapeBlocks, Turing.MultiTapeTM.initCfg}
\difficulty{3}
For every source machine, physical input, virtual input $y$, and horizon $T$, the configuration assembled from words, blank movement counters, empty record, a unary clock of $T$ true cells, the virtual input, and blank source work tapes, equals the recorder's prepared seam at logical time zero over the initial source configuration. No claim is made about native initialization from raw input.
\end{lemma}

\begin{lemma}[Polynomial ledger of the recording work]
\lean{Complexity.clRec_cost_bound}
\label{Complexity.clRec_cost_bound}
\leanok
\uses{Complexity.clRecFields, Turing.FinTM}
\difficulty{3}
For every source machine and horizon $T$, the complete stored-row ledger $(T+1)(F(3T+4) + 2T + 7)$, where $F$ is the field count, is at most $(4F + 7)(T+1)^2$, so native recording work is polynomial in the source horizon, separately from the formula output size.
\end{lemma}

\begin{definition}[Numeric value of a Boolean digit]
\lean{Complexity.clBit}
\label{Complexity.clBit}
\leanok
A Boolean digit has numeric value one when true and zero when false.
\end{definition}

\begin{definition}[Little-endian value of a binary word]
\lean{Complexity.clNum}
\label{Complexity.clNum}
\leanok
The value of an arbitrary little-endian bit list is obtained by folding the binary digit operation over its bits, so padded words with trailing zeros are included.
\end{definition}

\begin{lemma}[Canonical words decode to their values]
\lean{Complexity.clNum_bits}
\label{Complexity.clNum_bits}
\leanok
\uses{Complexity.clNum}
\difficulty{3}
For every natural number $n$, the little-endian value of the binary representation of $n$ is $n$ itself.
\end{lemma}

\begin{lemma}[Low digit and high digits]
\lean{Complexity.clNum_head}
\label{Complexity.clNum_head}
\leanok
\uses{Complexity.clBit, Complexity.clNum}
\difficulty{2}
For every bit list $w$, the value of $w$ equals twice the value of its tail plus the numeric value of its first digit, where an empty word contributes a zero digit.
\end{lemma}

\begin{definition}[One binary addition column]
\lean{Complexity.clAddColumn}
\label{Complexity.clAddColumn}
\leanok
For a carry bit and two addend bits, the addition column's low result bit is the exclusive or of the three inputs, and its outgoing carry is the majority of the three inputs.
\end{definition}

\begin{lemma}[The addition table is exact]
\lean{Complexity.clAddColumn_value}
\label{Complexity.clAddColumn_value}
\leanok
\uses{Complexity.clAddColumn, Complexity.clBit}
\difficulty{2}
For every carry bit $c$ and addend bits $u$ and $v$, the sum of the three digit values equals twice the outgoing carry's value plus the low result bit's value.
\end{lemma}

\begin{definition}[Comparator register update]
\lean{Complexity.clCmpUpdate}
\label{Complexity.clCmpUpdate}
\leanok
\uses{Complexity.clAddColumn}
Given three Boolean registers, two running carries and an equality flag, and the four scanned column bits, the updated registers are obtained by feeding each pair of bits through its own addition column: the two outgoing carries are stored, and the equality flag is conjoined with the comparison of the two low result bits.
\end{definition}

\begin{definition}[Comparator invariant]
\lean{Complexity.clCmpPred}
\label{Complexity.clCmpPred}
\leanok
\uses{Complexity.clBit, Complexity.clNum}
For comparator registers and four remaining high words, the invariant states that the equality flag is still set and that the value of the first pair of words plus the first carry equals the value of the second pair plus the second carry, given that all earlier low digits matched.
\end{definition}

\begin{lemma}[One column preserves the tested equality]
\lean{Complexity.clCmpUpdate_spec}
\label{Complexity.clCmpUpdate_spec}
\leanok
\uses{Complexity.clAddColumn, Complexity.clAddColumn_value, Complexity.clBit, Complexity.clCmpPred, Complexity.clCmpUpdate, Complexity.clNum, Complexity.clNum_head}
\difficulty{5}
For every comparator registers $q$ and four words, the comparator invariant on the full words holds exactly when it holds for the updated registers on the four tails: equal sums are equivalent to equal low bits together with equal remaining high sums, and unequal low bits have opposite parity, which no later column can repair.
\end{lemma}

\begin{definition}[Comparator control after scanned columns]
\lean{Complexity.clCmpOrbit}
\label{Complexity.clCmpOrbit}
\leanok
\uses{Complexity.clCmpUpdate}
For four words, the comparator control starts with clear carries and a set equality flag, and after $j+1$ columns is the register update applied to the control after $j$ columns, reading the four bits at position $j$ with blanks treated as absent.
\end{definition}

\begin{lemma}[The scanned invariant tests the equation]
\lean{Complexity.clCmpOrbit_spec}
\label{Complexity.clCmpOrbit_spec}
\leanok
\uses{Complexity.clBit, Complexity.clCmpOrbit, Complexity.clCmpPred, Complexity.clCmpUpdate_spec, Complexity.clNum}
\difficulty{3}
For every four words and every column count $j$, the comparator invariant at the control after $j$ columns, applied to the words with their first $j$ bits dropped, holds exactly when the value of the first word plus the second equals the value of the third plus the fourth; columns beyond a short word read zeros, never uncharged random access.
\end{lemma}

\begin{definition}[Maximum width of the compared words]
\lean{Complexity.clCmpSize}
\label{Complexity.clCmpSize}
\leanok
For four input words, the comparator size is the maximum of their four lengths; it is used only in the runtime accounting.
\end{definition}

\begin{lemma}[Each word fits the common scan width]
\lean{Complexity.clCmp_width}
\label{Complexity.clCmp_width}
\leanok
\uses{Complexity.clCmpSize}
\difficulty{1}
For every four comparator words and every index, the length of the chosen word is at most the maximum of the four lengths.
\end{lemma}

\begin{lemma}[Joint blanks mark the common boundary]
\lean{Complexity.clCmp_blank}
\label{Complexity.clCmp_blank}
\leanok
\uses{Complexity.clCmpSize}
\difficulty{1}
For every four comparator words and every position $j$, all four aligned reads at $j$ are blank exactly when the common maximum width is at most $j$, so a word attaining the maximum prevents an early return.
\end{lemma}

\begin{definition}[Final comparator verdict]
\lean{Complexity.clCmpVerdict}
\label{Complexity.clCmpVerdict}
\leanok
At the common right boundary, the verdict is true exactly when the equality flag is still set and the two final carry bits agree.
\end{definition}

\begin{lemma}[Digit values determine the digits]
\lean{Complexity.clBit_eq}
\label{Complexity.clBit_eq}
\leanok
\uses{Complexity.clBit}
\difficulty{2}
For all Boolean digits $a$ and $b$, their numeric values are equal exactly when the digits themselves are equal.
\end{lemma}

\begin{lemma}[The verdict decides the sum equation]
\lean{Complexity.clCmpVerdict_spec}
\label{Complexity.clCmpVerdict_spec}
\leanok
\uses{Complexity.clBit_eq, Complexity.clCmpOrbit, Complexity.clCmpOrbit_spec, Complexity.clCmpPred, Complexity.clCmpSize, Complexity.clCmpVerdict, Complexity.clCmp_width, Complexity.clNum}
\difficulty{3}
For every four words, the verdict computed from the comparator control after scanning the common maximum width is true exactly when the value of the first word plus the second equals the value of the third plus the fourth.
\end{lemma}

\begin{definition}[Native four-word comparator]
\lean{Complexity.clCmpTM}
\label{Complexity.clCmpTM}
\leanok
\uses{Complexity.clCmpUpdate, Complexity.clCmpVerdict, Turing.FinTM, Turing.FinTM.controlAction}
A machine with four work tapes scans four read-only words in lockstep, feeding each column through the finite addition tables, then rewinds all four heads to zero and idles in a live control carrying the Boolean verdict. No word is modified and nothing is physically emitted.
\end{definition}

\begin{definition}[Comparator configuration]
\lean{Complexity.clCmpCfg}
\label{Complexity.clCmpCfg}
\leanok
\uses{Complexity.clCmpTM, Turing.Cfg, Turing.FinTM.bufferTape}
For an input word, an input position, a comparator control, four stored words, and a common head position, the comparator configuration holds each word on its own buffer tape with all four heads aligned and empty output.
\end{definition}

\begin{lemma}[Aligned reads are the current column]
\lean{Complexity.clCmp_read}
\label{Complexity.clCmp_read}
\leanok
\uses{Complexity.clCmpCfg, Complexity.clCmpTM, Turing.Cfg.workTapeSymbols}
\difficulty{1}
For every comparator configuration with all four heads at a nonnegative aligned position $j$, the native reads are exactly the bits of the four words at position $j$, with blanks past each word's end.
\end{lemma}

\begin{lemma}[One forward step is one column]
\lean{Complexity.clCmp_forward_step}
\label{Complexity.clCmp_forward_step}
\leanok
\uses{Complexity.clCmpCfg, Complexity.clCmpOrbit, Complexity.clCmpSize, Complexity.clCmpTM, Complexity.clCmp_blank, Complexity.clCmp_read, Turing.Action.apply, Turing.MultiTapeTM.step, Turing.moveInputPos_zero}
\difficulty{3}
For every input word, input position, four words, and position $j$ strictly below the common width, one comparator step from the scanning control after $j$ columns yields exactly the next column: the control is the orbit value after $j+1$ columns and all four heads have moved right.
\end{lemma}

\begin{lemma}[The forward scan is a native run]
\lean{Complexity.clCmp_forward}
\label{Complexity.clCmp_forward}
\leanok
\uses{Complexity.clCmpCfg, Complexity.clCmpOrbit, Complexity.clCmpSize, Complexity.clCmpTM, Complexity.clCmp_forward_step, Turing.MultiTapeTM.runFrom, Turing.MultiTapeTM.runFrom_succ_eq_step'}
\difficulty{3}
For every input word, input position, four words, and $j$ at most the common width, running the comparator for $j$ steps from the start reaches the scanning control after $j$ columns with all heads at $j$; every prefix of the forward scan is an actual aligned native run.
\end{lemma}

\begin{lemma}[The boundary stores the verdict]
\lean{Complexity.clCmp_finish_scan}
\label{Complexity.clCmp_finish_scan}
\leanok
\uses{Complexity.clCmpCfg, Complexity.clCmpOrbit, Complexity.clCmpSize, Complexity.clCmpTM, Complexity.clCmpVerdict, Complexity.clCmp_blank, Complexity.clCmp_read, Turing.Action.apply, Turing.MultiTapeTM.step, Turing.moveInputPos_zero}
\difficulty{3}
For every input word, input position, and four words, one comparator step at the common right boundary reads the joint blank, stores the final carry verdict in finite control, and starts rewinding by moving all heads left; this includes the case where all four words are empty.
\end{lemma}

\begin{lemma}[The comparator rewinds all heads]
\lean{Complexity.clCmp_rewind}
\label{Complexity.clCmp_rewind}
\leanok
\uses{Complexity.clCmpCfg, Complexity.clCmpSize, Complexity.clCmpTM, Complexity.clCmp_blank, Complexity.clCmp_read, Turing.Action.apply, Turing.Cfg.workTapeSymbols, Turing.MultiTapeTM.runFrom, Turing.MultiTapeTM.runFrom_succ_eq_step, Turing.MultiTapeTM.runFrom_zero, Turing.MultiTapeTM.step, Turing.moveInputPos_zero}
\difficulty{4}
For every input word, input position, four words, stored verdict, and count $j$ at most the common width, running the comparator in its rewind control at aligned position $j-1$ for $j+1$ steps returns all four heads to zero and enters the completed control, preserving every word and the verdict.
\end{lemma}

\begin{lemma}[Exact comparator runtime]
\lean{Complexity.clCmp_run}
\label{Complexity.clCmp_run}
\leanok
\uses{Complexity.clCmpCfg, Complexity.clCmpOrbit, Complexity.clCmpSize, Complexity.clCmpTM, Complexity.clCmpVerdict, Complexity.clCmp_finish_scan, Complexity.clCmp_forward, Complexity.clCmp_rewind, Turing.MultiTapeTM.runFrom, Turing.MultiTapeTM.runFrom_add, Turing.MultiTapeTM.runFrom_succ_eq_step}
\difficulty{3}
For every input word, input position, and four words, running the comparator for twice the common maximum width plus two steps from the start yields the completed control carrying the orbit verdict, with every word unchanged and all heads at zero.
\end{lemma}

\begin{lemma}[The completed comparator is absorbing]
\lean{Complexity.clCmp_idle}
\label{Complexity.clCmp_idle}
\leanok
\uses{Complexity.clCmpTM, Turing.Action.apply, Turing.Cfg, Turing.FinTM.controlAction, Turing.FinTM.controlAction_apply, Turing.MultiTapeTM.step, Turing.moveInputPos_zero}
\difficulty{3}
Every comparator configuration in a completed control is left entirely unchanged by one step of the comparator.
\end{lemma}

\begin{lemma}[First completion with the arithmetic verdict]
\lean{Complexity.clCmp_first}
\label{Complexity.clCmp_first}
\leanok
\uses{Complexity.clCmpCfg, Complexity.clCmpOrbit, Complexity.clCmpSize, Complexity.clCmpTM, Complexity.clCmpVerdict, Complexity.clCmpVerdict_spec, Complexity.clCmp_idle, Complexity.clCmp_run, Complexity.clNum, Turing.MultiTapeTM.runFrom, Turing.MultiTapeTM.runFrom_add, Turing.MultiTapeTM.runFrom_zero}
\difficulty{5}
For every input word, input position, four words of length at most $W$, there is a time $t$ with $0 < t \le 2W + 2$ and a bit $b$, true exactly when the value of the first word plus the second equals the value of the third plus the fourth, such that no completed control appears before $t$ and at $t$ the comparator sits in the completed control carrying $b$ with all four words unchanged and heads at zero.
\end{lemma}

\begin{lemma}[Stored positions match the public schedule]
\lean{Complexity.clCounts_schedule}
\label{Complexity.clCounts_schedule}
\leanok
\uses{Complexity.clCounts, Complexity.clCounts_positions, Complexity.clPositions, Complexity.clSigned, Complexity.inputPosAt, Complexity.workPosAt, Turing.FinTM}
\difficulty{3}
For every oblivious machine, reference length $m$, and time $t$, the signed value of the stored input-head count pair over the all-false reference input, plus one, equals the public input position at time $t$, and the signed work count pairs equal the public work positions, including the shifted input origin.
\end{lemma}

\begin{definition}[Packing a list of fields]
\lean{Complexity.clFields}
\label{Complexity.clFields}
\leanok
\uses{Turing.pairEncode}
A list of binary words is packed by encoding each word as a self-delimiting doubled-bit field, the same separators as the native copier, concatenated in order; empty binary zero fields are preserved and no fixed width is assumed.
\end{definition}

\begin{lemma}[Appending past an encoded field]
\lean{Complexity.clPair_append}
\label{Complexity.clPair_append}
\leanok
\uses{Turing.pairEncode}
\difficulty{1}
For all words $a$, $b$, and $c$, the self-delimiting encoding of $a$ with payload $b$, followed by $c$, equals the encoding of $a$ with payload $b$ followed by $c$; the boundary of the first field is preserved.
\end{lemma}

\begin{lemma}[Packing respects concatenation]
\lean{Complexity.clFields_append}
\label{Complexity.clFields_append}
\leanok
\uses{Complexity.clFields, Complexity.clPair_append}
\difficulty{3}
For all lists of words $a$ and $b$, the packed encoding of their concatenation is the packed encoding of $a$ followed by the packed encoding of $b$.
\end{lemma}

\begin{lemma}[Stored rows are packed field lists]
\lean{Complexity.clRowPrefix_fields}
\label{Complexity.clRowPrefix_fields}
\leanok
\uses{Complexity.clFields, Complexity.clFields_append, Complexity.clRowPrefix}
\difficulty{3}
For every family of counter words on $l$ coordinates and every $j \le l$, the stored row prefix of the first $j$ fields equals the packed encoding of the first $j$ words taken in the fixed tape order.
\end{lemma}

\begin{definition}[Sequential field parser]
\lean{Complexity.clReadFields}
\label{Complexity.clReadFields}
\leanok
\uses{Turing.pairDecode}
Given a count $n$ and a bit stream, the parser repeatedly decodes one self-delimiting field, returning after $n$ fields the list of decoded words together with the exact remaining suffix; malformed streams yield no result.
\end{definition}

\begin{definition}[Native two-tape field reader]
\lean{Complexity.clReadTM}
\label{Complexity.clReadTM}
\leanok
\uses{Complexity.clTwo, Turing.FinTM, Turing.FinTM.controlAction}
A machine with a read-only stream tape and a target tape consumes one aligned doubled-bit field: it reads each stream bit into finite control, checks the repeated bit, overwrites one target cell per decoded bit, detects the false-true separator, rewinds only the target head, and idles in a completed control. Its contract assumes the old target is no longer than the decoded field, as holds for movement counts read in increasing source-time order.
\end{definition}

\begin{definition}[Stream-reader configuration]
\lean{Complexity.clReadCfg}
\label{Complexity.clReadCfg}
\leanok
\uses{Complexity.clReadTM, Complexity.clTwo, Turing.Cfg, Turing.FinTM.bufferTape}
For an input word, an input position, a reader control, a stream word, a target word, and independent stream and target cursors, the reader configuration holds the stream and target as buffers with the given heads; physical input and output are fixed.
\end{definition}

\begin{lemma}[Two reads decode one doubled bit]
\lean{Complexity.clRead_pair}
\label{Complexity.clRead_pair}
\leanok
\uses{Complexity.clCountTape_eq, Complexity.clCountTape_write, Complexity.clReadCfg, Complexity.clReadTM, Complexity.clTwo, Turing.Action.apply, Turing.Cfg.workTapeSymbols, Turing.FinTM.bufferTape, Turing.FinTM.bufferTape_nat, Turing.MultiTapeTM.runFrom, Turing.MultiTapeTM.runFrom_succ_eq_step', Turing.MultiTapeTM.runFrom_zero, Turing.MultiTapeTM.step, Turing.moveInputPos_zero}
\difficulty{4}
For every input word, input position, stream consisting of a prefix $u$ followed by a doubled bit and a remainder, and target split as a written prefix and an old suffix, two reader steps from the scanning control consume the two equal stream bits, overwrite the first old target cell with the decoded bit, and advance the stream cursor by two and the target cursor by one, using real adjacent head moves.
\end{lemma}

\begin{lemma}[The forward parse decodes the payload]
\lean{Complexity.clRead_forward}
\label{Complexity.clRead_forward}
\leanok
\uses{Complexity.clReadCfg, Complexity.clReadTM, Complexity.clRead_pair, Turing.MultiTapeTM.runFrom, Turing.MultiTapeTM.runFrom_add, Turing.pairEncode}
\difficulty{4}
For every input word, input position, field word $w$, stream prefix, stream tail, written target prefix, and old suffix, running the reader for $2|w|$ steps decodes all of $w$ into the target, one overwritten cell per decoded bit, preserves the exact unread stream suffix, and leaves the separator for the following steps; both stream bits of each pair are charged individually.
\end{lemma}

\begin{lemma}[Consuming the field separator]
\lean{Complexity.clRead_separator}
\label{Complexity.clRead_separator}
\leanok
\uses{Complexity.clReadCfg, Complexity.clReadTM, Complexity.clTwo, Turing.Action.apply, Turing.Cfg.workTapeSymbols, Turing.FinTM.bufferTape, Turing.FinTM.bufferTape_nat, Turing.MultiTapeTM.runFrom, Turing.MultiTapeTM.runFrom_succ_eq_step', Turing.MultiTapeTM.runFrom_zero, Turing.MultiTapeTM.step, Turing.moveInputPos_zero}
\difficulty{4}
For every input word, input position, stream prefix followed by the false-true separator and a tail, and target word, two reader steps yield the rewinding control, with the separator consumed, the stream cursor advanced past it to the next field, and the target head moved once left, without writing the separator and without changing the stream or target words.
\end{lemma}

\begin{lemma}[The reader rewinds only the target]
\lean{Complexity.clRead_rewind}
\label{Complexity.clRead_rewind}
\leanok
\uses{Complexity.clReadCfg, Complexity.clReadTM, Complexity.clTwo, Turing.Action.apply, Turing.Cfg.workTapeSymbols, Turing.FinTM.bufferTape, Turing.FinTM.bufferTape_nat, Turing.MultiTapeTM.runFrom, Turing.MultiTapeTM.runFrom_succ_eq_step, Turing.MultiTapeTM.runFrom_zero, Turing.MultiTapeTM.step, Turing.moveInputPos_zero}
\difficulty{4}
For every input word, input position, stream, decoded target $w$, stream cursor, and count $j \le |w|$, running the reader in its rewind control at target cell $j-1$ for $j+1$ steps returns the target head to zero and enters the completed control; the stream cursor never moves during the rewind.
\end{lemma}

\begin{lemma}[Complete native parse of one field]
\lean{Complexity.clRead_run}
\label{Complexity.clRead_run}
\leanok
\uses{Complexity.clReadCfg, Complexity.clReadTM, Complexity.clRead_forward, Complexity.clRead_rewind, Complexity.clRead_separator, Turing.MultiTapeTM.runFrom, Turing.MultiTapeTM.runFrom_add, Turing.pairEncode}
\difficulty{4}
For every input word, input position, stream prefix $u$, field word $w$, stream tail, and old target of length at most $|w|$, running the reader for $3|w| + 3$ steps from the scanning control at stream cursor $|u|$ yields the completed control, in which exactly one encoded field is consumed, the target is $w$ with its head restored to zero, the whole stream is preserved, and the stream cursor sits just past the consumed field; the empty field takes three steps.
\end{lemma}

\begin{lemma}[Size of the stored trajectory]
\lean{Complexity.clRecords_length}
\label{Complexity.clRecords_length}
\leanok
\uses{Complexity.clCounts, Complexity.clCounts_bound, Complexity.clElapsed_width, Complexity.clRecFields, Complexity.clRecords, Complexity.clRowPrefix_length, Turing.FinTM}
\difficulty{4}
For every source machine, virtual input $y$, horizon $T$, and count $t \le T + 1$, the serialized trajectory prefix of the first $t$ rows has length at most $t$ times the field count times $2T + 2$: every stored count has width at most the horizon and every field is exactly a doubled payload plus its separator.
\end{lemma}

\begin{lemma}[Cutting a run at its first absorbing return]
\lean{Complexity.clFirst}
\label{Complexity.clFirst}
\leanok
\uses{Turing.Cfg, Turing.MultiTapeTM, Turing.MultiTapeTM.runFrom, Turing.MultiTapeTM.runFrom_add, Turing.MultiTapeTM.runFrom_zero, Turing.MultiTapeTM.step}
\difficulty{4}
Let a machine run from a start configuration to a finish configuration in $B$ steps, where the finish control is a distinguished absorbing return state, every configuration in that state is fixed by one step, and the start is not in the return state. Then there is a strictly positive time $t \le B$ at which the run first reaches the return state, and the configuration at $t$ is already the exact finish configuration.
\end{lemma}

\begin{lemma}[Guarded state embedding on unchanged tapes]
\lean{Complexity.clMap_run}
\label{Complexity.clMap_run}
\leanok
\uses{Complexity.clSlotCfg, Complexity.clSlot_run, Turing.Action.mapState, Turing.Cfg, Turing.Cfg.mapState, Turing.MultiTapeTM, Turing.MultiTapeTM.runFrom}
\difficulty{2}
Suppose a host machine with the same tapes agrees, on every embedded control satisfying a guard, with a smaller machine's transition after relabeling states. Then for every start configuration and horizon $t$ whose strictly earlier controls all satisfy the guard, the host run equals the state-relabeled run of the smaller machine; all tape data is preserved exactly.
\end{lemma}

\begin{lemma}[Clearing the last occupied cell]
\lean{Complexity.clErase_last}
\label{Complexity.clErase_last}
\leanok
\uses{Turing.FinTM.bufferTape, Turing.FinTM.bufferTape_append, Turing.FinTM.bufferTape_nat}
\difficulty{3}
For every word $w$ and bit $b$, updating the buffer tape of $w$ extended by $b$ at cell $|w|$ with a blank yields exactly the buffer tape of $w$.
\end{lemma}

\begin{definition}[Two-tape scratch reset machine]
\lean{Complexity.clWipeTM}
\label{Complexity.clWipeTM}
\leanok
\uses{Complexity.clTwo, Turing.FinTM, Turing.FinTM.controlAction}
A three-state machine with a stream tape and a target tape resets its scratch target: it scans the target to its right blank, erases every occupied cell on the way back to the left blank, and returns at target cell zero in an absorbing control. The stream tape and its cursor are never read or moved.
\end{definition}

\begin{definition}[Reset seam configuration]
\lean{Complexity.clWipeCfg}
\label{Complexity.clWipeCfg}
\leanok
\uses{Complexity.clTwo, Turing.Cfg, Turing.FinTM.bufferTape}
For an input word, an input position, a reset control, arbitrary preserved stream contents with a stream head, a target word, and a target head, the reset configuration holds the stream and target as buffers with empty physical output.
\end{definition}

\begin{lemma}[Forward scratch scan]
\lean{Complexity.clWipe_forward}
\label{Complexity.clWipe_forward}
\leanok
\uses{Complexity.clTwo, Complexity.clWipeCfg, Complexity.clWipeTM, Turing.Action.apply, Turing.Cfg.workTapeSymbols, Turing.FinTM.bufferTape, Turing.FinTM.bufferTape_nat, Turing.MultiTapeTM.runFrom, Turing.MultiTapeTM.runFrom_succ_eq_step, Turing.MultiTapeTM.runFrom_zero, Turing.MultiTapeTM.step, Turing.moveInputPos_zero}
\difficulty{4}
For every input word, input position, stream, target word $w$, stream head, and start cell $j \le |w|$, running the reset machine for $|w| - j + 1$ steps from its scanning control moves the target head right across each occupied cell, turns at the right blank, and enters the erasing control at cell $|w| - 1$; the target is still unchanged.
\end{lemma}

\begin{lemma}[Backward erasing pass]
\lean{Complexity.clWipe_backward}
\label{Complexity.clWipe_backward}
\leanok
\uses{Complexity.clErase_last, Complexity.clTwo, Complexity.clWipeCfg, Complexity.clWipeTM, Turing.Action.apply, Turing.Cfg.workTapeSymbols, Turing.FinTM.bufferTape, Turing.FinTM.bufferTape_nat, Turing.MultiTapeTM.runFrom, Turing.MultiTapeTM.runFrom_succ_eq_step, Turing.MultiTapeTM.runFrom_zero, Turing.MultiTapeTM.step, Turing.moveInputPos_zero}
\difficulty{4}
For every input word, input position, stream, stream head, and target word $w$, running the reset machine in its erasing control at cell $|w| - 1$ for $|w| + 1$ steps erases every occupied target cell, including the first one, and returns the target head to zero in the completed control; the stream is neither read nor modified.
\end{lemma}

\begin{lemma}[Exact reset cost]
\lean{Complexity.clWipe_run}
\label{Complexity.clWipe_run}
\leanok
\uses{Complexity.clWipeCfg, Complexity.clWipeTM, Complexity.clWipe_backward, Complexity.clWipe_forward, Turing.MultiTapeTM.runFrom, Turing.MultiTapeTM.runFrom_add}
\difficulty{3}
For every input word, input position, stream, target word $w$, and stream head, running the reset machine for $2|w| + 2$ steps from the scanning control at target cell zero leaves an empty target with its head at zero, in the completed control; both blank transitions are included and no bound on any later field is required.
\end{lemma}

\begin{lemma}[The completed reset is absorbing]
\lean{Complexity.clWipe_idle}
\label{Complexity.clWipe_idle}
\leanok
\uses{Complexity.clWipeTM, Turing.Action.apply, Turing.Cfg, Turing.FinTM.controlAction, Turing.FinTM.controlAction_apply, Turing.MultiTapeTM.step, Turing.moveInputPos_zero}
\difficulty{3}
Every reset-machine configuration in the completed control is left entirely unchanged by one step.
\end{lemma}

\begin{lemma}[First completion of a scratch reset]
\lean{Complexity.clWipe_first}
\label{Complexity.clWipe_first}
\leanok
\uses{Complexity.clFirst, Complexity.clWipeCfg, Complexity.clWipeTM, Complexity.clWipe_idle, Complexity.clWipe_run, Turing.MultiTapeTM.runFrom}
\difficulty{1}
For every input word, input position, stream, target word $w$, and stream head, there is a time $t$ with $0 < t \le 2|w| + 2$ such that the reset machine first reaches its completed control exactly at $t$, with the target erased, the target head at zero, and the stream untouched.
\end{lemma}

\begin{definition}[Unrestricted field loader]
\lean{Complexity.clFreshTM}
\label{Complexity.clFreshTM}
\leanok
\uses{Complexity.clReadTM, Complexity.clWipeTM, Turing.Action.mapState, Turing.FinTM, Turing.FinTM.controlAction}
A two-tape machine first runs the scratch reset on its target, then dispatches once into the banked field reader on the same tapes: the old target is actually cleared before parsing, so the new field may be empty, shorter, or unrelated to the old contents. Stream positions and all reset and read transitions are charged.
\end{definition}

\begin{lemma}[Reset-then-read whole-configuration run]
\lean{Complexity.clFresh_run}
\label{Complexity.clFresh_run}
\leanok
\uses{Complexity.clFreshTM, Complexity.clMap_run, Complexity.clReadCfg, Complexity.clReadTM, Complexity.clRead_run, Complexity.clWipeCfg, Complexity.clWipeTM, Complexity.clWipe_first, Turing.Action.apply, Turing.Cfg.mapState, Turing.FinTM.controlAction, Turing.FinTM.controlAction_apply, Turing.MultiTapeTM.runFrom, Turing.MultiTapeTM.runFrom_add, Turing.MultiTapeTM.runFrom_succ_eq_step', Turing.MultiTapeTM.step, Turing.moveInputPos_zero, Turing.pairEncode}
\difficulty{5}
For every input word, input position, stream prefix $u$, field word $w$, stream tail, and old target, there is a time $t \le 2\,(\text{old length}) + 3|w| + 6$ after which the unrestricted loader, started at the reset phase with stream cursor $|u|$, holds exactly $w$ on its target with head zero and stream cursor just past the consumed field; the native input and physical output are unchanged, and the stream moves only during reading.
\end{lemma}

\begin{lemma}[The completed loader is absorbing]
\lean{Complexity.clFresh_idle}
\label{Complexity.clFresh_idle}
\leanok
\uses{Complexity.clFreshTM, Turing.Action.apply, Turing.Cfg, Turing.FinTM.controlAction, Turing.FinTM.controlAction_apply, Turing.MultiTapeTM.step, Turing.moveInputPos_zero}
\difficulty{3}
Every configuration of the unrestricted loader in its completed reading control is left entirely unchanged by one step.
\end{lemma}

\begin{lemma}[First return of the unrestricted loader]
\lean{Complexity.clFresh_first}
\label{Complexity.clFresh_first}
\leanok
\uses{Complexity.clFirst, Complexity.clFreshTM, Complexity.clFresh_idle, Complexity.clFresh_run, Complexity.clReadCfg, Complexity.clWipeCfg, Turing.Cfg.mapState, Turing.MultiTapeTM.runFrom, Turing.pairEncode}
\difficulty{2}
For every input word, input position, stream prefix, field word $w$, stream tail, and old target, there is a strictly positive time, at most twice the old target length plus $3|w| + 6$, at which the unrestricted loader first reaches its completed control, holding exactly $w$ with restored target head and advanced stream cursor; no bound relating the old and new lengths is assumed.
\end{lemma}

\begin{definition}[Tape slots of one row-load call]
\lean{Complexity.clLoadIndex}
\label{Complexity.clLoadIndex}
\leanok
\uses{Complexity.clTwo}
For a target field coordinate $i$ among $l$ fields, the loader's two tapes are embedded with the shared stream first, the final extra tape, and the target field second, counter tape $i$.
\end{definition}

\begin{definition}[Partial inverse of the row-load slots]
\lean{Complexity.clLoadSelect}
\label{Complexity.clLoadSelect}
\leanok
For a target field coordinate $i$, the selection maps field tape $i$ back to the loader's target slot and the final stream tape back to its stream slot, leaving every other field tape unselected.
\end{definition}

\begin{lemma}[Row-load selection inverts the injection]
\lean{Complexity.clLoad_inverse}
\label{Complexity.clLoad_inverse}
\leanok
\uses{Complexity.clLoadIndex, Complexity.clLoadSelect, Complexity.clTwo}
\difficulty{2}
For every target field coordinate $i$ and each of the two local loader tapes, the row-load selection applied to the embedded slot returns that local tape, and the two selected tapes are physically distinct.
\end{lemma}

\begin{definition}[Sequential row loader]
\lean{Complexity.clLoadTM}
\label{Complexity.clLoadTM}
\leanok
\uses{Complexity.clFreshTM, Complexity.clLoadIndex, Complexity.clLoadSelect, Complexity.clSlotAction, Turing.FinTM, Turing.FinTM.controlAction}
For $l$ field tapes plus one shared stream tape, the sequential loader is a controller that loads a fixed row in increasing field order: for each field it runs the unrestricted loader, reset followed by parse, on that field and the stream, dispatching once after each completed field; every field clears its own target first, so repeated queries need no monotone-width assumption.
\end{definition}

\begin{definition}[Row-loader seam configuration]
\lean{Complexity.clLoadCfg}
\label{Complexity.clLoadCfg}
\leanok
\uses{Complexity.clFreshTM, Complexity.clLoadTM, Turing.Cfg, Turing.FinTM.bufferTape}
For an input word, an input position, a field index, a loader control, field words, a stream, and a stream cursor, the row-loader seam holds every field as a buffer at head zero and the stream with its explicit cursor; the physical output is empty.
\end{definition}

\begin{lemma}[Inactive frame of one field call]
\lean{Complexity.clLoad_frame}
\label{Complexity.clLoad_frame}
\leanok
\uses{Complexity.clFreshTM, Complexity.clLoadCfg, Complexity.clLoadSelect, Complexity.clSlotCfg, Complexity.clTwo, Turing.Cfg, Turing.FinTM.bufferTape}
\difficulty{3}
For every input word, input position, target field $i$, loader control, old field words, stream, replacement word, and cursor, relocating the two-tape loader configuration into the row loader's slots, retaining all other fields, yields the row seam in which only field $i$ has been replaced; every other field remains untouched.
\end{lemma}

\begin{lemma}[One charged row-field load]
\lean{Complexity.clLoad_field}
\label{Complexity.clLoad_field}
\leanok
\uses{Complexity.clFreshTM, Complexity.clFresh_first, Complexity.clLoadCfg, Complexity.clLoadIndex, Complexity.clLoadSelect, Complexity.clLoadTM, Complexity.clLoad_frame, Complexity.clLoad_inverse, Complexity.clSlot_run, Complexity.clWipeCfg, Turing.Action.apply, Turing.Cfg.mapState, Turing.FinTM.controlAction_apply, Turing.MultiTapeTM.runFrom, Turing.MultiTapeTM.runFrom_succ_eq_step', Turing.moveInputPos_zero, Turing.pairEncode}
\difficulty{5}
For every input word, input position, target field $i$, old field words, stream prefix $u$, field word $w$, and stream tail, there is a strictly positive time $t$, at most twice the old field's length plus $3|w| + 7$, after which the row loader has replaced field $i$ by $w$, advanced its field index to $i + 1$, and moved the stream cursor exactly past the consumed field, from $|u|$ to $|u| + 2|w| + 2$; the cost includes target reset, native parsing, target rewind, and one dispatch.
\end{lemma}

\begin{definition}[Partially loaded row bank]
\lean{Complexity.clLoadWords}
\label{Complexity.clLoadWords}
\leanok
For old and new field words on $l$ coordinates and a count $j$, the intermediate bank holds the new word at every coordinate below $j$ and the old word elsewhere.
\end{definition}

\begin{lemma}[Loading the next field extends the prefix]
\lean{Complexity.clLoadWords_step}
\label{Complexity.clLoadWords_step}
\leanok
\uses{Complexity.clLoadWords}
\difficulty{3}
For all old and new field words and every coordinate $i$, replacing field $i$ of the bank in which the first $i$ fields are new yields the bank in which the first $i + 1$ fields are new.
\end{lemma}

\begin{lemma}[Exposing the next encoded field]
\lean{Complexity.clLoad_split}
\label{Complexity.clLoad_split}
\leanok
\uses{Complexity.clFields, Complexity.clFields_append, Complexity.clPair_append, Turing.pairEncode}
\difficulty{4}
For every row of field words, coordinate $i$, stream prefix, and trailer, the stream consisting of the prefix, the packed encoding of the whole row, and the trailer equals the prefix extended by the packed first $i$ fields, followed by the self-delimiting encoding of field $i$, whose payload is the packed remaining fields and the trailer. The suffix is preserved literally.
\end{lemma}

\begin{lemma}[Charged loading of a row prefix]
\lean{Complexity.clLoad_prefix}
\label{Complexity.clLoad_prefix}
\leanok
\uses{Complexity.clFields, Complexity.clFields_append, Complexity.clLoadCfg, Complexity.clLoadTM, Complexity.clLoadWords, Complexity.clLoadWords_step, Complexity.clLoad_field, Complexity.clLoad_split, Turing.MultiTapeTM.runFrom, Turing.MultiTapeTM.runFrom_add, Turing.MultiTapeTM.runFrom_zero, Turing.pairEncode}
\difficulty{5}
For every input word, input position, old fields of length at most $U$, new fields of length at most $W$, stream prefix and trailer, and count $j \le l$, there is a time $t \le j\,(2U + 3W + 7)$ after which the row loader, started at field zero with the stream holding the packed new row, has loaded exactly the first $j$ new fields, restored every field head to zero, and advanced the stream cursor by the exact length of the packed first $j$ fields. Every field is reset even when revisiting a short row after a longer one.
\end{lemma}

\begin{lemma}[Full sequential row load]
\lean{Complexity.clLoad_complete}
\label{Complexity.clLoad_complete}
\leanok
\uses{Complexity.clFields, Complexity.clLoadCfg, Complexity.clLoadTM, Complexity.clLoadWords, Complexity.clLoad_prefix, Turing.MultiTapeTM.runFrom}
\difficulty{2}
For every input word, input position, old fields of length at most $U$, new fields of length at most $W$, and stream prefix and trailer, there is a time $t \le l\,(2U + 3W + 7)$ after which the row loader holds exactly the new row on its field tapes, with every field head restored, the full record stream retained, and the stream cursor immediately after the packed row; the load is a native call, not indexed list access.
\end{lemma}

\begin{lemma}[The completed row loader is absorbing]
\lean{Complexity.clLoad_idle}
\label{Complexity.clLoad_idle}
\leanok
\uses{Complexity.clLoadTM, Turing.Cfg, Turing.FinTM.controlAction_apply, Turing.MultiTapeTM.step, Turing.moveInputPos_zero}
\difficulty{3}
Every row-loader configuration whose control is the final field index with a fresh loader control is left entirely unchanged by one step, including the empty-row case.
\end{lemma}

\begin{lemma}[First return of a nonempty row load]
\lean{Complexity.clLoad_first}
\label{Complexity.clLoad_first}
\leanok
\uses{Complexity.clFields, Complexity.clFirst, Complexity.clLoadCfg, Complexity.clLoadTM, Complexity.clLoad_complete, Complexity.clLoad_idle, Turing.MultiTapeTM.runFrom}
\difficulty{3}
For a positive field count, every input word, input position, old fields of length at most $U$, new fields of length at most $W$, and stream prefix and trailer, there is a strictly positive time $t \le l\,(2U + 3W + 7)$ at which the row loader first reaches its complete-row return, holding the new row with all heads restored and the exact sequential stream cursor.
\end{lemma}

\begin{definition}[Native input copier]
\lean{Complexity.clInputTM}
\label{Complexity.clInputTM}
\leanok
\uses{Turing.FinTM, Turing.FinTM.controlAction}
A three-state machine with one work tape copies its entire native input onto the work tape, rewinds the input head against the physical left boundary with the destination head following exactly, and idles in a silent live return control.
\end{definition}

\begin{definition}[Input-copy configuration]
\lean{Complexity.clInputCfg}
\label{Complexity.clInputCfg}
\leanok
\uses{Turing.Cfg, Turing.FinTM.bufferTape}
For an input word, a copier control, an input position, a copied prefix, and a work head, the copy configuration holds the copied prefix as a buffer with aligned heads and empty output.
\end{definition}

\begin{lemma}[Forward native input copy]
\lean{Complexity.clInput_forward}
\label{Complexity.clInput_forward}
\leanok
\uses{Complexity.clInputCfg, Complexity.clInputTM, Turing.Action.apply, Turing.Cfg.inputSymbol, Turing.FinTM.bufferTape, Turing.FinTM.bufferTape_append, Turing.MultiTapeTM.runFrom, Turing.MultiTapeTM.runFrom_succ_eq_step', Turing.inputSymbolInner, Turing.moveInputPos_pos_of_ne_right}
\difficulty{4}
For every input word $x$ and every $j \le |x|$, running the input copier for $j$ steps from its start configuration copies exactly the first $j$ bits of $x$ onto the work buffer, with the input head at position $j + 1$ and the work head at $j$; every source read and target write is charged.
\end{lemma}

\begin{lemma}[Rewinding both aligned heads]
\lean{Complexity.clInput_backward}
\label{Complexity.clInput_backward}
\leanok
\uses{Complexity.clInputCfg, Complexity.clInputTM, Turing.Action.apply, Turing.Cfg.inputSymbol, Turing.FinTM.moveInputPos_neg_val, Turing.MultiTapeTM.runFrom, Turing.MultiTapeTM.runFrom_succ_eq_step, Turing.MultiTapeTM.runFrom_zero, Turing.MultiTapeTM.step, Turing.inputSymbolInner, Turing.moveInputPos}
\difficulty{4}
For every input word $x$, copied word $w$, and prefix length $j \le |x|$, running the copier in its rewind control with the input head at $j$ and work head at $j - 1$ for $j + 1$ steps returns both heads to their origins and enters the return control; the physical left boundary identifies the origin without an unbounded control counter.
\end{lemma}

\begin{lemma}[Genuine retained-input initialization]
\lean{Complexity.clInput_run}
\label{Complexity.clInput_run}
\leanok
\uses{Complexity.clInputCfg, Complexity.clInputTM, Complexity.clInput_backward, Complexity.clInput_forward, Turing.Action.apply, Turing.Cfg.init, Turing.Cfg.inputSymbol, Turing.Cfg.ofWords, Turing.FinTM.moveInputPos_neg_val, Turing.MultiTapeTM.initCfg, Turing.MultiTapeTM.runFrom, Turing.MultiTapeTM.runFrom_add, Turing.MultiTapeTM.runFrom_succ_eq_step, Turing.MultiTapeTM.step}
\difficulty{4}
For every input word $x$, running the input copier from its native initial configuration for $2|x| + 2$ steps reaches the configuration assembled from words in which the single work tape holds exactly $x$, both heads are restored, and the physical output is empty.
\end{lemma}

\begin{lemma}[The copy return is absorbing]
\lean{Complexity.clInput_idle}
\label{Complexity.clInput_idle}
\leanok
\uses{Complexity.clInputTM, Turing.Action.apply, Turing.Cfg, Turing.FinTM.controlAction, Turing.FinTM.controlAction_apply, Turing.MultiTapeTM.step, Turing.moveInputPos_zero}
\difficulty{3}
Every input-copier configuration in the return control is left entirely unchanged by one step, so the return can be used at its first visit.
\end{lemma}

\begin{lemma}[First arrival at the retained-input seam]
\lean{Complexity.clInput_first}
\label{Complexity.clInput_first}
\leanok
\uses{Complexity.clFirst, Complexity.clInputTM, Complexity.clInput_idle, Complexity.clInput_run, Turing.Cfg.init, Turing.Cfg.ofWords, Turing.MultiTapeTM.initCfg, Turing.MultiTapeTM.runFrom}
\difficulty{1}
For every input word $x$, there is a time $t$ with $0 < t \le 2|x| + 2$ at which the input copier, started from its native initial configuration, first reaches its return control, holding exactly $x$ on its work tape with both heads restored.
\end{lemma}

\begin{definition}[Stream slot of the preparation copier]
\lean{Complexity.clPrepareIndex}
\label{Complexity.clPrepareIndex}
\leanok
During native preparation, the input copier's single work tape is embedded as the final stream slot of the field bank of $l + 1$ tapes.
\end{definition}

\begin{definition}[Partial inverse of the preparation slot]
\lean{Complexity.clPrepareSelect}
\label{Complexity.clPrepareSelect}
\leanok
The preparation selection maps the final stream slot back to the copier's only tape and leaves every future field target inactive during the input-copy phase.
\end{definition}

\begin{definition}[Preparation transducer]
\lean{Complexity.clPrepareTM}
\label{Complexity.clPrepareTM}
\leanok
\uses{Complexity.clInputTM, Complexity.clLoadTM, Complexity.clPrepareIndex, Complexity.clPrepareSelect, Complexity.clSlotAction, Turing.Action.mapState, Turing.FinTM, Turing.FinTM.controlAction}
For $l$ field tapes plus one stream tape, the preparation machine first runs the native input copier into the stream slot, then dispatches once into the sequential row loader, which parses the fixed number of self-delimiting fields from the copied stream.
\end{definition}

\begin{lemma}[Reaching the loader entry from native start]
\lean{Complexity.clPrepare_start}
\label{Complexity.clPrepare_start}
\leanok
\uses{Complexity.clInputTM, Complexity.clInput_first, Complexity.clLoadCfg, Complexity.clPrepareIndex, Complexity.clPrepareSelect, Complexity.clPrepareTM, Complexity.clSlotCfg, Complexity.clSlot_run, Turing.Action.apply, Turing.Cfg.init, Turing.Cfg.mapState, Turing.Cfg.ofWords, Turing.FinTM.controlAction, Turing.FinTM.controlAction_apply, Turing.MultiTapeTM.initCfg, Turing.MultiTapeTM.runFrom, Turing.MultiTapeTM.runFrom_succ_eq_step', Turing.MultiTapeTM.step, Turing.moveInputPos_zero}
\difficulty{5}
For every field count $l$ and input word $x$, there is a time $t \le 2|x| + 3$ after which the preparation machine, started from its own native initial configuration with no assumed work-tape words, sits at the row loader's entry: all field targets blank, the stream holding exactly $x$ with cursor zero, and both heads canonical.
\end{lemma}

\begin{lemma}[Exact field-bank preparation]
\lean{Complexity.clPrepare_complete}
\label{Complexity.clPrepare_complete}
\leanok
\uses{Complexity.clFields, Complexity.clLoadCfg, Complexity.clLoadTM, Complexity.clLoad_complete, Complexity.clMap_run, Complexity.clPrepareTM, Complexity.clPrepare_start, Turing.Cfg.mapState, Turing.MultiTapeTM.initCfg, Turing.MultiTapeTM.runFrom, Turing.MultiTapeTM.runFrom_add}
\difficulty{4}
Let a row of field words of length at most $W$ and an arbitrary trailer be given, and let the native input be the packed encoding of the row followed by the trailer. Then there is a time $t \le 2|x| + 3 + l(3W + 7)$, where $x$ is that input, after which the preparation machine holds every parsed field at head zero, retains the original encoded input on its stream tape, and sits with the cursor immediately before the trailer.
\end{lemma}

\begin{lemma}[Completed preparation is absorbing]
\lean{Complexity.clPrepare_idle}
\label{Complexity.clPrepare_idle}
\leanok
\uses{Complexity.clLoadTM, Complexity.clPrepareTM, Turing.Action.apply, Turing.Cfg, Turing.FinTM.controlAction, Turing.FinTM.controlAction_apply, Turing.MultiTapeTM.step, Turing.moveInputPos_zero}
\difficulty{3}
Every preparation-machine configuration at the loader's complete-row return is left entirely unchanged by one step.
\end{lemma}

\begin{lemma}[First arrival at the prepared seam]
\lean{Complexity.clPrepare_first}
\label{Complexity.clPrepare_first}
\leanok
\uses{Complexity.clFields, Complexity.clFirst, Complexity.clLoadCfg, Complexity.clPrepareTM, Complexity.clPrepare_complete, Complexity.clPrepare_idle, Turing.Cfg.init, Turing.Cfg.mapState, Turing.MultiTapeTM.initCfg, Turing.MultiTapeTM.runFrom}
\difficulty{2}
For every row of field words of length at most $W$ and trailer, there is a strictly positive time, within the preparation ledger $2|x| + 3 + l(3W+7)$, at which the preparation machine first reaches the prepared seam with all fields parsed; native input copying and the dispatch ensure positivity even with zero fields.
\end{lemma}

\begin{lemma}[Packing a list of native computations]
\lean{Complexity.clNative_fields}
\label{Complexity.clNative_fields}
\leanok
\uses{Complexity.PolyTimeComputable, Complexity.clFields, Complexity.clNative_linear, Complexity.clNative_pair, Turing.FinTM.computesFunInTime_const}
\difficulty{3}
If every function in a finite list of string functions is polynomial-time computable, then so is the function sending an input $x$ to the packed self-delimiting encoding of the list of their values at $x$; all repeated input scans are charged by composition.
\end{lemma}

\begin{definition}[Iterated header tail]
\lean{Complexity.clHeaderTail}
\label{Complexity.clHeaderTail}
\leanok
\uses{Turing.pairDecode}
For a count $j$ and an encoded nested-pair header, the $j$-th tail follows the payload projection $j$ times, with the empty word as default for malformed input; the projections are actual native scans.
\end{definition}

\begin{definition}[Fixed header field]
\lean{Complexity.clHeaderField}
\label{Complexity.clHeaderField}
\leanok
\uses{Complexity.clHeaderTail, Turing.pairDecode}
For a count $j$ and an encoded nested-pair header, the $j$-th field is the first component of the $j$-th tail, with the empty word as default; selection is implemented by charged scans, not free indexing.
\end{definition}

\begin{lemma}[Header tails are native]
\lean{Complexity.clHeaderTail_native}
\label{Complexity.clHeaderTail_native}
\leanok
\uses{Complexity.PolyTimeComputable, Complexity.PolyTimeComputable.comp, Complexity.clHeaderTail, Complexity.clNative_linear, Complexity.polyTimeComputable_id, Turing.FinTM.computesFunInTime_pairSnd}
\difficulty{2}
For every fixed count $j$, the function taking an encoded header to its $j$-th tail is polynomial-time computable.
\end{lemma}

\begin{lemma}[Header fields are native]
\lean{Complexity.clHeaderField_native}
\label{Complexity.clHeaderField_native}
\leanok
\uses{Complexity.PolyTimeComputable, Complexity.PolyTimeComputable.comp, Complexity.clHeaderField, Complexity.clHeaderTail_native, Complexity.clNative_linear, Turing.FinTM.computesFunInTime_pairFst}
\difficulty{1}
For every fixed count $j$, the function taking an encoded header to its $j$-th field is polynomial-time computable, with all projections charged.
\end{lemma}

\begin{definition}[Prepared work words of the recorder]
\lean{Complexity.clRecWords}
\label{Complexity.clRecWords}
\leanok
\uses{Complexity.clRecFields, Complexity.clRecTM, Turing.FinTM, Turing.FinTM.tapeBlocks}
For a source machine, virtual input $y$, and horizon $T$, the recorder's prepared work words are: empty movement counters, an empty record, a unary clock of $T$ true bits, the virtual input $y$, and blank source tapes, in the recorder's established tape layout.
\end{definition}

\begin{definition}[Four retained header fields]
\lean{Complexity.clHeaderKeep}
\label{Complexity.clHeaderKeep}
\leanok
\uses{Complexity.clHeaderField}
From an encoded preparation header, the four retained fields are its first four components: the instance, the binary certificate length, the binary reference-input length, and the binary source horizon.
\end{definition}

\begin{definition}[Recorder layout of the unpacked header]
\lean{Complexity.clHeaderLayout}
\label{Complexity.clHeaderLayout}
\leanok
\uses{Complexity.clHeaderField, Complexity.clHeaderKeep, Complexity.clHeaderTail, Complexity.clRecFields, Complexity.clRecTM, Turing.FinTM, Turing.FinTM.tapeBlocks}
From an encoded preparation header, the full layout places blank counters, an empty record, the original fifth tail as the unary clock, and the fourth field as the virtual reference input into the recorder bank, followed by the four retained arithmetic and instance fields in a disjoint bank; the clock is the original horizon, not a guessed or enlarged one.
\end{definition}

\begin{lemma}[The layout is native polynomial time]
\lean{Complexity.clHeaderLayout_native}
\label{Complexity.clHeaderLayout_native}
\leanok
\uses{Complexity.PolyTimeComputable, Complexity.clFields, Complexity.clHeaderField_native, Complexity.clHeaderKeep, Complexity.clHeaderLayout, Complexity.clHeaderTail_native, Complexity.clNative_fields, Complexity.clNative_linear, Complexity.clRecTM, Turing.FinTM, Turing.FinTM.computesFunInTime_const, Turing.FinTM.tapeBlocks}
\difficulty{4}
For every source machine, the function sending an encoded header to the packed self-delimiting encoding of its full recorder layout is polynomial-time computable: every field is a constant empty word or a fixed composition of header projections, and the number of fields depends only on the machine.
\end{lemma}

\begin{lemma}[Exact layout on the proved header]
\lean{Complexity.clHeaderLayout_exact}
\label{Complexity.clHeaderLayout_exact}
\leanok
\uses{Complexity.clHeaderField, Complexity.clHeaderKeep, Complexity.clHeaderLayout, Complexity.clHeaderTail, Complexity.clPrepHeader, Complexity.clRecWords, Turing.FinTM, Turing.pairDecode_pairEncode}
\difficulty{3}
For every source machine, parameters $C, e, c, A, d$, and instance $x$, with $Q = C(|x|+1)^e$, $m = |x| + Q$, and $T = c(A(m+1)^d+1)^2$, the recorder layout of the arithmetic header equals the recorder's prepared words over the all-false reference input of length $m$ at horizon $T$, extended by the four retained fields: the instance $x$ and the binary representations of $Q$, $m$, and $T$. The certificate length stays exact.
\end{lemma}

\begin{lemma}[The recorder argument is native]
\lean{Complexity.clRecordArgument_native}
\label{Complexity.clRecordArgument_native}
\leanok
\uses{Complexity.PolyTimeComputable, Complexity.PolyTimeComputable.comp, Complexity.clFields, Complexity.clHeaderLayout, Complexity.clHeaderLayout_native, Complexity.clPrepHeader, Complexity.clPrepHeader_native, Turing.FinTM}
\difficulty{1}
For every source machine and parameters $C, e, c, A, d$, the function sending an instance to the packed encoding of the recorder layout of its arithmetic header is polynomial-time computable.
\end{lemma}

\begin{definition}[Protected slots of the recorder host]
\lean{Complexity.clRecordSelect}
\label{Complexity.clRecordSelect}
\leanok
\uses{Complexity.clRecTM, Turing.FinTM}
In the recorder host's bank, the selection maps the first slots back to the active recorder tapes and leaves the four retained header fields and the native-input buffer unselected, protecting them from recorder actions.
\end{definition}

\begin{definition}[Injection into the recorder host]
\lean{Complexity.clRecordIndex}
\label{Complexity.clRecordIndex}
\leanok
\uses{Complexity.clRecTM, Turing.FinTM}
Each active recorder tape is embedded into the host bank before the four retained header fields and the final input-buffer slot.
\end{definition}

\begin{lemma}[Recorder selection inverts its injection]
\lean{Complexity.clRecord_inverse}
\label{Complexity.clRecord_inverse}
\leanok
\uses{Complexity.clRecTM, Complexity.clRecordIndex, Complexity.clRecordSelect, Turing.FinTM}
\difficulty{1}
For every active recorder tape coordinate, the host selection applied to the embedded coordinate returns that coordinate, with no aliasing of the protected fields.
\end{lemma}

\begin{definition}[Prepared recorder host]
\lean{Complexity.clRecordTM}
\label{Complexity.clRecordTM}
\leanok
\uses{Complexity.clPrepareTM, Complexity.clRecTM, Complexity.clRecordIndex, Complexity.clRecordSelect, Complexity.clSlotAction, Turing.Action.mapState, Turing.FinTM, Turing.FinTM.controlAction}
The host machine runs native field preparation on the full bank, dispatches once at the prepared seam into the unchanged inclusive trajectory recorder, and then runs the recorder on its active slots while the retained header fields and input buffer stay protected; the completed recorder remains live so that later search and packing phases can be attached.
\end{definition}

\begin{definition}[Retained configuration around the recorder]
\lean{Complexity.clRecordCfg}
\label{Complexity.clRecordCfg}
\leanok
\uses{Complexity.clRecTM, Complexity.clRecordSelect, Complexity.clRecordTM, Complexity.clSlotCfg, Turing.Cfg, Turing.FinTM, Turing.FinTM.bufferTape}
For an input word, four retained header fields, and a recorder configuration, the host configuration is the recorder configuration embedded on the active slots, with the four header fields kept at head zero and the original encoded input kept on the final buffer with its cursor at the input's full length.
\end{definition}

\begin{lemma}[Parsed words form the prepared seam]
\lean{Complexity.clRecord_prepare_frame}
\label{Complexity.clRecord_prepare_frame}
\leanok
\uses{Complexity.clLoadCfg, Complexity.clRecCfg, Complexity.clRecTM, Complexity.clRecWords, Complexity.clRec_prepared, Complexity.clRecordCfg, Complexity.clRecordSelect, Complexity.clRecordTM, Complexity.clSlotCfg, Turing.Cfg.mapState, Turing.Cfg.ofWords, Turing.FinTM, Turing.MultiTapeTM.initCfg}
\difficulty{4}
For every source machine, input word, virtual input $y$, horizon $T$, and retained fields, the loader's completed configuration over the parsed recorder words, with its control replaced by the recorder's start, equals the host configuration around the recorder's prepared seam at logical time zero; source initialization, counters, clock, and record bank are all fixed by the established prepared-seam identity.
\end{lemma}

\begin{lemma}[Recording from genuine native input]
\lean{Complexity.clRecord_complete}
\label{Complexity.clRecord_complete}
\leanok
\uses{Complexity.clCounts, Complexity.clFields, Complexity.clLoadCfg, Complexity.clMap_run, Complexity.clPrepareTM, Complexity.clPrepare_first, Complexity.clRecCfg, Complexity.clRecFields, Complexity.clRecTM, Complexity.clRecWords, Complexity.clRec_complete, Complexity.clRec_cost_bound, Complexity.clRecordCfg, Complexity.clRecordIndex, Complexity.clRecordSelect, Complexity.clRecordTM, Complexity.clRecord_inverse, Complexity.clRecord_prepare_frame, Complexity.clRecords, Complexity.clSlot_run, Turing.Action.apply, Turing.Cfg, Turing.Cfg.mapState, Turing.FinTM, Turing.FinTM.bufferTape, Turing.FinTM.controlAction_apply, Turing.MultiTapeTM.initCfg, Turing.MultiTapeTM.runFrom, Turing.MultiTapeTM.runFrom_add, Turing.MultiTapeTM.runFrom_succ_eq_step', Turing.MultiTapeTM.step, Turing.moveInputPos_zero}
\difficulty{5}
Let a source machine, virtual input $y$, horizon $T$, and four retained fields be given, with every laid-out field of length at most $W$, and let the native input be the packed encoding of the layout. Then within $2|x| + 4 + (k'+4)(3W+7) + (4F+7)(T+1)^2$ steps, where $x$ is that input, $k'$ is the recorder tape count, and $F$ the field count, the host reaches the completed recorder configuration holding all stored rows of times $0,\dots,T$, the horizon source configuration, and the protected header fields and input buffer.
\end{lemma}

\begin{definition}[Comparator words of one row query]
\lean{Complexity.clMatchWords}
\label{Complexity.clMatchWords}
\leanok
For positive and negative counter coordinates and a stored row, together with two retained target count words, the four comparator words are, in order: the row's positive count, the second target word, the first target word, and the row's negative count, exactly the cross-sum interface of the banked comparator.
\end{definition}

\begin{definition}[Loader slots inside the search]
\lean{Complexity.clMatchLoadIndex}
\label{Complexity.clMatchLoadIndex}
\leanok
The row loader's $l + 1$ tapes are embedded into the search's $l + 5$ tapes with the candidate fields first and the record stream as the first administrative slot.
\end{definition}

\begin{definition}[Partial inverse of the search-loader slots]
\lean{Complexity.clMatchLoadSelect}
\label{Complexity.clMatchLoadSelect}
\leanok
The search's selection maps each candidate field and the record stream back to their loader coordinates, and leaves the two target words, the query clock, and the flag tape unselected during loading.
\end{definition}

\begin{lemma}[Search-loader selection inverts its injection]
\lean{Complexity.clMatchLoad_inverse}
\label{Complexity.clMatchLoad_inverse}
\leanok
\uses{Complexity.clMatchLoadIndex, Complexity.clMatchLoadSelect}
\difficulty{2}
For every loader tape coordinate, the search's loading selection applied to the embedded slot returns that coordinate, with no aliasing.
\end{lemma}

\begin{definition}[Comparator slots inside the search]
\lean{Complexity.clMatchCmpIndex}
\label{Complexity.clMatchCmpIndex}
\leanok
For distinct positive and negative counter coordinates, the four comparator tapes are embedded as: the candidate's positive count field, the second target slot, the first target slot, and the candidate's negative count field; the injection is fixed, with no data-dependent or uncharged record lookup.
\end{definition}

\begin{definition}[Partial inverse of the search-comparator slots]
\lean{Complexity.clMatchCmpSelect}
\label{Complexity.clMatchCmpSelect}
\leanok
The search's comparison selection maps the positive and negative candidate fields and the two target slots back to their comparator coordinates, and leaves the stream, clock, flag tape, and all other fields fixed.
\end{definition}

\begin{lemma}[Comparator selection inverts its injection]
\lean{Complexity.clMatchCmp_inverse}
\label{Complexity.clMatchCmp_inverse}
\leanok
\uses{Complexity.clMatchCmpIndex, Complexity.clMatchCmpSelect}
\difficulty{2}
For distinct positive and negative coordinates and each of the four comparator tapes, the search's comparison selection applied to the embedded slot returns that tape, so the call is a genuine four-tape comparison.
\end{lemma}

\begin{definition}[Control phases of the matcher]
\lean{Complexity.clMatchState}
\label{Complexity.clMatchState}
\leanok
\uses{Complexity.clCmpTM, Complexity.clLoadTM}
The matcher's finite control is a disjoint sum of four phases: the clock test, the sequential row load, the native comparison, and the completed live return.
\end{definition}

\begin{definition}[Native all-earlier-row matcher]
\lean{Complexity.clMatchTM}
\label{Complexity.clMatchTM}
\leanok
\uses{Complexity.clCmpTM, Complexity.clLoadTM, Complexity.clMatchCmpIndex, Complexity.clMatchCmpSelect, Complexity.clMatchLoadIndex, Complexity.clMatchLoadSelect, Complexity.clMatchState, Complexity.clSlotAction, Turing.FinTM, Turing.FinTM.controlAction}
For distinct positive and negative counter coordinates, the matcher repeatedly tests its unary query clock, loads one stored row from the record stream, compares cross-sums against two retained target words, and appends one flag per row in source-time order; when the clock is exhausted it stops before inspecting the next row, so exactly the clock's length of rows is inspected and the row at the target time is excluded.
\end{definition}

\begin{definition}[Complete search seam]
\lean{Complexity.clMatchCfg}
\label{Complexity.clMatchCfg}
\leanok
\uses{Complexity.clMatchState, Complexity.clMatchTM, Turing.Cfg, Turing.FinTM.bufferTape}
For a matcher control, candidate fields, two immutable target count words, the record stream with an explicit cursor, a unary query clock of $N$ tokens with its head at $j$, and the stored flag prefix, the search configuration holds every binary field at head zero, the flag tape head at the flag count, and empty physical output.
\end{definition}

\begin{lemma}[A live clock token starts one row query]
\lean{Complexity.clMatch_tick}
\label{Complexity.clMatch_tick}
\leanok
\uses{Complexity.clMatchCfg, Complexity.clMatchTM, Turing.Action.apply, Turing.Cfg.workTapeSymbols, Turing.FinTM.controlAction_apply, Turing.MultiTapeTM.step, Turing.moveInputPos_zero}
\difficulty{3}
For every search seam with query-clock position $j < N$, one matcher step reads the live clock token and dispatches to the row-loading phase; no source-time bookkeeping, stream bit, or stored flag changes.
\end{lemma}

\begin{lemma}[The exhausted query clock stops the search]
\lean{Complexity.clMatch_stop}
\label{Complexity.clMatch_stop}
\leanok
\uses{Complexity.clMatchCfg, Complexity.clMatchTM, Turing.Action.apply, Turing.Cfg.workTapeSymbols, Turing.FinTM.controlAction_apply, Turing.MultiTapeTM.step, Turing.moveInputPos_zero}
\difficulty{3}
For every search seam whose query clock is exhausted, one matcher step reads the clock blank and enters the completed live control without reading any stream bit; this includes the empty time-zero query.
\end{lemma}

\begin{lemma}[Loader frame inside the search]
\lean{Complexity.clMatch_load_frame}
\label{Complexity.clMatch_load_frame}
\leanok
\uses{Complexity.clLoadCfg, Complexity.clLoadTM, Complexity.clMatchCfg, Complexity.clMatchLoadSelect, Complexity.clMatchState, Complexity.clSlotCfg}
\difficulty{3}
For every loader control, old and new candidate rows, target words, stream, cursors, clock data, and flags, relocating the row-loader configuration into the search's slots, retaining the search's own administrative tapes, yields the search seam in the loading phase; both target words, the clock, and all stored flags are framed, while the candidate row and stream cursor may change.
\end{lemma}

\begin{lemma}[One full row is loaded into the search]
\lean{Complexity.clMatch_load}
\label{Complexity.clMatch_load}
\leanok
\uses{Complexity.clFields, Complexity.clLoadCfg, Complexity.clLoadTM, Complexity.clLoad_first, Complexity.clMatchCfg, Complexity.clMatchLoadIndex, Complexity.clMatchLoadSelect, Complexity.clMatchLoad_inverse, Complexity.clMatchState, Complexity.clMatchTM, Complexity.clMatch_load_frame, Complexity.clSlot_run, Turing.Action.apply, Turing.FinTM.controlAction_apply, Turing.MultiTapeTM.runFrom, Turing.MultiTapeTM.runFrom_succ_eq_step', Turing.moveInputPos_zero}
\difficulty{5}
For a positive field count, distinct coordinates, old and new candidate rows of field length at most $W$, target words, and a stream exposing the packed new row after a prefix, there is a time $t \le l(5W+7) + 1$ after which the matcher has loaded the new row, cleared all old field targets, advanced the stream cursor exactly past the packed row, and dispatched to the signed comparison with fresh comparator registers.
\end{lemma}

\begin{lemma}[Comparator frame inside the search]
\lean{Complexity.clMatch_cmp_frame}
\label{Complexity.clMatch_cmp_frame}
\leanok
\uses{Complexity.clCmpCfg, Complexity.clCmpTM, Complexity.clMatchCfg, Complexity.clMatchCmpSelect, Complexity.clMatchState, Complexity.clMatchWords, Complexity.clSlotCfg}
\difficulty{4}
For distinct positive and negative coordinates, every comparator control, candidate row, target words, stream, cursor, clock data, and flags, relocating the four-word comparator configuration into the search's slots yields the search seam in the comparison phase; the remaining row fields, stream cursor, clock, and accumulated flags are restored literally from the frame.
\end{lemma}

\begin{lemma}[Committing one comparison flag]
\lean{Complexity.clMatch_commit}
\label{Complexity.clMatch_commit}
\leanok
\uses{Complexity.clMatchCfg, Complexity.clMatchTM, Turing.Action.apply, Turing.FinTM.bufferTape_append, Turing.MultiTapeTM.step, Turing.moveInputPos_zero}
\difficulty{3}
For every verdict bit $b$ and search seam in the comparison-finished control, one matcher step appends $b$ to the flag tape by one real write, advances the query clock by one token, and returns to the clock-test phase; the physical output stays silent.
\end{lemma}

\begin{lemma}[Signed comparison with stored flag]
\lean{Complexity.clMatch_compare}
\label{Complexity.clMatch_compare}
\leanok
\uses{Complexity.clCmpCfg, Complexity.clCmpTM, Complexity.clCmp_first, Complexity.clMatchCfg, Complexity.clMatchCmpIndex, Complexity.clMatchCmpSelect, Complexity.clMatchCmp_inverse, Complexity.clMatchState, Complexity.clMatchTM, Complexity.clMatchWords, Complexity.clMatch_cmp_frame, Complexity.clMatch_commit, Complexity.clNum, Complexity.clSlot_run, Turing.MultiTapeTM.runFrom, Turing.MultiTapeTM.runFrom_succ_eq_step'}
\difficulty{5}
For distinct positive and negative coordinates, a candidate row and target words of length at most $W$, and any stream, cursor, clock data, and flags, there is a time $t \le 2W + 3$ and a bit $b$, true exactly when the row's positive count plus the second target equals the first target plus the row's negative count, such that the matcher completes the comparison, restores all four word heads, stores $b$ on the flag tape, and advances the clock; cross-sum equality, not the raw words, determines the flag.
\end{lemma}

\begin{definition}[Generic ordered row stream]
\lean{Complexity.clRows}
\label{Complexity.clRows}
\leanok
\uses{Complexity.clFields}
For every family of rows of $l$ binary fields indexed by time and every count $n$, the stream of the first $n$ rows is the concatenation of the packed self-delimiting encodings of rows $0, \dots, n-1$ in order, in exactly the banked field format.
\end{definition}

\begin{lemma}[Splitting the row stream]
\lean{Complexity.clRows_add}
\label{Complexity.clRows_add}
\leanok
\uses{Complexity.clRows}
\difficulty{3}
For every row family and all counts $n$ and $m$, the stream of the first $n + m$ rows is the stream of the first $n$ rows followed by the stream of $m$ rows of the shifted family; the literal suffix is preserved.
\end{lemma}

\begin{lemma}[Factoring out the next row]
\lean{Complexity.clRows_split}
\label{Complexity.clRows_split}
\leanok
\uses{Complexity.clFields, Complexity.clRows, Complexity.clRows_add}
\difficulty{3}
For every row family, bound $N$, time $j < N$, and trailer, there is a remainder word such that the full stream of $N$ rows plus trailer equals the stream of the first $j$ rows, followed by the packed encoding of row $j$, followed by that remainder; obtaining the row itself still requires the native loader.
\end{lemma}

\begin{lemma}[Row streams specialize to the stored trajectory]
\lean{Complexity.clRows_records}
\label{Complexity.clRows_records}
\leanok
\uses{Complexity.clCounts, Complexity.clRecords, Complexity.clRowPrefix_fields, Complexity.clRows, Turing.FinTM}
\difficulty{3}
For every source machine, virtual input $y$, and count $n$, the generic row stream of the binary movement counts at each time equals the recorder's serialized trajectory prefix of $n$ rows.
\end{lemma}

\begin{definition}[Flag of one signed-position comparison]
\lean{Complexity.clMatchFlag}
\label{Complexity.clMatchFlag}
\leanok
\uses{Complexity.clNum}
For positive and negative coordinates, a stored row, and two target count words, the flag is true exactly when the row's positive count plus the second target value equals the first target value plus the row's negative count.
\end{definition}

\begin{lemma}[One complete search iteration]
\lean{Complexity.clMatch_round}
\label{Complexity.clMatch_round}
\leanok
\uses{Complexity.clFields, Complexity.clMatchCfg, Complexity.clMatchFlag, Complexity.clMatchTM, Complexity.clMatch_compare, Complexity.clMatch_load, Complexity.clMatch_tick, Turing.MultiTapeTM.runFrom, Turing.MultiTapeTM.runFrom_add, Turing.MultiTapeTM.runFrom_succ_eq_step}
\difficulty{4}
For distinct coordinates, old and new candidate rows and target words of field length at most $W$, a stream exposing the packed new row, and a clock position $j < N$, there is a time $t \le l(5W+7) + 2W + 5$ after which the matcher has consumed one clock token, loaded row $j$, compared cross-sums, appended exactly the comparison flag of that row, and advanced its cursor past the row; every reset, stream read, comparison, dispatch, and flag write is included in the bound.
\end{lemma}

\begin{definition}[Candidate bank before the next load]
\lean{Complexity.clPriorRow}
\label{Complexity.clPriorRow}
\leanok
For a row family and a time $j$, the candidate bank before loading row $j$ holds empty fields when $j = 0$ and row $j - 1$ otherwise.
\end{definition}

\begin{lemma}[The strict-prefix scan stores every flag]
\lean{Complexity.clMatch_prefix}
\label{Complexity.clMatch_prefix}
\leanok
\uses{Complexity.clMatchCfg, Complexity.clMatchFlag, Complexity.clMatchTM, Complexity.clMatch_round, Complexity.clPriorRow, Complexity.clRows, Complexity.clRows_split, Turing.MultiTapeTM.runFrom, Turing.MultiTapeTM.runFrom_add}
\difficulty{5}
For distinct coordinates, a row family with all fields of length at most $W$ among the first $N$ rows, target words of length at most $W$, a trailer, and every $j \le N$, there is a time $t \le j\,(l(5W+7) + 2W + 5)$ after which the matcher, started with empty candidates, cursor zero, and no flags, holds exactly the comparison flags of rows $0,\dots,j-1$ in order, with the stream cursor at the end of the first $j$ rows and the immutable target words framed throughout.
\end{lemma}

\begin{lemma}[Complete bounded search]
\lean{Complexity.clMatch_complete}
\label{Complexity.clMatch_complete}
\leanok
\uses{Complexity.clMatchCfg, Complexity.clMatchFlag, Complexity.clMatchTM, Complexity.clMatch_prefix, Complexity.clMatch_stop, Complexity.clPriorRow, Complexity.clRows, Turing.MultiTapeTM.runFrom, Turing.MultiTapeTM.runFrom_succ_eq_step'}
\difficulty{2}
For distinct coordinates, a row family with fields of length at most $W$ among the first $N$ rows, target words of length at most $W$, and a trailer, there is a time $t \le N(l(5W+7) + 2W + 5) + 1$ after which the matcher is in its completed control holding exactly the $N$ comparison flags in source-time order; the next row and the arbitrary tail are untouched, and at $N = 0$ only the silent stop executes.
\end{lemma}

\begin{lemma}[Search flags test the public schedule]
\lean{Complexity.clMatchFlag_schedule}
\label{Complexity.clMatchFlag_schedule}
\leanok
\uses{Complexity.clCounts, Complexity.clCounts_schedule, Complexity.clMatchFlag, Complexity.clNum_bits, Complexity.clSigned_eq, Complexity.clTwo, Complexity.workPosAt, Turing.FinTM}
\difficulty{3}
For every oblivious machine, reference length $m$, times $s$ and $t$, and work tape $\tau$, the comparison flag of the stored row at time $s$ against the target counts at time $t$, taken at the positive and negative slots of $\tau$, is true exactly when the public work-head position on $\tau$ at time $s$ equals the position at time $t$.
\end{lemma}

\begin{definition}[Optional-time code]
\lean{Complexity.clVisitCode}
\label{Complexity.clVisitCode}
\leanok
\uses{Turing.pairEncode}
An absent visit is encoded as the empty word; a present visit at time $s$ is encoded as an aligned separator followed by the canonical little-endian binary representation of $s$. The code is unambiguous and compact.
\end{definition}

\begin{definition}[Code of the last marked time]
\lean{Complexity.clLastCode}
\label{Complexity.clLastCode}
\leanok
\uses{Complexity.clVisitCode, Turing.splitAtLastTrue}
For an ordered flag word, the last-visit code encodes the index of the last true flag, obtained by a reverse marker scan followed by a binary length computation, or the empty word when no flag is set.
\end{definition}

\begin{lemma}[The last-visit code is native]
\lean{Complexity.clLastCode_native}
\label{Complexity.clLastCode_native}
\leanok
\uses{Complexity.PolyTimeComputable, Complexity.PolyTimeComputable.comp, Complexity.clLastCode, Complexity.clNative_linear, Complexity.clNative_map, Complexity.clNative_pair, Complexity.clVisitCode, Complexity.polyTimeComputable_id, Turing.FinTM.computesFunInTime_const, Turing.FinTM.computesFunInTime_lengthBits, Turing.FinTM.computesFunInTime_stripLast, Turing.pairDecode, Turing.pairDecode_pairEncode, Turing.pairEncode, Turing.splitAtLastTrue}
\difficulty{4}
The function taking an ordered flag word to the code of its last marked index is polynomial-time computable: the flag word is threaded through the last-true stripper, an absent result maps to the empty code, and otherwise the stripped prefix's binary length is exactly the index of the last marked row.
\end{lemma}

\begin{lemma}[Streaming recurrence of the marker scan]
\lean{Complexity.clLastMarker_step}
\label{Complexity.clLastMarker_step}
\leanok
\uses{Turing.splitAtLastTrue}
\difficulty{2}
For every flag word $w$ and flag $b$, splitting at the last true flag of $w$ extended by $b$ yields the whole prefix $w$ when $b$ is true, and otherwise the split of $w$ itself.
\end{lemma}

\begin{definition}[Pure last-match tracker]
\lean{Complexity.clLastIndex}
\label{Complexity.clLastIndex}
\leanok
For a predicate on times and a bound, the tracker scans all indices below the bound in order, replacing its candidate at every match, and returns the last matching index, or nothing when no index matches.
\end{definition}

\begin{lemma}[The marker scan computes the tracker]
\lean{Complexity.clLastIndex_marker}
\label{Complexity.clLastIndex_marker}
\leanok
\uses{Complexity.clLastIndex, Complexity.clLastMarker_step, Turing.splitAtLastTrue}
\difficulty{3}
For every predicate $P$ and bound $n$, mapping the length over the last-true split of the flag word of $P$ on indices $0,\dots,n-1$ yields exactly the tracker's last matching index below $n$.
\end{lemma}

\begin{lemma}[The last match is the filtered maximum]
\lean{Complexity.clLastIndex_max}
\label{Complexity.clLastIndex_max}
\leanok
\uses{Complexity.clLastIndex}
\difficulty{4}
For every predicate $P$ and bound $n$, the tracker's candidate below $n$ equals the maximum of the indices below $n$ satisfying $P$; in particular the empty candidate set at time zero is certified.
\end{lemma}

\begin{lemma}[Flag codes compute the last-visit schedule]
\lean{Complexity.clLastCode_prev}
\label{Complexity.clLastCode_prev}
\leanok
\uses{Complexity.clLastCode, Complexity.clLastIndex_marker, Complexity.clLastIndex_max, Complexity.clVisitCode, Complexity.prevVisit, Complexity.workPosAt, Turing.FinTM}
\difficulty{2}
For every oblivious machine, reference length $m$, time $t$, and work tape $\tau$, the last-visit code of the flag word testing, for each $s < t$, whether the work position at $s$ equals the position at $t$, is exactly the optional-time code of the public last-visit value; the time is in canonical binary.
\end{lemma}

\begin{definition}[Buffer replay machine]
\lean{Complexity.clReplayTM}
\label{Complexity.clReplayTM}
\leanok
\uses{Turing.FinTM}
A three-state machine with one work tape replays a stored buffer to the physical output: started at the buffer's right blank, it rewinds explicitly to the origin, emits every stored bit left to right, and physically halts on the final blank; no verdict or terminator is added.
\end{definition}

\begin{definition}[Replay configuration]
\lean{Complexity.clReplayCfg}
\label{Complexity.clReplayCfg}
\leanok
\uses{Turing.Cfg, Turing.FinTM.bufferTape}
For an input word, an input position, an optional control, a stored buffer, a head cell, and current output, the replay configuration is the one-tape configuration retaining the buffer and the native input head.
\end{definition}

\begin{lemma}[Replay's backward rewind]
\lean{Complexity.clReplay_back}
\label{Complexity.clReplay_back}
\leanok
\uses{Complexity.clReplayCfg, Complexity.clReplayTM, Turing.Action.apply, Turing.Cfg.workTapeSymbols, Turing.FinTM.bufferTape, Turing.FinTM.bufferTape_nat, Turing.MultiTapeTM.runFrom, Turing.MultiTapeTM.runFrom_succ_eq_step, Turing.MultiTapeTM.runFrom_zero, Turing.MultiTapeTM.step}
\difficulty{4}
For every input word, input position, buffer $w$, output, and count $j \le |w|$, running the replay machine in its rewind control at cell $j - 1$ for $j + 1$ steps restores the buffer head to zero and enters the emitting control; each occupied prefix cell is read once and the empty-buffer case is included.
\end{lemma}

\begin{lemma}[Forward replay emits the suffix]
\lean{Complexity.clReplay_forward}
\label{Complexity.clReplay_forward}
\leanok
\uses{Complexity.clReplayCfg, Complexity.clReplayTM, Turing.Action.apply, Turing.Cfg.workTapeSymbols, Turing.FinTM.bufferTape, Turing.FinTM.bufferTape_nat, Turing.MultiTapeTM.runFrom, Turing.MultiTapeTM.runFrom_succ_eq_step, Turing.MultiTapeTM.runFrom_zero, Turing.MultiTapeTM.step}
\difficulty{4}
For every input word, input position, buffer $w$, position $j \le |w|$, and output prefix, running the replay machine in its emitting control for $|w| - j + 1$ steps appends exactly the suffix of $w$ from position $j$ to the output and halts on the right blank, without any spurious trailing bit.
\end{lemma}

\begin{lemma}[Complete replay ledger]
\lean{Complexity.clReplay_run}
\label{Complexity.clReplay_run}
\leanok
\uses{Complexity.clReplayCfg, Complexity.clReplayTM, Complexity.clReplay_back, Complexity.clReplay_forward, Turing.Action.apply, Turing.MultiTapeTM.runFrom, Turing.MultiTapeTM.runFrom_add, Turing.MultiTapeTM.runFrom_succ_eq_step, Turing.MultiTapeTM.step}
\difficulty{3}
For every input word, input position, buffer $w$, and output prefix, running the replay machine from its start control at the buffer's right blank for $2|w| + 3$ steps appends exactly $w$ to the output and physically halts; the initial left move, full rewind, every emitted bit, and the halt are all charged.
\end{lemma}

\begin{lemma}[Quantitative buffered composition]
\lean{Complexity.clCompute_comp}
\label{Complexity.clCompute_comp}
\leanok
\uses{Turing.Cfg.init, Turing.FinTM, Turing.FinTM.ComputesInTime, Turing.FinTM.ComputesInTime.mono, Turing.FinTM.VirtualTag, Turing.FinTM.bufferedCompTM, Turing.FinTM.bufferedComp_start, Turing.FinTM.bufferedSecondCfg, Turing.FinTM.bufferedSecondCfg_run, Turing.FinTM.computesInTime_iff, Turing.MultiTapeTM.initCfg, Turing.MultiTapeTM.output_length_le, Turing.MultiTapeTM.runFrom_add}
\difficulty{4}
Let a machine compute output $y$ from input $x$ within time $a$, and a second machine compute $z$ from $y$ within time $b$. Then the buffered composition of the two machines computes $z$ from $x$ within time $2a + b + 2$; the second machine's contract is needed only on the actual intermediate word.
\end{lemma}

\begin{definition}[Selecting one physical tape]
\lean{Complexity.clOneSelect}
\label{Complexity.clOneSelect}
\leanok
For a chosen tape coordinate, the selection maps that tape to the single local slot and leaves every other tape and head framed.
\end{definition}

\begin{definition}[Attaching a physical replay to a completion]
\lean{Complexity.clOutputTM}
\label{Complexity.clOutputTM}
\leanok
\uses{Complexity.clOneSelect, Complexity.clReplayTM, Complexity.clSlotAction, Turing.Action.mapState, Turing.FinTM, Turing.FinTM.controlAction}
Given a machine with an absorbing live completion control and a chosen buffer tape, the extended machine runs the original computation unchanged, dispatches once at the first completion, and then replays the chosen buffer to the physical output; the source computation and its startup are untouched before the return.
\end{definition}

\begin{lemma}[A completed buffer is physically returned]
\lean{Complexity.clOutput_compute}
\label{Complexity.clOutput_compute}
\leanok
\uses{Complexity.clFirst, Complexity.clMap_run, Complexity.clOneSelect, Complexity.clOutputTM, Complexity.clReplayCfg, Complexity.clReplayTM, Complexity.clReplay_run, Complexity.clSlotCfg, Complexity.clSlot_run, Turing.Cfg, Turing.Cfg.mapState, Turing.FinTM, Turing.FinTM.ComputesInTime, Turing.FinTM.ComputesInTime.mono, Turing.FinTM.bufferTape, Turing.FinTM.computesInTime_iff, Turing.FinTM.controlAction_apply, Turing.MultiTapeTM.initCfg, Turing.MultiTapeTM.runFrom, Turing.MultiTapeTM.runFrom_add, Turing.MultiTapeTM.runFrom_succ_eq_step', Turing.MultiTapeTM.step, Turing.moveInputPos_zero}
\difficulty{5}
Let a machine reach, in $a$ steps from its native start, a configuration in an absorbing completion control with empty output, whose chosen tape holds exactly the word $w$ as a buffer with its head at the right blank, and whose start control differs from the completion. Then the replay-extended machine computes $w$ from the same input within $a + 2|w| + 4$ steps, with an explicit replay ledger and every other tape framed.
\end{lemma}

\begin{definition}[Native loading before a prepared machine]
\lean{Complexity.clPreparedTM}
\label{Complexity.clPreparedTM}
\leanok
\uses{Complexity.clPrepareTM, Complexity.clSlotAction, Turing.Action.mapState, Turing.FinTM, Turing.FinTM.controlAction}
For a machine expecting prepared work words, the extended machine first parses arbitrary exact work words from its self-delimiting native input into the field bank, keeping the retained input on its own extra tape, then dispatches once and runs the supplied machine on the loaded bank.
\end{definition}

\begin{definition}[Frame of a prepared computation]
\lean{Complexity.clPreparedCfg}
\label{Complexity.clPreparedCfg}
\leanok
\uses{Complexity.clPreparedTM, Complexity.clSlotCfg, Turing.Cfg, Turing.FinTM, Turing.FinTM.bufferTape}
For a prepared machine, an input word, a stream cursor, and an inner configuration, the frame is the inner configuration embedded on the field bank, while the retained input buffer keeps the original word at the given cursor.
\end{definition}

\begin{lemma}[Prepared computation from genuine input]
\lean{Complexity.clPrepared_run}
\label{Complexity.clPrepared_run}
\leanok
\uses{Complexity.clFields, Complexity.clLoadCfg, Complexity.clMap_run, Complexity.clPrepareTM, Complexity.clPrepare_first, Complexity.clPreparedCfg, Complexity.clPreparedTM, Complexity.clSlotCfg, Complexity.clSlot_run, Turing.Action.apply, Turing.Cfg, Turing.Cfg.mapState, Turing.Cfg.ofWords, Turing.FinTM, Turing.FinTM.bufferTape, Turing.FinTM.controlAction_apply, Turing.MultiTapeTM.initCfg, Turing.MultiTapeTM.runFrom, Turing.MultiTapeTM.runFrom_add, Turing.MultiTapeTM.runFrom_succ_eq_step', Turing.MultiTapeTM.step, Turing.moveInputPos_zero}
\difficulty{5}
Let work words of length at most $W$ and a trailer be given, and let the native input be their packed encoding followed by the trailer. Then there is a startup time $a \le 2|x| + 4 + k(3W+7)$, where $x$ is that input and $k$ the field count, such that for every duration $b$, running the loading-extended machine for $a + b$ steps equals the frame around the supplied machine's $b$-step run from the configuration assembled over the parsed words; no cleanup or random access is assumed.
\end{lemma}

\begin{lemma}[Absorbing returns survive native loading]
\lean{Complexity.clPrepared_idle}
\label{Complexity.clPrepared_idle}
\leanok
\uses{Complexity.clPreparedTM, Complexity.clSlotAction, Turing.Cfg, Turing.FinTM, Turing.FinTM.controlAction, Turing.FinTM.controlAction_apply, Turing.MultiTapeTM.step, Turing.moveInputPos_zero}
\difficulty{3}
If a control of the prepared machine is absorbing, its transition is the stationary control action, then the corresponding control of the loading-extended machine is absorbing as well, with the complete retained-input frame fixed.
\end{lemma}

\begin{definition}[Argument words of one native query]
\lean{Complexity.clQueryWords}
\label{Complexity.clQueryWords}
\leanok
For a field count $l$, two target count words, a record stream, and a strict-prefix bound $N$, the query argument words are: blank candidate fields, the retained stream, the two target counts, a unary clock of $N$ tokens, and a blank flag tape.
\end{definition}

\begin{lemma}[Prepared seam of a native query]
\lean{Complexity.clQuery_initial}
\label{Complexity.clQuery_initial}
\leanok
\uses{Complexity.clMatchCfg, Complexity.clMatchTM, Complexity.clQueryWords, Turing.Cfg.ofWords}
\difficulty{3}
For all coordinates, input word, target words, stream, and bound $N$, the configuration assembled from the query argument words at the matcher's start control equals the search seam with empty candidates, cursor zero, full clock, and no flags, including every scratch tape and head.
\end{lemma}

\begin{lemma}[The matcher's return is stationary]
\lean{Complexity.clMatch_return}
\label{Complexity.clMatch_return}
\leanok
\uses{Complexity.clMatchTM, Turing.FinTM.controlAction}
\difficulty{1}
On every input symbol and work reads, the matcher's transition at its completed control is the stationary control action returning to the same control.
\end{lemma}

\begin{definition}[Complete native query transducer]
\lean{Complexity.clQueryFlagsTM}
\label{Complexity.clQueryFlagsTM}
\leanok
\uses{Complexity.clMatchTM, Complexity.clOutputTM, Complexity.clPreparedTM, Turing.FinTM}
For distinct positive and negative coordinates, the query transducer prepares its field bank from the actual encoded input, runs the all-earlier-row matcher over exactly the clocked strict prefix, and then physically replays the stored flag word as its output.
\end{definition}

\begin{lemma}[Output and ledger of the native query]
\lean{Complexity.clQueryFlags_compute}
\label{Complexity.clQueryFlags_compute}
\leanok
\uses{Complexity.clFields, Complexity.clMatchCfg, Complexity.clMatchFlag, Complexity.clMatchTM, Complexity.clMatch_complete, Complexity.clMatch_return, Complexity.clOutput_compute, Complexity.clPreparedCfg, Complexity.clPreparedTM, Complexity.clPrepared_idle, Complexity.clPrepared_run, Complexity.clPriorRow, Complexity.clQueryFlagsTM, Complexity.clQueryWords, Complexity.clQuery_initial, Complexity.clRows, Complexity.clSlotCfg, Turing.FinTM.ComputesInTime, Turing.FinTM.ComputesInTime.mono, Turing.MultiTapeTM.initCfg, Turing.MultiTapeTM.runFrom}
\difficulty{5}
For distinct coordinates, a row family with fields of length at most $W$ among the first $N$ rows, target words of length at most $W$, a trailer, and argument words of length at most $U$, the query transducer computes, from the packed encoding of the query argument, exactly the list of the $N$ comparison flags in source-time order, within $2|a| + 9 + (l+5)(3U+7) + N(l(5W+7) + 2W + 7)$ steps, where $a$ is the argument; argument copying, preparation, every row query, comparison, dispatch, flag write, rewind, replay, and the halt are all charged.
\end{lemma}

\begin{lemma}[One machine returns the last marked time]
\lean{Complexity.clQueryCode_machine}
\label{Complexity.clQueryCode_machine}
\leanok
\uses{Complexity.clCompute_comp, Complexity.clFields, Complexity.clLastCode, Complexity.clLastCode_native, Complexity.clMatchFlag, Complexity.clQueryFlagsTM, Complexity.clQueryFlags_compute, Complexity.clQueryWords, Complexity.clRows, Turing.FinTM, Turing.FinTM.ComputesInTime, Turing.FinTM.bufferedCompTM}
\difficulty{3}
For distinct positive and negative coordinates, there exist a single machine $Q$ and constants $c, e$ such that on every canonical query argument, a packed row family, target words, and strict-prefix clock, $Q$ computes the optional-time code of the last row whose cross-sums match the targets, within twice the query ledger plus $c(N+1)^e + 2$ steps. Only canonical inputs are promised; callers must establish them.
\end{lemma}

\begin{lemma}[Fields are dominated by their packing]
\lean{Complexity.clFields_width}
\label{Complexity.clFields_width}
\leanok
\uses{Complexity.clFields, Turing.pairEncode}
\difficulty{3}
For every list of words and every member $w$ of the list, the length of $w$ is at most the length of the packed encoding of the whole list; the doubled bits and separators make the bound immediate without random access.
\end{lemma}

\begin{lemma}[The completed recorder host is absorbing]
\lean{Complexity.clRecord_idle}
\label{Complexity.clRecord_idle}
\leanok
\uses{Complexity.clRecTM, Complexity.clRecordSelect, Complexity.clRecordTM, Complexity.clSlotAction, Turing.Cfg, Turing.FinTM, Turing.FinTM.controlAction, Turing.FinTM.controlAction_apply, Turing.MultiTapeTM.step, Turing.moveInputPos_zero}
\difficulty{3}
Every configuration of the recorder host whose control is the recorder's completed return is left entirely unchanged by one step, with the full retained input and header frame unchanged.
\end{lemma}

\begin{definition}[Physically returning the trajectory]
\lean{Complexity.clRecordOutputTM}
\label{Complexity.clRecordOutputTM}
\leanok
\uses{Complexity.clOutputTM, Complexity.clRecFields, Complexity.clRecordIndex, Complexity.clRecordTM, Turing.FinTM}
The trajectory-output machine runs the prepared recorder host to completion and then replays its actual record tape as the entire physical output, preserving the recorder's effects-first silence until the replay.
\end{definition}

\begin{lemma}[Arithmetic of the replay ledger]
\lean{Complexity.clRecord_replay_budget}
\label{Complexity.clRecord_replay_budget}
\leanok
\difficulty{2}
For all naturals, if a recording time is at most $2n + 4 + k(3W+7) + (4l+7)(T+1)^2$ and a replayed length is at most $2l(T+1)^2$, then the recording time plus twice the replayed length plus four is at most $2n + 8 + k(3W+7) + (8l+7)(T+1)^2$.
\end{lemma}

\begin{lemma}[Quadratic envelope of recording]
\lean{Complexity.clRecord_size_budget}
\label{Complexity.clRecord_size_budget}
\leanok
\difficulty{3}
For all naturals with $T \le n$, the recording-and-replay ledger $2n + 8 + k(3n+7) + (8l+7)(T+1)^2$ is at most $(10 + 10k + 8l + 7)(n+1)^2$, a quadratic envelope covering preparation, inclusive recording, and replay.
\end{lemma}

\begin{lemma}[Envelope for one final count]
\lean{Complexity.clCount_size_budget}
\label{Complexity.clCount_size_budget}
\leanok
\difficulty{3}
For all naturals with $R \le T \le n$ and a recording time within the standard recording ledger, the recording time plus $R + 4$ is at most $(10 + 10k + 8l + 7)(n+1)^2$: the same argument envelope pays for returning one final count instead of the full trajectory, including its physical output and halt.
\end{lemma}

\begin{lemma}[Recording with physical trajectory output]
\lean{Complexity.clRecordOutput_compute}
\label{Complexity.clRecordOutput_compute}
\leanok
\uses{Complexity.clFields, Complexity.clOutput_compute, Complexity.clRecCfg, Complexity.clRecFields, Complexity.clRecTM, Complexity.clRecWords, Complexity.clRecordCfg, Complexity.clRecordIndex, Complexity.clRecordOutputTM, Complexity.clRecordSelect, Complexity.clRecordTM, Complexity.clRecord_complete, Complexity.clRecord_idle, Complexity.clRecord_replay_budget, Complexity.clRecords, Complexity.clRecords_length, Complexity.clSlotCfg, Turing.FinTM, Turing.FinTM.ComputesInTime, Turing.FinTM.ComputesInTime.mono, Turing.FinTM.tapeBlocks}
\difficulty{4}
For every source machine, virtual input $y$, horizon $T$, retained fields, and width bound $W$ covering all laid-out words, the trajectory-output machine computes, from the packed encoding of the layout, the complete serialized inclusive trajectory of rows $0,\dots,T$, within $2|x| + 8 + (k'+4)(3W+7) + (8F+7)(T+1)^2$ steps, where $x$ is the input, $k'$ the recorder tape count, and $F$ the field count; preparation, recording, rewind, replay, and the terminal transition are all charged.
\end{lemma}

\begin{lemma}[Quadratic trajectory bound]
\lean{Complexity.clRecordOutput_quadratic}
\label{Complexity.clRecordOutput_quadratic}
\leanok
\uses{Complexity.clFields, Complexity.clFields_width, Complexity.clRecClockIndex, Complexity.clRecFields, Complexity.clRecTM, Complexity.clRecWords, Complexity.clRecordOutputTM, Complexity.clRecordOutput_compute, Complexity.clRecord_size_budget, Complexity.clRecords, Turing.FinTM, Turing.FinTM.ComputesInTime, Turing.FinTM.ComputesInTime.mono, Turing.FinTM.tapeBlocks}
\difficulty{3}
For every source machine, virtual input $y$, horizon $T$, and retained fields, the trajectory-output machine computes the inclusive trajectory from the packed layout within a constant times the square of the encoded argument length plus one; every initial word and the unary horizon occur inside the argument, so the bound covers a zero horizon as well.
\end{lemma}

\begin{lemma}[Per-image polynomial composition]
\lean{Complexity.clNative_image}
\label{Complexity.clNative_image}
\leanok
\uses{Complexity.PolyTimeComputable, Complexity.clCompute_comp, Turing.FinTM, Turing.FinTM.ComputesInTime, Turing.FinTM.ComputesInTime.mono, Turing.FinTM.bufferedCompTM, Turing.FinTM.computesInTime_iff, Turing.MultiTapeTM.output_length_le}
\difficulty{5}
Let $f$ be polynomial-time computable, and let a machine compute $g(x)$ from $f(x)$ within $b\,(|f(x)|+1)^r$ steps for every input $x$, correctness being required only on the image of $f$. Then $g$ is polynomial-time computable; the produced intermediate word is bounded by the first machine's runtime and the full buffered-composition ledger is charged.
\end{lemma}

\begin{lemma}[The trajectory is a native output of the instance]
\lean{Complexity.clRecords_native}
\label{Complexity.clRecords_native}
\leanok
\uses{Complexity.PolyTimeComputable, Complexity.clHeaderLayout_exact, Complexity.clNative_image, Complexity.clRecFields, Complexity.clRecTM, Complexity.clRecordArgument_native, Complexity.clRecordOutputTM, Complexity.clRecordOutput_quadratic, Complexity.clRecords, Turing.FinTM}
\difficulty{3}
For every source machine and parameters $C, e, c, A, d$, the function sending an instance $x$ to the inclusive serialized trajectory over the all-false reference input of length $m = |x| + C(|x|+1)^e$ at horizon $T = c(A(m+1)^d+1)^2$, with $T+1$ rows, is polynomial-time computable; the output is produced by the actual recorder, not read off a prepared configuration.
\end{lemma}

\begin{lemma}[Replay from a prefix boundary]
\lean{Complexity.clReplay_from}
\label{Complexity.clReplay_from}
\leanok
\uses{Complexity.clReplayCfg, Complexity.clReplayTM, Complexity.clReplay_back, Complexity.clReplay_forward, Turing.Action.apply, Turing.MultiTapeTM.runFrom, Turing.MultiTapeTM.runFrom_add, Turing.MultiTapeTM.runFrom_succ_eq_step, Turing.MultiTapeTM.step}
\difficulty{3}
For every input word, input position, buffer $w$, output prefix, and head position $j \le |w|$, running the replay machine from its start control at cell $j$ for $j + |w| + 3$ steps appends exactly $w$ to the output and physically halts; only the prefix rewind, the full emission, and the halt are charged.
\end{lemma}

\begin{lemma}[Physical return from a left-edge buffer]
\lean{Complexity.clOutputAt_compute}
\label{Complexity.clOutputAt_compute}
\leanok
\uses{Complexity.clFirst, Complexity.clMap_run, Complexity.clOneSelect, Complexity.clOutputTM, Complexity.clReplayCfg, Complexity.clReplayTM, Complexity.clReplay_from, Complexity.clSlotCfg, Complexity.clSlot_run, Turing.Cfg, Turing.Cfg.mapState, Turing.FinTM, Turing.FinTM.ComputesInTime, Turing.FinTM.ComputesInTime.mono, Turing.FinTM.bufferTape, Turing.FinTM.computesInTime_iff, Turing.FinTM.controlAction_apply, Turing.MultiTapeTM.initCfg, Turing.MultiTapeTM.runFrom, Turing.MultiTapeTM.runFrom_add, Turing.MultiTapeTM.runFrom_succ_eq_step', Turing.MultiTapeTM.step, Turing.moveInputPos_zero}
\difficulty{5}
Let a machine reach, in $a$ steps from its native start, an absorbing completion whose chosen tape holds the word $w$ with its head at some position $j \le |w|$, with empty output and a start control distinct from the completion. Then the replay-extended machine computes $w$ from the same input within $a + j + |w| + 4$ steps; the whole source frame is preserved and every rewind and output transition appears in the ledger.
\end{lemma}

\begin{definition}[Physically returning one final count]
\lean{Complexity.clCountOutputTM}
\label{Complexity.clCountOutputTM}
\leanok
\uses{Complexity.clOutputTM, Complexity.clRecFields, Complexity.clRecordIndex, Complexity.clRecordTM, Turing.FinTM}
For a selected movement-counter coordinate, the count-output machine runs the prepared recorder host to completion and then replays that counter's final canonical binary word as the physical output, starting from its restored head at zero.
\end{definition}

\begin{lemma}[Quadratic bound for one final count]
\lean{Complexity.clCountOutput_quadratic}
\label{Complexity.clCountOutput_quadratic}
\leanok
\uses{Complexity.clCountOutputTM, Complexity.clCount_size_budget, Complexity.clCounts, Complexity.clCounts_bound, Complexity.clElapsed_width, Complexity.clFields, Complexity.clFields_width, Complexity.clOutputAt_compute, Complexity.clRecCfg, Complexity.clRecClockIndex, Complexity.clRecFields, Complexity.clRecTM, Complexity.clRecWords, Complexity.clRecordCfg, Complexity.clRecordIndex, Complexity.clRecordSelect, Complexity.clRecordTM, Complexity.clRecord_complete, Complexity.clRecord_idle, Complexity.clSlotCfg, Turing.FinTM, Turing.FinTM.ComputesInTime, Turing.FinTM.ComputesInTime.mono, Turing.FinTM.tapeBlocks}
\difficulty{4}
For every source machine, counter coordinate, virtual input $y$, horizon $T$, and retained fields, the count-output machine computes the binary representation of the selected movement count at time $T$ from the packed layout, and its time bound is the quadratic envelope in the encoded argument length that also covers the full trajectory output.
\end{lemma}

\begin{lemma}[Native assembly of recorder arguments]
\lean{Complexity.clRecordWords_native}
\label{Complexity.clRecordWords_native}
\leanok
\uses{Complexity.PolyTimeComputable, Complexity.clFields, Complexity.clNative_fields, Complexity.clNative_linear, Complexity.clRecTM, Complexity.clRecWords, Turing.FinTM, Turing.FinTM.computesFunInTime_const, Turing.FinTM.tapeBlocks}
\difficulty{4}
If a reference-input producer and a unary-horizon producer are polynomial-time computable, then so is the function packing the recorder's full argument: blank counters, empty record, the produced unary clock, the produced virtual input, blank source tapes, and four empty retained fields, as self-delimiting fields.
\end{lemma}

\begin{lemma}[General native trajectory production]
\lean{Complexity.clTrajectory_native}
\label{Complexity.clTrajectory_native}
\leanok
\uses{Complexity.PolyTimeComputable, Complexity.clNative_image, Complexity.clRecFields, Complexity.clRecTM, Complexity.clRecordOutputTM, Complexity.clRecordOutput_quadratic, Complexity.clRecordWords_native, Complexity.clRecords, Turing.FinTM}
\difficulty{2}
If a reference-input producer and a unary-horizon producer are polynomial-time computable, then the function sending an input $x$ to the inclusive serialized trajectory of the source over the produced reference input, with the produced horizon plus one rows, is polynomial-time computable; no recomputed semantics or unbounded numeric conversion is used.
\end{lemma}

\begin{lemma}[Final counts are native functions]
\lean{Complexity.clFinalCount_native}
\label{Complexity.clFinalCount_native}
\leanok
\uses{Complexity.PolyTimeComputable, Complexity.clCountOutputTM, Complexity.clCountOutput_quadratic, Complexity.clCounts, Complexity.clNative_image, Complexity.clRecFields, Complexity.clRecTM, Complexity.clRecordWords_native, Turing.FinTM}
\difficulty{2}
If a reference-input producer and a unary-horizon producer are polynomial-time computable, then for every counter coordinate the function sending an input to the binary representation of that movement count at the produced horizon is polynomial-time computable; keeping the time unary avoids any exponential clock.
\end{lemma}

\begin{lemma}[Uniform polynomial query envelope]
\lean{Complexity.clQuery_budget}
\label{Complexity.clQuery_budget}
\leanok
\difficulty{4}
For all naturals with $N \le L$, the complete query ledger, twice the scan cost $2L + 9 + (l+5)(3L+7) + N(l(5L+7) + 2L + 7)$ plus $c(N+1)^e + 2$, is at most $(100(l+10) + c)(L+1)^{e+2}$, one common monomial in the encoded argument length $L$.
\end{lemma}

\begin{definition}[Signed target fields of a search]
\lean{Complexity.clSearchTarget}
\label{Complexity.clSearchTarget}
\leanok
\uses{Complexity.clCounts, Complexity.clRecFields, Complexity.clTwo, Turing.FinTM}
For a source machine, virtual input $y$, time $t$, and positive and negative counter coordinates, the two target fields are the binary representations of the two selected movement counts at time $t$, taken from the final row of an actual source run; they are compared by cross-sums, never by raw word equality.
\end{definition}

\begin{definition}[Complete native query argument]
\lean{Complexity.clSearchArg}
\label{Complexity.clSearchArg}
\leanok
\uses{Complexity.clFields, Complexity.clQueryWords, Complexity.clRecFields, Complexity.clRecords, Complexity.clSearchTarget, Turing.FinTM}
For a source machine, virtual input $y$, time $t$, and two counter coordinates, the query argument packs the query words over the inclusive trajectory of $t+1$ rows with a strict candidate clock of $t$ tokens, so the final row is retained as an unconsumed suffix beyond the scanned prefix.
\end{definition}

\begin{lemma}[Query arguments are native]
\lean{Complexity.clSearchArg_native}
\label{Complexity.clSearchArg_native}
\leanok
\uses{Complexity.PolyTimeComputable, Complexity.clFinalCount_native, Complexity.clNative_fields, Complexity.clNative_linear, Complexity.clQueryWords, Complexity.clRecFields, Complexity.clRecords, Complexity.clSearchArg, Complexity.clSearchTarget, Complexity.clTrajectory_native, Complexity.clTwo, Turing.FinTM, Turing.FinTM.computesFunInTime_const}
\difficulty{4}
If a reference-input producer and a unary-time producer are polynomial-time computable, then for any two counter coordinates the function sending an input to its complete packed query argument is polynomial-time computable; each field comes from the banked recorder, a native final-count output, the unary clock, or a native constant, and the field packer charges every retained copy.
\end{lemma}

\begin{lemma}[The native last-visit search is polynomial]
\lean{Complexity.clSearch_native}
\label{Complexity.clSearch_native}
\leanok
\uses{Complexity.PolyTimeComputable, Complexity.clCounts, Complexity.clCounts_bound, Complexity.clElapsed_width, Complexity.clFields, Complexity.clFields_width, Complexity.clLastCode, Complexity.clMatchFlag, Complexity.clNative_image, Complexity.clQueryCode_machine, Complexity.clQueryWords, Complexity.clQuery_budget, Complexity.clRecFields, Complexity.clRecords, Complexity.clRows, Complexity.clRows_records, Complexity.clSearchArg, Complexity.clSearchArg_native, Complexity.clSearchTarget, Complexity.clTwo, Turing.FinTM, Turing.FinTM.ComputesInTime.mono}
\difficulty{5}
For a source machine, distinct positive and negative counter coordinates, and polynomial-time reference-input and unary-time producers, the function sending an input to the last-visit code of the flags comparing each strictly earlier row's cross-sums against the final targets is polynomial-time computable; the scanner consumes exactly the first $T$ rows of the inclusive $T+1$-row trajectory, and the uniform query envelope charges argument construction, comparisons, reverse search, and output.
\end{lemma}

\begin{lemma}[The public last visit is natively computable]
\lean{Complexity.clPrev_native}
\label{Complexity.clPrev_native}
\leanok
\uses{Complexity.PolyTimeComputable, Complexity.clLastCode_prev, Complexity.clMatchFlag_schedule, Complexity.clRecFields, Complexity.clSearch_native, Complexity.clVisitCode, Complexity.prevVisit, Turing.FinTM}
\difficulty{3}
For a source machine, work tape $\tau$, and polynomial-time producers of a reference length (as an all-false word) and a unary time, the function sending an input to the optional-time code of the public greatest strictly earlier visit to the currently scanned cell of $\tau$ is polynomial-time computable, with the canonical binary time encoding.
\end{lemma}

\begin{definition}[Row of last-visit codes]
\lean{Complexity.clVisitRow}
\label{Complexity.clVisitRow}
\leanok
\uses{Complexity.clFields, Complexity.clHeaderField, Complexity.clHeaderTail, Complexity.clVisitCode, Complexity.prevVisit, Turing.FinTM}
For a source machine, certificate parameters $C$ and $e$, and a request pairing an instance $x$ with a unary query-time word, the visit row is the packed list, in tape order, of the self-delimiting last-visit codes of every work tape at the query time, over the all-false reference input of length $|x| + C(|x|+1)^e$.
\end{definition}

\begin{lemma}[Visit rows are native]
\lean{Complexity.clVisitRow_native}
\label{Complexity.clVisitRow_native}
\leanok
\uses{Complexity.PolyTimeComputable, Complexity.PolyTimeComputable.comp, Complexity.clHeaderField, Complexity.clHeaderField_native, Complexity.clHeaderTail, Complexity.clHeaderTail_native, Complexity.clNative_append, Complexity.clNative_fields, Complexity.clNative_fill, Complexity.clNative_unary, Complexity.clPrev_native, Complexity.clVisitCode, Complexity.clVisitRow, Complexity.prevVisit, Turing.FinTM}
\difficulty{4}
For every source machine and certificate parameters $C$ and $e$, the function sending a retained-instance and unary-time request to the packed row of last-visit codes for all work tapes is polynomial-time computable: native projections recover the instance and time, exact preparation builds the reference input, each tape's search is native, and finite pairing charges the whole ordered row.
\end{lemma}

\begin{lemma}[Visit rows on canonical requests]
\lean{Complexity.clVisitRow_exact}
\label{Complexity.clVisitRow_exact}
\leanok
\uses{Complexity.clFields, Complexity.clHeaderField, Complexity.clHeaderTail, Complexity.clVisitCode, Complexity.clVisitRow, Complexity.prevVisit, Turing.FinTM, Turing.pairDecode_pairEncode, Turing.pairEncode}
\difficulty{1}
For every source machine, parameters $C$ and $e$, instance $x$, and time $t$, the visit row of the encoded pair of $x$ with $t$ unary tokens is exactly the packed list, in source tape order, of the optional-time codes of the greatest strictly earlier visits at time $t$ over the reference length $|x| + C(|x|+1)^e$.
\end{lemma}

\begin{lemma}[Clean install call of a native function]
\lean{Complexity.clNative_cleanCall}
\label{Complexity.clNative_cleanCall}
\leanok
\uses{Complexity.PolyTimeComputable, Turing.Cfg.ofWords, Turing.FinTM, Turing.FinTM.computesInTime_iff, Turing.FinTM.exists_installCallTM, Turing.MultiTapeTM.output_length_le, Turing.MultiTapeTM.runFrom, Turing.stateWord}
\difficulty{4}
For every polynomial-time computable string function $f$, there are a machine with at least one work tape, entry and exit controls, and constants $a, d$ such that for every physical input and argument word, the machine started at the entry over the argument's state word reaches, at a strictly positive first exit within $a(|arg|+1)^d$ steps, exactly the state word of $f$ of the argument, with all scratch cells and heads restored after each request.
\end{lemma}

\begin{definition}[Retained state of the visit producer]
\lean{Complexity.clVisitState}
\label{Complexity.clVisitState}
\leanok
\uses{Turing.pairEncode}
For an instance $x$, a time $t$, and accumulated rows, the producer's internal state is the nested pair encoding of $x$, the unary word of $t$ tokens, and the stored rows; it is the loop state, not the final emitter state.
\end{definition}

\begin{definition}[One producer time step]
\lean{Complexity.clVisitStep}
\label{Complexity.clVisitStep}
\leanok
\uses{Complexity.clHeaderField, Complexity.clHeaderTail, Complexity.clVisitRow, Turing.FinTM, Turing.pairEncode}
One internal step maps a retained state to the state with the same instance, the unary clock extended by one token, and the stored rows extended by the complete visit row of the current instance and time; every operation is a native word computation.
\end{definition}

\begin{lemma}[The producer step is native]
\lean{Complexity.clVisitStep_native}
\label{Complexity.clVisitStep_native}
\leanok
\uses{Complexity.PolyTimeComputable, Complexity.PolyTimeComputable.comp, Complexity.clHeaderField_native, Complexity.clHeaderTail_native, Complexity.clNative_append, Complexity.clNative_linear, Complexity.clNative_pair, Complexity.clVisitRow_native, Complexity.clVisitStep, Turing.FinTM, Turing.FinTM.computesFunInTime_const}
\difficulty{3}
For every source machine and certificate parameters $C$ and $e$, the producer's one-step function on retained state words is polynomial-time computable, including retained copies, every per-tape search, and the clock increment.
\end{lemma}

\begin{definition}[Ordered table of last-visit rows]
\lean{Complexity.clVisitRows}
\label{Complexity.clVisitRows}
\leanok
\uses{Complexity.clVisitRow, Turing.FinTM, Turing.pairEncode}
For a source machine, parameters $C$ and $e$, an instance $x$, and a count, the table concatenates the visit rows of times $0, 1, \dots$ in order; time zero is present and the completed horizon table has $T + 1$ rows.
\end{definition}

\begin{lemma}[The step advances the retained state]
\lean{Complexity.clVisitStep_state}
\label{Complexity.clVisitStep_state}
\leanok
\uses{Complexity.clHeaderField, Complexity.clHeaderTail, Complexity.clVisitRow, Complexity.clVisitState, Complexity.clVisitStep, Turing.FinTM, Turing.pairDecode_pairEncode, Turing.pairEncode}
\difficulty{1}
For every source machine, parameters, instance $x$, stored rows, and time $t$, one producer step on the retained state at time $t$ yields the retained state at time $t+1$ whose rows are extended by the visit row of time $t$ in fixed work-tape order.
\end{lemma}

\begin{lemma}[The orbit is the ordered table]
\lean{Complexity.clVisitStep_orbit}
\label{Complexity.clVisitStep_orbit}
\leanok
\uses{Complexity.clVisitRows, Complexity.clVisitState, Complexity.clVisitStep, Complexity.clVisitStep_state, Turing.FinTM}
\difficulty{3}
For every source machine, parameters, instance $x$, and count $t$, iterating the producer step $t$ times from the empty retained state yields the retained state at time $t$ holding exactly the ordered table of the first $t$ visit rows.
\end{lemma}

\begin{lemma}[Clean call of the producer step]
\lean{Complexity.clVisitStep_cleanCall}
\label{Complexity.clVisitStep_cleanCall}
\leanok
\uses{Complexity.clNative_cleanCall, Complexity.clVisitStep, Complexity.clVisitStep_native, Turing.Cfg.ofWords, Turing.FinTM, Turing.MultiTapeTM.runFrom, Turing.stateWord}
\difficulty{1}
For every source machine and parameters $C$ and $e$, there are a machine, entry and exit controls, and constants such that each producer step has a clean reusable native call: a strictly positive first exit within a polynomial of the argument length, returning exactly the stepped state word with all scratch restored. The variable number of such calls is not yet assembled here.
\end{lemma}

\begin{definition}[Unary-bounded repetition controller]
\lean{Complexity.clRepeatTM}
\label{Complexity.clRepeatTM}
\leanok
\uses{Complexity.clSlotAction, Turing.FinTM, Turing.FinTM.controlAction}
Given a clean-call machine with entry and exit controls, the repetition controller tests a protected unary clock, and for each live token executes one complete clean call on the shared argument bank, then dispatches back; a fresh begin phase executes the call's first transition before testing the positive return, so entry and exit controls may coincide. On a blank clock it completes silently.
\end{definition}

\begin{definition}[Repetition seam]
\lean{Complexity.clRepeatCfg}
\label{Complexity.clRepeatCfg}
\leanok
\uses{Complexity.clRepeatTM, Turing.Cfg, Turing.FinTM, Turing.FinTM.bufferTape, Turing.stateWord}
For a controller state, an argument word, a clock budget $N$, and a logical cursor $j$, the repetition seam holds the clean argument bank at blank scratch with the state word installed, and the protected unary clock with its head at $j$; the physical output is empty.
\end{definition}

\begin{lemma}[Clean calls agree with the protected-clock seam]
\lean{Complexity.clRepeat_frame}
\label{Complexity.clRepeat_frame}
\leanok
\uses{Complexity.clRepeatCfg, Complexity.clRepeatTM, Complexity.clSlotCfg, Turing.Cfg.ofWords, Turing.FinTM, Turing.FinTM.bufferTape, Turing.stateWord}
\difficulty{3}
For every clean-call machine, controls, input, state word, clock budget, cursor, and inner control, relocating the configuration assembled over the state word into the repetition controller's slots, with the protected clock retained, yields exactly the repetition seam at that inner control.
\end{lemma}

\begin{lemma}[One relocated clean call]
\lean{Complexity.clRepeat_call}
\label{Complexity.clRepeat_call}
\leanok
\uses{Complexity.clRepeatCfg, Complexity.clRepeatTM, Complexity.clRepeat_frame, Complexity.clSlotAction, Complexity.clSlotCfg, Complexity.clSlot_apply, Complexity.clSlot_run, Turing.Action.apply, Turing.Cfg.inputSymbol, Turing.Cfg.ofWords, Turing.Cfg.workTapeSymbols, Turing.FinTM, Turing.FinTM.bufferTape, Turing.MultiTapeTM.runFrom, Turing.MultiTapeTM.runFrom_succ_eq_step, Turing.MultiTapeTM.step, Turing.stateWord}
\difficulty{5}
Suppose a clean call runs from the entry over a state word to its exit over a new state word in a strictly positive time $t$, never showing the exit control strictly before $t$. Then the repetition controller, started in its begin phase over the same bank, reaches in exactly $t$ steps the seam at the exit control with the new state word; the first source action is executed before any return test, so an entry equal to the exit is handled, and the unary clock is preserved.
\end{lemma}

\begin{lemma}[One clocked repetition round]
\lean{Complexity.clRepeat_round}
\label{Complexity.clRepeat_round}
\leanok
\uses{Complexity.clRepeatCfg, Complexity.clRepeatTM, Complexity.clRepeat_call, Turing.Action.apply, Turing.Cfg.ofWords, Turing.Cfg.workTapeSymbols, Turing.FinTM, Turing.FinTM.bufferTape, Turing.FinTM.bufferTape_nat, Turing.FinTM.controlAction_apply, Turing.MultiTapeTM.runFrom, Turing.MultiTapeTM.runFrom_succ_eq_step, Turing.MultiTapeTM.runFrom_succ_eq_step', Turing.MultiTapeTM.step, Turing.moveInputPos_zero, Turing.stateWord}
\difficulty{4}
Suppose a clean call runs from the entry over a state word to its exit over a new state word in a strictly positive time $t$, never showing the exit strictly before $t$, and let the clock cursor be $j < N$. Then the controller, started at its test phase, advances the clock by one token, executes the entire positive clean call, and performs one return dispatch, reaching the test phase over the new state word at cursor $j + 1$ in exactly $t + 2$ steps; calls returning empty words are included.
\end{lemma}

\begin{lemma}[Bounded repetition of clean calls]
\lean{Complexity.clRepeat_complete}
\label{Complexity.clRepeat_complete}
\leanok
\uses{Complexity.clRepeatCfg, Complexity.clRepeatTM, Complexity.clRepeat_round, Turing.Action.apply, Turing.Cfg.ofWords, Turing.Cfg.workTapeSymbols, Turing.FinTM, Turing.FinTM.bufferTape_nat, Turing.FinTM.controlAction_apply, Turing.MultiTapeTM.runFrom, Turing.MultiTapeTM.runFrom_add, Turing.MultiTapeTM.runFrom_succ_eq_step', Turing.moveInputPos_zero, Turing.stateWord}
\difficulty{5}
Suppose every iterate $j < N$ of a word function $f$ has a clean call of positive duration at most $B$ from the controller's bank. Then within $N(B+2) + 1$ steps the controller, started at its test phase over the initial word with an $N$-token clock, reaches its completed control holding exactly the $N$-fold iterate of $f$; every call starts from the fully restored seam and the final blank clock executes one silent completion.
\end{lemma}

\begin{definition}[Argument words of a repetition]
\lean{Complexity.clRepeatWords}
\label{Complexity.clRepeatWords}
\leanok
\uses{Turing.FinTM, Turing.stateWord}
For a clean-call machine, an initial state word, and a count $N$, the repetition argument words are the call's state word with blank scratch, followed by the unary clock of $N$ tokens in a separate protected field.
\end{definition}

\begin{lemma}[Prepared repetition seam]
\lean{Complexity.clRepeat_initial}
\label{Complexity.clRepeat_initial}
\leanok
\uses{Complexity.clRepeatCfg, Complexity.clRepeatTM, Complexity.clRepeatWords, Turing.Cfg.ofWords, Turing.FinTM}
\difficulty{2}
For every clean-call machine, controls, input, state word, and count $N$, the configuration assembled from the repetition argument words at the controller's start equals the repetition seam at the test phase with cursor zero.
\end{lemma}

\begin{definition}[Repetition with physical state replay]
\lean{Complexity.clRepeatOutputTM}
\label{Complexity.clRepeatOutputTM}
\leanok
\uses{Complexity.clOutputTM, Complexity.clPreparedTM, Complexity.clRepeatTM, Turing.FinTM}
The repetition-output machine loads its argument bank natively from self-delimiting input, runs the unary-bounded repetition to completion, and then physically replays the completed state word from the first scratch tape as its output.
\end{definition}

\begin{lemma}[Complete native repetition ledger]
\lean{Complexity.clRepeatOutput_compute}
\label{Complexity.clRepeatOutput_compute}
\leanok
\uses{Complexity.clFields, Complexity.clFields_width, Complexity.clOutputAt_compute, Complexity.clPreparedCfg, Complexity.clPreparedTM, Complexity.clPrepared_idle, Complexity.clPrepared_run, Complexity.clRepeatCfg, Complexity.clRepeatOutputTM, Complexity.clRepeatTM, Complexity.clRepeatWords, Complexity.clRepeat_complete, Complexity.clRepeat_initial, Complexity.clSlotCfg, Turing.Cfg.ofWords, Turing.FinTM, Turing.FinTM.ComputesInTime, Turing.FinTM.ComputesInTime.mono, Turing.MultiTapeTM.runFrom, Turing.stateWord}
\difficulty{5}
Suppose every argument has a clean call computing one application of a word function $f$ within $a(|arg|+1)^d$ steps, and every iterate of $f$ from the initial word $s$ up to $N$ has length at most $W$. Then the repetition-output machine computes the $N$-fold iterate of $f$ at $s$ from the packed repetition argument within $2|x| + 9 + (k+1)(3|x|+7) + N(a(W+1)^d + 2) + W$ steps, where $x$ is the argument and $k$ the call machine's tape count; input parsing, all positive calls, every dispatch, and the final replay are charged.
\end{lemma}

\begin{lemma}[Storage envelope of a field list]
\lean{Complexity.clFields_size}
\label{Complexity.clFields_size}
\leanok
\uses{Complexity.clFields, Turing.pairEncode}
\difficulty{3}
For every list of words with all payloads of length at most $W$, the packed encoding has length at most the list length times $2W + 2$; each field stores its doubled payload plus two separator bits.
\end{lemma}

\begin{lemma}[Size of one visit code]
\lean{Complexity.clVisitCode_size}
\label{Complexity.clVisitCode_size}
\leanok
\uses{Complexity.clElapsed_width, Complexity.clPrev_spec, Complexity.clVisitCode, Complexity.prevVisit, Turing.FinTM, Turing.pairEncode}
\difficulty{3}
For every source machine, reference length $m$, time $t$, and tape $\tau$, the optional-time code of the greatest strictly earlier visit has length at most $t + 2$: a present predecessor is strictly below the query time, and absence at time zero is included.
\end{lemma}

\begin{lemma}[Size of one visit row]
\lean{Complexity.clVisitRow_size}
\label{Complexity.clVisitRow_size}
\leanok
\uses{Complexity.clFields_size, Complexity.clVisitCode, Complexity.clVisitCode_size, Complexity.clVisitRow, Complexity.clVisitRow_exact, Complexity.prevVisit, Turing.FinTM, Turing.pairEncode}
\difficulty{3}
For every source machine, parameters $C$ and $e$, instance $x$, and times $t \le N$, the packed visit row at time $t$ has length at most $k(2N + 6)$, where $k$ is the number of work tapes; a delimiter is stored even for an absent visit.
\end{lemma}

\begin{lemma}[Size of the accumulated table]
\lean{Complexity.clVisitRows_size}
\label{Complexity.clVisitRows_size}
\leanok
\uses{Complexity.clVisitRow_size, Complexity.clVisitRows, Turing.FinTM}
\difficulty{3}
For every source machine, parameters, instance, horizon $N$, and count $t \le N$, the ordered table of the first $t$ visit rows has length at most $t \cdot k(2N+6)$, quadratic in the unary outer horizon.
\end{lemma}

\begin{lemma}[Quadratic bound on the retained state]
\lean{Complexity.clVisitState_size}
\label{Complexity.clVisitState_size}
\leanok
\uses{Complexity.clVisitRows, Complexity.clVisitRows_size, Complexity.clVisitState, Complexity.clVisitStep, Complexity.clVisitStep_orbit, Turing.FinTM, Turing.pairEncode}
\difficulty{4}
For every source machine, parameters, instance $x$ of length at most $L$, and horizon $N \le L$, every iterate up to $N$ of the producer step from the empty retained state has length at most $(8 + 8k)(L+1)^2$, where $k$ is the number of work tapes.
\end{lemma}

\begin{lemma}[Polynomial envelope of the repetition]
\lean{Complexity.clRepeat_budget}
\label{Complexity.clRepeat_budget}
\leanok
\difficulty{5}
For all naturals with $N \le L$, the complete repetition ledger, preparation $2L + 9 + k(3L+7)$, plus $N$ calls each costing the clean-call polynomial evaluated at the quadratic state bound $K(L+1)^2$, plus the final replay $K(L+1)^2$, is at most $(13 + 7k + a(K+1)^d + K)(L+1)^{2d+3}$.
\end{lemma}

\begin{lemma}[Repetition arguments are native]
\lean{Complexity.clRepeatArgument_native}
\label{Complexity.clRepeatArgument_native}
\leanok
\uses{Complexity.PolyTimeComputable, Complexity.clFields, Complexity.clNative_fields, Complexity.clNative_linear, Complexity.clRepeatWords, Turing.FinTM, Turing.FinTM.computesFunInTime_const, Turing.stateWord}
\difficulty{3}
If an initial-state producer and a unary-count producer are polynomial-time computable, then so is the function packing the repetition argument: the produced state word, blank call scratch, and the produced unary clock.
\end{lemma}

\begin{definition}[The producer's exact horizon]
\lean{Complexity.clProducerHorizon}
\label{Complexity.clProducerHorizon}
\leanok
For parameters $C, e, c, A, d$ and an input length $n$, the producer horizon is $c\,(A(n + C(n+1)^e + 1)^d + 1)^2$, the inherited verifier time bound at the reference length; it remains distinct from the emitter's fuel.
\end{definition}

\begin{lemma}[The unary producer clock is native]
\lean{Complexity.clProducerClock_native}
\label{Complexity.clProducerClock_native}
\leanok
\uses{Complexity.PolyTimeComputable, Complexity.PolyTimeComputable.comp, Complexity.clHeaderTail, Complexity.clHeaderTail_native, Complexity.clNative_append, Complexity.clNative_linear, Complexity.clPrepHeader, Complexity.clPrepHeader_native, Complexity.clProducerHorizon, Turing.FinTM.computesFunInTime_const, Turing.pairDecode_pairEncode}
\difficulty{3}
For all parameters $C, e, c, A, d$, the function sending an input $x$ to the unary word of the producer horizon at $|x|$ plus one true token is polynomial-time computable; the horizon is read off the inherited arithmetic header, and the extra token accounts for the inclusive final row.
\end{lemma}

\begin{lemma}[The complete visit table is native]
\lean{Complexity.clVisitRows_native}
\label{Complexity.clVisitRows_native}
\leanok
\uses{Complexity.PolyTimeComputable, Complexity.PolyTimeComputable.comp, Complexity.clFields, Complexity.clFields_width, Complexity.clHeaderTail, Complexity.clHeaderTail_native, Complexity.clNative_image, Complexity.clNative_linear, Complexity.clNative_pair, Complexity.clProducerClock_native, Complexity.clProducerHorizon, Complexity.clRepeatArgument_native, Complexity.clRepeatOutputTM, Complexity.clRepeatOutput_compute, Complexity.clRepeatWords, Complexity.clRepeat_budget, Complexity.clVisitRows, Complexity.clVisitState, Complexity.clVisitState_size, Complexity.clVisitStep, Complexity.clVisitStep_cleanCall, Complexity.clVisitStep_orbit, Complexity.polyTimeComputable_id, Turing.FinTM, Turing.FinTM.ComputesInTime.mono, Turing.FinTM.computesFunInTime_const, Turing.pairDecode_pairEncode, Turing.pairEncode, Turing.stateWord}
\difficulty{5}
For every source machine and parameters $C, e, c, A, d$, the function sending an instance $x$ to the ordered table of visit rows for all times up to the producer horizon plus one is polynomial-time computable: the retained instance and empty table start the loop, a clean call executes each full row step, the repetition controller runs the exact unary number of calls with a quadratic bound on all intermediate state words, and a native projection returns precisely the stored rows.
\end{lemma}

\begin{definition}[Complete packed preparation data]
\lean{Complexity.clPackedRecords}
\label{Complexity.clPackedRecords}
\leanok
\uses{Complexity.clPrepHeader, Complexity.clProducerHorizon, Complexity.clRecords, Complexity.clVisitRows, Turing.FinTM, Turing.pairEncode}
For a source machine, parameters $C, e, c, A, d$, and an instance $x$, the packed record is the nested pair of: the untouched exact arithmetic header of $x$; the inclusive chronological movement-count trajectory over the all-false reference input of length $|x| + C(|x|+1)^e$ at the producer horizon plus one; and the inclusive time-ordered, work-tape-ordered table of canonical last-visit codes.
\end{definition}

\begin{lemma}[The packed producer is native]
\lean{Complexity.clPackedRecords_native}
\label{Complexity.clPackedRecords_native}
\leanok
\uses{Complexity.PolyTimeComputable, Complexity.clNative_pair, Complexity.clPackedRecords, Complexity.clPrepHeader_native, Complexity.clRecords_native, Complexity.clVisitRows_native, Turing.FinTM}
\difficulty{1}
For every source machine and parameters $C, e, c, A, d$, the function sending an instance to its complete packed preparation data, header, trajectory, and visit table, is polynomial-time computable; every retained copy and capture is charged by the native composition calculus.
\end{lemma}

\begin{lemma}[One machine returns the packed result]
\lean{Complexity.clPackedRecords_machine}
\label{Complexity.clPackedRecords_machine}
\leanok
\uses{Complexity.clPackedRecords, Complexity.clPackedRecords_native, Turing.FinTM, Turing.FinTM.ComputesFunInTime, Turing.FinTM.computesInTime_iff, Turing.MultiTapeTM.output_length_le}
\difficulty{2}
For every source machine and parameters, there are a machine and constants $a, r$ such that the machine computes the complete packed preparation data within $a(n+1)^r$ steps on inputs of length $n$, and the packed answer itself has length at most the same bound.
\end{lemma}

\begin{lemma}[Cutting a padded run at its first halt]
\lean{Complexity.clA5_call_first_halt}
\label{Complexity.clA5_call_first_halt}
\leanok
\uses{Turing.Cfg, Turing.Cfg.Halted, Turing.MultiTapeTM, Turing.MultiTapeTM.runFrom, Turing.MultiTapeTM.runFrom_add, Turing.MultiTapeTM.runFrom_of_halt, Turing.MultiTapeTM.runFrom_zero}
\difficulty{4}
Let a machine run from a live configuration to a halted configuration in $t$ steps. Then there is a strictly positive time $u \le t$ at which the run first halts, and the configuration at $u$ equals the configuration at $t$; padding by absorbed halted steps changes no endpoint field.
\end{lemma}

\begin{definition}[Padding an action into a shared bank]
\lean{Complexity.clA5PadAction}
\label{Complexity.clA5PadAction}
\leanok
\uses{Turing.Action}
An action on $k$ tapes is padded to a larger bank of $K$ tapes by keeping its input move, output emission, and successor state, applying its writes and moves on the first $k$ tapes, and leaving every additional tape stationary.
\end{definition}

\begin{definition}[Padded module over the shared bank]
\lean{Complexity.clA5PadTM}
\label{Complexity.clA5PadTM}
\leanok
\uses{Complexity.clA5PadAction, Turing.FinTM}
For a machine with $k \le K$ tapes, the padded machine on $K$ tapes runs the same control, reading only the first $k$ work symbols and padding each action; all modules can therefore share tape zero and one common bank.
\end{definition}

\begin{definition}[Padded configuration]
\lean{Complexity.clA5PadCfg}
\label{Complexity.clA5PadCfg}
\leanok
\uses{Turing.Cfg}
The padded form of a configuration on $k$ tapes is the configuration on $K$ tapes that keeps its state, input position, and output, places its tapes and heads on the first $k$ slots, and leaves every additional tape genuinely blank with head zero.
\end{definition}

\begin{lemma}[Padding commutes with one action]
\lean{Complexity.clA5Pad_apply}
\label{Complexity.clA5Pad_apply}
\leanok
\uses{Complexity.clA5PadAction, Complexity.clA5PadCfg, Turing.Action, Turing.Action.apply, Turing.Cfg}
\difficulty{3}
For every action and configuration on the source tapes, applying the padded action to the padded configuration equals padding the result of the original application, including writes at arbitrary cells; inactive slots neither write nor move.
\end{lemma}

\begin{lemma}[Padded runs project to source runs]
\lean{Complexity.clA5Pad_run}
\label{Complexity.clA5Pad_run}
\leanok
\uses{Complexity.clA5PadAction, Complexity.clA5PadCfg, Complexity.clA5PadTM, Complexity.clA5Pad_apply, Turing.Action.apply, Turing.Cfg, Turing.Cfg.inputSymbol, Turing.Cfg.workTapeSymbols, Turing.FinTM, Turing.MultiTapeTM.runFrom, Turing.MultiTapeTM.runFrom_comm_of_step, Turing.MultiTapeTM.step}
\difficulty{4}
For every machine with $k \le K$ tapes, source configuration, and duration $t$, running the padded machine for $t$ steps from the padded configuration equals padding the source's $t$-step run; the padded source observes exactly its original work symbols.
\end{lemma}

\begin{lemma}[Padded canonical seams]
\lean{Complexity.clA5Pad_seam}
\label{Complexity.clA5Pad_seam}
\leanok
\uses{Complexity.clA5PadCfg, Turing.Cfg.ofWords, Turing.FinTM.bufferTape, Turing.stateWord}
\difficulty{3}
For a positive source tape count, every control, stored word, and output, padding the configuration assembled over the source's state word yields exactly the configuration assembled over the larger bank's state word with the same output.
\end{lemma}

\begin{definition}[Release wrapper of a clean call]
\lean{Complexity.clA5StopTM}
\label{Complexity.clA5StopTM}
\leanok
\uses{Turing.FinTM, Turing.FinTM.controlAction}
Given a machine with entry and exit controls, the wrapped machine carries one release flag: it executes the underlying transitions, raising the flag after the first step, and physically halts exactly when the flag is set and the control is the exit; an entry equal to the exit therefore still executes at least one step.
\end{definition}

\begin{definition}[Release-flag configuration]
\lean{Complexity.clA5StopCfg}
\label{Complexity.clA5StopCfg}
\leanok
\uses{Turing.Cfg}
A configuration of the underlying machine embeds into the wrapped machine by pairing its control with a release flag; all physical fields, tapes, heads, input position, and output, are unchanged.
\end{definition}

\begin{lemma}[The wrapper steps with the module]
\lean{Complexity.clA5Stop_step}
\label{Complexity.clA5Stop_step}
\leanok
\uses{Complexity.clA5StopCfg, Complexity.clA5StopTM, Turing.Cfg, Turing.FinTM, Turing.MultiTapeTM.step}
\difficulty{3}
For every flag value and live configuration whose control is not the designated exit when the flag is set, one step of the release wrapper equals the module's physical step with the flag raised.
\end{lemma}

\begin{lemma}[Live endpoints have live prefixes]
\lean{Complexity.clA5_live_prefix}
\label{Complexity.clA5_live_prefix}
\leanok
\uses{Turing.Cfg, Turing.MultiTapeTM, Turing.MultiTapeTM.runFrom, Turing.MultiTapeTM.runFrom_add, Turing.MultiTapeTM.runFrom_of_halt}
\difficulty{2}
For every machine, configuration, and time $t$ whose configuration at $t$ is live, every configuration at a time $u \le t$ is live as well, since halting is absorbing.
\end{lemma}

\begin{lemma}[Stopping one step after the first exit]
\lean{Complexity.clA5Stop_clean}
\label{Complexity.clA5Stop_clean}
\leanok
\uses{Complexity.clA5StopCfg, Complexity.clA5StopTM, Complexity.clA5Stop_step, Complexity.clA5_live_prefix, Turing.Action.apply, Turing.Cfg.ofWords, Turing.FinTM, Turing.FinTM.controlAction, Turing.MultiTapeTM.runFrom, Turing.MultiTapeTM.runFrom_succ_eq_step', Turing.MultiTapeTM.step, Turing.stateWord}
\difficulty{5}
Suppose a clean call runs from its entry over an argument word to its exit over a result word with some emitted output, in positive time $t$, never showing the exit strictly between $0$ and $t$. Then the release wrapper, started unreleased at the entry, halts physically at time $t + 1$ with exactly the result's canonical seam, the same emitted output, and no further transition.
\end{lemma}

\begin{definition}[Control of the five-module host]
\lean{Complexity.CLA5HostState}
\label{Complexity.CLA5HostState}
\leanok
\uses{Turing.FinTM}
The host's finite control is either the unique loop anchor or a state of one of the five modules: native input copy, packed-record installation, cursor initialization, emission, and cursor update.
\end{definition}

\begin{definition}[Module return destinations]
\lean{Complexity.clA5HostNext}
\label{Complexity.clA5HostNext}
\leanok
Module zero continues to module one, module one to module two, and module three to module four; modules two and four return to the unique loop anchor.
\end{definition}

\begin{definition}[Entry state of a module]
\lean{Complexity.clA5HostEntry}
\label{Complexity.clA5HostEntry}
\leanok
\uses{Complexity.CLA5HostState, Turing.FinTM}
For each of the five modules, its entry state inside the host is the module's own start control tagged by its index.
\end{definition}

\begin{definition}[Resolved return of a module]
\lean{Complexity.clA5HostRet}
\label{Complexity.clA5HostRet}
\leanok
\uses{Complexity.CLA5HostState, Complexity.clA5HostEntry, Complexity.clA5HostNext, Turing.FinTM}
A module's return state is the entry of its successor module when one exists, and the loop anchor otherwise; resolving the return performs no tape operation.
\end{definition}

\begin{definition}[The five-module emission host]
\lean{Complexity.clA5HostTM}
\label{Complexity.clA5HostTM}
\leanok
\uses{Complexity.CLA5HostState, Complexity.clA5HostEntry, Complexity.clA5HostRet, Complexity.clA5PadTM, Turing.FinTM, Turing.FinTM.controlAction, Turing.emitAction}
Over one shared padded work bank, the host runs the five modules through their call chain: the anchor dispatches into the emission module, and each module's transitions are forwarded with halts redirected to the module's resolved return; the clean-call contracts restore all scratch before each next call.
\end{definition}

\begin{lemma}[Forwarded seams of the host]
\lean{Complexity.clA5Host_seam}
\label{Complexity.clA5Host_seam}
\leanok
\uses{Complexity.CLA5HostState, Complexity.clA5HostRet, Complexity.clA5PadCfg, Complexity.clA5Pad_seam, Turing.Cfg, Turing.Cfg.ofWords, Turing.FinTM, Turing.emitCfg, Turing.stateWord}
\difficulty{3}
For every module with a positive tape count, control, stored word, output prefix, and optional state, forwarding the padded clean configuration yields the host's canonical seam over the shared bank's state word, with output exactly the prefix followed by the module's output.
\end{lemma}

\begin{lemma}[Running one clean module inside the host]
\lean{Complexity.clA5Host_call}
\label{Complexity.clA5Host_call}
\leanok
\uses{Complexity.CLA5HostState, Complexity.clA5HostEntry, Complexity.clA5HostRet, Complexity.clA5HostTM, Complexity.clA5Host_seam, Complexity.clA5PadCfg, Complexity.clA5PadTM, Complexity.clA5Pad_run, Turing.Cfg.Halted, Turing.Cfg.ofWords, Turing.FinTM, Turing.MultiTapeTM.runFrom, Turing.emitCfg, Turing.emit_run, Turing.stateWord}
\difficulty{5}
Suppose a module runs from its start over an argument word to a halted canonical seam with a result word and output chunk in time $t$, never halting earlier. Then the host, started at that module's entry over the shared bank with an output prefix, never shows the loop anchor before $t$, and at $t$ sits at the module's resolved return over the result word with the output prefix extended by exactly the module's chunk.
\end{lemma}

\begin{lemma}[Joining two guarded segments]
\lean{Complexity.clA5_guard_add}
\label{Complexity.clA5_guard_add}
\leanok
\uses{Turing.Cfg, Turing.MultiTapeTM, Turing.MultiTapeTM.runFrom, Turing.MultiTapeTM.runFrom_add}
\difficulty{3}
If a run avoids an anchor control for $a$ steps, reaches a configuration $d$, and the run from $d$ avoids the anchor for $b$ steps, then the whole run avoids the anchor for $a + b$ steps.
\end{lemma}

\begin{definition}[Uniform clean-module contract]
\lean{Complexity.CLA5Clean}
\label{Complexity.CLA5Clean}
\leanok
\uses{Turing.Cfg.Halted, Turing.Cfg.ofWords, Turing.FinTM, Turing.MultiTapeTM.runFrom, Turing.stateWord}
A machine satisfies the clean contract for a result map, an emission map, and a budget when it has a positive tape count and, for every native input and argument word, it first halts at some positive time within the budget of the argument length, at the canonical seam holding the result word and having emitted exactly the prescribed chunk.
\end{definition}

\begin{lemma}[Clean calls give uniform first-halting modules]
\lean{Complexity.clA5Clean_stop}
\label{Complexity.clA5Clean_stop}
\leanok
\uses{Complexity.CLA5Clean, Complexity.clA5StopTM, Complexity.clA5Stop_clean, Complexity.clA5_call_first_halt, Turing.Cfg.ofWords, Turing.FinTM, Turing.MultiTapeTM.runFrom, Turing.stateWord}
\difficulty{3}
If a machine has, for every input and argument, a positive first exit within a budget $B$ at the canonical seam with a result word and emitted chunk, then its release wrapper satisfies the uniform clean contract with budget $B + 1$; the minimal-halt cut retains the entire clean configuration.
\end{lemma}

\begin{lemma}[Degree bound of the bridge overhead]
\lean{Complexity.clA5_bridge_bound}
\label{Complexity.clA5_bridge_bound}
\leanok
\difficulty{4}
For all naturals with a result length at most $A(n+1)^d$, the bridge cost $c\,(A(n+1)^d + n + \mathrm{len} + 1) + 1$ is at most $(c(2A+1) + 1)(n+1)^{d+1}$; the final unit pays for the stopping transition.
\end{lemma}

\begin{lemma}[Clean install module of a native function]
\lean{Complexity.clA5Clean_install}
\label{Complexity.clA5Clean_install}
\leanok
\uses{Complexity.CLA5Clean, Complexity.PolyTimeComputable, Complexity.clA5Clean_stop, Complexity.clA5StopTM, Complexity.clA5_bridge_bound, Turing.FinTM.computesInTime_iff, Turing.FinTM.exists_installCallTM, Turing.MultiTapeTM.output_length_le}
\difficulty{4}
Every polynomial-time computable string function $f$ has a first-halting clean module that replaces the argument word by $f$ of the argument, emits nothing, and runs within $A(n+1)^d$ steps of the argument length for some constants $A, d$.
\end{lemma}

\begin{lemma}[Clean emit module of a native function]
\lean{Complexity.clA5Clean_emit}
\label{Complexity.clA5Clean_emit}
\leanok
\uses{Complexity.CLA5Clean, Complexity.PolyTimeComputable, Complexity.clA5Clean_stop, Complexity.clA5StopTM, Complexity.clA5_bridge_bound, Turing.FinTM.computesInTime_iff, Turing.FinTM.exists_emitCallTM, Turing.MultiTapeTM.output_length_le}
\difficulty{4}
Every polynomial-time computable string function $f$ has a first-halting clean module that preserves the argument word, emits exactly $f$ of the argument as output, and runs within a polynomial budget of the argument length, including cleanup and the final halt.
\end{lemma}

\begin{lemma}[Clean contracts specialize to a call site]
\lean{Complexity.clA5Host_clean}
\label{Complexity.clA5Host_clean}
\leanok
\uses{Complexity.CLA5Clean, Complexity.clA5HostEntry, Complexity.clA5HostRet, Complexity.clA5HostTM, Complexity.clA5Host_call, Turing.Cfg.ofWords, Turing.FinTM, Turing.MultiTapeTM.runFrom, Turing.stateWord}
\difficulty{2}
If the module at index $i$ of the host satisfies the uniform clean contract for result map $f$, emission map $g$, and budget $B$, then from the host entry of $i$ over any argument and output prefix, the host reaches, within $B$ of the argument length and without touching the loop anchor, the resolved return over $f$ of the argument with the output extended by exactly $g$ of the argument.
\end{lemma}

\begin{lemma}[Startup copy as a first-halting module]
\lean{Complexity.clA5Copy_clean}
\label{Complexity.clA5Copy_clean}
\leanok
\uses{Complexity.clA5StopTM, Complexity.clA5Stop_clean, Complexity.clA5_call_first_halt, Complexity.clInputTM, Complexity.clInput_first, Turing.Cfg.Halted, Turing.Cfg.init, Turing.Cfg.ofWords, Turing.FinTM.bufferTape, Turing.MultiTapeTM.initCfg, Turing.MultiTapeTM.runFrom, Turing.stateWord}
\difficulty{4}
For every input word $x$, the release-wrapped native input copier, started over the blank state word, first halts at a positive time at most $2|x| + 3$ at the canonical seam whose stored word is exactly $x$, with silent output and unchanged heads.
\end{lemma}

\begin{lemma}[Install module of the packed producer]
\lean{Complexity.clA5Packed_install}
\label{Complexity.clA5Packed_install}
\leanok
\uses{Complexity.CLA5Clean, Complexity.clA5Clean_install, Complexity.clPackedRecords, Complexity.clPackedRecords_machine, Turing.FinTM}
\difficulty{1}
For every source machine and parameters, there are a clean module and constants such that the module replaces its argument word, the genuine retained native input, by the complete packed preparation data, header, trajectory, and visit table, emitting nothing, within a polynomial budget; the completed producer itself is consumed, not substituted.
\end{lemma}

\begin{definition}[Persistent body word with emission cursor]
\lean{Complexity.clA5Pack}
\label{Complexity.clA5Pack}
\leanok
\uses{Complexity.clPackedRecords, Turing.FinTM, Turing.pairEncode}
For a source machine, parameters, instance $x$, and cursor value $i$, the persistent body word is the pair of the unary cursor of $i$ tokens with the entire unchanged packed producer answer; the original instance remains inside its header.
\end{definition}

\begin{lemma}[Install module of the initial cursor]
\lean{Complexity.clA5Cursor_install}
\label{Complexity.clA5Cursor_install}
\leanok
\uses{Complexity.CLA5Clean, Complexity.clA5Clean_install, Complexity.clNative_linear, Turing.FinTM.computesFunInTime_pairEncodeFixed, Turing.pairEncode}
\difficulty{1}
There are a clean module and constants such that the module replaces its argument word by the pair encoding of the empty word with that argument, installing an empty unary emission cursor in front of the packed data, while emitting nothing; neither the header nor either stored table changes.
\end{lemma}

\begin{definition}[Concrete body module bank]
\lean{Complexity.clA5Modules}
\label{Complexity.clA5Modules}
\leanok
\uses{Complexity.clA5StopTM, Complexity.clInputTM, Turing.FinTM}
Given install, cursor, emission, and update machines, the five body modules are: the release-wrapped native input copier, the packed-record installer, the cursor installer, the emission module, and the cursor-update module.
\end{definition}

\begin{lemma}[A shared bank fits all modules]
\lean{Complexity.clA5Modules_bound}
\label{Complexity.clA5Modules_bound}
\leanok
\uses{Complexity.clA5Modules, Complexity.clA5StopTM, Complexity.clInputTM, Turing.FinTM}
\difficulty{2}
For all four supplied machines, every one of the five body modules has at most $1$ plus the sum of the four machines' tape counts many tapes, so one finite shared bank with a real data tape contains every module.
\end{lemma}

\begin{lemma}[Silent startup to the first anchor]
\lean{Complexity.clA5Startup}
\label{Complexity.clA5Startup}
\leanok
\uses{Complexity.clA5Copy_clean, Complexity.clA5Cursor_install, Complexity.clA5HostEntry, Complexity.clA5HostTM, Complexity.clA5Host_call, Complexity.clA5Host_clean, Complexity.clA5Modules, Complexity.clA5Modules_bound, Complexity.clA5Pack, Complexity.clA5Packed_install, Complexity.clA5StopTM, Complexity.clA5_guard_add, Complexity.clInputTM, Complexity.clPackedRecords, Complexity.clPackedRecords_machine, Turing.Cfg.init, Turing.Cfg.ofWords, Turing.FinTM, Turing.FinTM.bufferTape, Turing.MultiTapeTM.initCfg, Turing.MultiTapeTM.runFrom, Turing.MultiTapeTM.runFrom_add, Turing.pairEncode, Turing.stateWord}
\difficulty{5}
For every source machine and parameters, there are install and cursor machines with polynomial constants such that, for arbitrary emission and update modules, the five-module host started from its native initial configuration reaches, within the summed startup budget and without ever showing the loop anchor earlier, the anchor over the persistent body word with cursor zero: the complete packed producer answer behind an empty unary cursor. Startup is silent and never executes the emission or update modules.
\end{lemma}

\begin{lemma}[Constant words are native]
\lean{Complexity.clA5_pt_const}
\label{Complexity.clA5_pt_const}
\leanok
\uses{Complexity.PolyTimeComputable, Complexity.clNative_linear, Turing.FinTM.computesFunInTime_const}
\difficulty{1}
For every fixed word $w$, the constant function emitting $w$ on all inputs is polynomial-time computable from finite control.
\end{lemma}

\begin{lemma}[Polynomial-time conditional]
\lean{Complexity.clA5_pt_cond}
\label{Complexity.clA5_pt_cond}
\leanok
\uses{Complexity.PolyTimeComputable, Turing.FinTM.ComputesInTime.mono, Turing.FinTM.computesFunInTime_cond}
\difficulty{4}
If a single-bit test $p$ and two string functions $f$ and $g$ are polynomial-time computable, then so is the function returning $f$ of the input when the test holds and $g$ of the input otherwise; the three budgets fit one common maximal degree.
\end{lemma}

\begin{lemma}[Short-circuit conjunction of tests]
\lean{Complexity.clA5_pt_and}
\label{Complexity.clA5_pt_and}
\leanok
\uses{Complexity.PolyTimeComputable, Complexity.clA5_pt_cond, Complexity.clA5_pt_const}
\difficulty{2}
If two single-bit tests are polynomial-time computable, then so is their conjunction, evaluated in order with the second test skipped when the first fails.
\end{lemma}

\begin{definition}[Semantics of a one-way transducer]
\lean{Complexity.clA5MapWord}
\label{Complexity.clA5MapWord}
\leanok
For a finite state set with a step function, an optional per-bit emission, and an optional final emission, the transduction of a word processes each bit in order, appending its optional emitted bit and updating the state, and appends the optional final bit at the end; at most one bit is emitted per input bit.
\end{definition}

\begin{definition}[Finite one-way transducer machine]
\lean{Complexity.clA5MapTM}
\label{Complexity.clA5MapTM}
\leanok
\uses{Turing.FinTM}
For a finite state set with step, emission, and finishing functions and a start state, the transducer is a machine with no work tapes that scans its input left to right, emits the per-bit outputs, and halts at the right boundary after the optional final emission.
\end{definition}

\begin{definition}[Scanner frame of the transducer]
\lean{Complexity.clA5MapCfg}
\label{Complexity.clA5MapCfg}
\leanok
\uses{Turing.Cfg}
For an input word, a state, a scanned-prefix length, and an already-emitted output, the scanner frame is the configuration with the input head just past the prefix and no work tapes.
\end{definition}

\begin{lemma}[The transducer matches its semantics]
\lean{Complexity.clA5Map_run}
\label{Complexity.clA5Map_run}
\leanok
\uses{Complexity.clA5MapCfg, Complexity.clA5MapTM, Complexity.clA5MapWord, Turing.Action.apply, Turing.Cfg.ext_zero_tapes, Turing.Cfg.inputSymbol, Turing.FinTM.inputSymbol_at, Turing.MultiTapeTM.runFrom, Turing.MultiTapeTM.runFrom_succ_eq_step, Turing.MultiTapeTM.step, Turing.moveInputPos_pos_of_ne_right}
\difficulty{4}
For every finite transducer, input split into a prefix and an unread suffix, state, and emitted output, running the machine for the suffix length plus one steps halts, and its output is exactly the prior emitted output followed by the recursive transduction of the suffix from the current state.
\end{lemma}

\begin{lemma}[Transducers run in linear time]
\lean{Complexity.clA5Map_poly}
\label{Complexity.clA5Map_poly}
\leanok
\uses{Complexity.PolyTimeComputable, Complexity.clA5MapCfg, Complexity.clA5MapTM, Complexity.clA5MapWord, Complexity.clA5Map_run, Turing.Cfg.ext_zero_tapes, Turing.FinTM.computesInTime_iff, Turing.MultiTapeTM.initCfg}
\difficulty{3}
Every finite one-way word transducer computes its recursive word semantics within time equal to the input length plus one, and is in particular polynomial-time computable.
\end{lemma}

\begin{lemma}[Unary input length is native]
\lean{Complexity.clA5_pt_unaryLength}
\label{Complexity.clA5_pt_unaryLength}
\leanok
\uses{Complexity.PolyTimeComputable, Complexity.clA5MapWord, Complexity.clA5Map_poly}
\difficulty{3}
The function sending a word to the unary word of true bits of the same length is polynomial-time computable, by the one-state transducer replacing every bit by true.
\end{lemma}

\begin{lemma}[Dropping the first bit is native]
\lean{Complexity.clA5_pt_tail}
\label{Complexity.clA5_pt_tail}
\leanok
\uses{Complexity.PolyTimeComputable, Complexity.clA5MapWord, Complexity.clA5Map_poly}
\difficulty{3}
The function removing the first bit of a word, returning the empty word on empty input, is polynomial-time computable, by a two-state finite transduction.
\end{lemma}

\begin{lemma}[Taking the first bit is native]
\lean{Complexity.clA5_pt_head}
\label{Complexity.clA5_pt_head}
\leanok
\uses{Complexity.PolyTimeComputable, Complexity.clA5MapWord, Complexity.clA5Map_poly}
\difficulty{3}
The function keeping at most the first bit of a word is polynomial-time computable, by a two-state finite transduction.
\end{lemma}

\begin{lemma}[Equality to a fixed word is native]
\lean{Complexity.clA5_pt_eq}
\label{Complexity.clA5_pt_eq}
\leanok
\uses{Complexity.PolyTimeComputable, Complexity.PolyTimeComputable.comp, Complexity.clNative_linear, Turing.FinTM.computesFunInTime_ifEq}
\difficulty{2}
For every polynomial-time computable string function $f$ and fixed word $w$, the single-bit test deciding whether $f$ of the input equals $w$ is polynomial-time computable.
\end{lemma}

\begin{lemma}[Native unary-bounded iteration]
\lean{Complexity.clA5Iter_native}
\label{Complexity.clA5Iter_native}
\leanok
\uses{Complexity.PolyTimeComputable, Complexity.clFields, Complexity.clFields_width, Complexity.clNative_cleanCall, Complexity.clNative_image, Complexity.clRepeatArgument_native, Complexity.clRepeatOutputTM, Complexity.clRepeatOutput_compute, Complexity.clRepeatWords, Complexity.clRepeat_budget, Turing.FinTM.ComputesInTime.mono, Turing.stateWord}
\difficulty{5}
Let $f$, an initial-word producer $s$, and a unary count producer $N$ be polynomial-time computable, and suppose every iterate of $f$ up to $N$ of the input has length at most $K((|s| + N + 1))^2$ in the produced values. Then the function sending an input $x$ to the $N(x)$-fold iterate of $f$ at $s(x)$ is polynomial-time computable; each iteration is one clean call inside the real unary repetition host.
\end{lemma}

\begin{lemma}[Iterating a shrinking step]
\lean{Complexity.clA5Iter_shrinking}
\label{Complexity.clA5Iter_shrinking}
\leanok
\uses{Complexity.PolyTimeComputable, Complexity.clA5Iter_native}
\difficulty{3}
If a polynomial-time computable step function never increases word length, then for any polynomial-time initial-word and unary-count producers, the counted iterate is polynomial-time computable; repeated sequential projections and tail scans therefore need no random access.
\end{lemma}

\begin{lemma}[Counted dropping of leading bits]
\lean{Complexity.clA5Drop_native}
\label{Complexity.clA5Drop_native}
\leanok
\uses{Complexity.PolyTimeComputable, Complexity.PolyTimeComputable.comp, Complexity.clA5Iter_shrinking, Complexity.clA5_pt_tail, Complexity.clA5_pt_unaryLength}
\difficulty{3}
If a word producer and a length-witness producer are polynomial-time computable, then so is the function dropping, from the produced word, a number of leading bits equal to the witness's length; each tail call is charged even after the word is exhausted.
\end{lemma}

\begin{lemma}[Decoded components fit the source word]
\lean{Complexity.clA5Pair_sizes}
\label{Complexity.clA5Pair_sizes}
\leanok
\uses{Complexity.clHeaderField, Complexity.clHeaderTail, Turing.pairDecode}
\difficulty{5}
For every word $w$, the first component and the remaining payload of its pair decoding each have length at most $|w|$, including malformed input, where the defaults are empty. The estimate concerns storage only.
\end{lemma}

\begin{lemma}[Variable header-field access is native]
\lean{Complexity.clA5Field_native}
\label{Complexity.clA5Field_native}
\leanok
\uses{Complexity.PolyTimeComputable, Complexity.PolyTimeComputable.comp, Complexity.clA5Iter_shrinking, Complexity.clA5Pair_sizes, Complexity.clA5_pt_unaryLength, Complexity.clHeaderField, Complexity.clHeaderField_native, Complexity.clHeaderTail, Complexity.clHeaderTail_native}
\difficulty{3}
If a header producer and a unary index producer are polynomial-time computable, then so is the function selecting the header field at the produced index: a counted sequence of native tail projections followed by one first projection.
\end{lemma}

\begin{lemma}[Fixed multiples of unary values]
\lean{Complexity.clA5Times_native}
\label{Complexity.clA5Times_native}
\leanok
\uses{Complexity.PolyTimeComputable, Complexity.clA5_pt_const, Complexity.clNative_append}
\difficulty{3}
If a unary value is polynomial-time computable, then for every fixed constant $k$ the unary word of $k$ times that value is polynomial-time computable by $k$ native concatenations; the coefficient lives in finite control.
\end{lemma}

\begin{lemma}[Unary comparison is native]
\lean{Complexity.clA5Le_native}
\label{Complexity.clA5Le_native}
\leanok
\uses{Complexity.PolyTimeComputable, Complexity.clA5Drop_native, Complexity.clA5_pt_eq}
\difficulty{2}
If two unary values are polynomial-time computable, then the single-bit test deciding whether the first is at most the second is polynomial-time computable, by charged subtraction followed by a finite empty-word test.
\end{lemma}

\begin{lemma}[Native serialization of a wired template]
\lean{Complexity.clA5Template_native}
\label{Complexity.clA5Template_native}
\leanok
\uses{Complexity.PolyTimeComputable, Complexity.clA5_pt_const, Complexity.clFragment, Complexity.clNative_append, Std.Sat.CNF.serializeClause, Std.Sat.CNF.serializeLit}
\difficulty{4}
For every fixed CNF table and every wiring whose unary index is polynomial-time computable at each of the finitely many template variables, the function emitting the serialized fragment of the relabeled table is polynomial-time computable; every literal occurrence pays its variable-index cost, and the formula terminator is deliberately absent from the fragment.
\end{lemma}

\begin{lemma}[Unary packed addresses are native]
\lean{Complexity.clA5Address_native}
\label{Complexity.clA5Address_native}
\leanok
\uses{Complexity.PolyTimeComputable, Complexity.clA5Times_native, Complexity.clA5_pt_const, Complexity.clNative_append, Complexity.clPack}
\difficulty{2}
For a fixed block width $B$ and bit position $j$, if the unary input-variable count and the unary time are polynomial-time computable, then so is the unary packed address $m + tB + j$, by native addition and fixed multiplication.
\end{lemma}

\begin{lemma}[Native serialization of one clause group]
\lean{Complexity.clA5Group_native}
\label{Complexity.clA5Group_native}
\leanok
\uses{Complexity.CLFieldCode, Complexity.CLTemplateKind, Complexity.PolyTimeComputable, Complexity.clA5Address_native, Complexity.clA5Template_native, Complexity.clFragment, Complexity.clGroup, Complexity.clWidth, Complexity.clWire, Turing.FinTM}
\difficulty{3}
For a fixed machine, state code, and template kind, if the unary input count, source time, target time, and selected index are polynomial-time computable, then the serialized fragment of the wired clause group is polynomial-time computable; only the unary wiring varies with the input.
\end{lemma}

\begin{definition}[Stored header of the packed word]
\lean{Complexity.clA5StoredHeader}
\label{Complexity.clA5StoredHeader}
\leanok
\uses{Complexity.clHeaderField, Complexity.clHeaderTail}
From the persistent body word, the stored header is the first field of the first payload: the untouched arithmetic header behind the unary emission cursor, selected by native projections.
\end{definition}

\begin{definition}[Retained instance of the packed word]
\lean{Complexity.clA5Instance}
\label{Complexity.clA5Instance}
\leanok
\uses{Complexity.clA5StoredHeader, Complexity.clHeaderField}
From the persistent body word, the retained original instance equals the first field of the stored header, recovered verbatim.
\end{definition}

\begin{definition}[Stored last-round index]
\lean{Complexity.clA5StoredRound}
\label{Complexity.clA5StoredRound}
\leanok
\uses{Complexity.clA5Instance, Complexity.clA5StoredHeader, Complexity.clHeaderTail, Complexity.clLastRound, Turing.FinTM}
From the persistent body word, the stored last-round index is computed from the retained instance's length, the machine's tape count, and the length of the stored unary horizon, header field five; the horizon is a unary length, not an unproved binary-to-unary conversion.
\end{definition}

\begin{lemma}[Stored projections match the producer word]
\lean{Complexity.clA5Stored_exact}
\label{Complexity.clA5Stored_exact}
\leanok
\uses{Complexity.clA5Instance, Complexity.clA5Pack, Complexity.clA5StoredHeader, Complexity.clA5StoredRound, Complexity.clHeaderField, Complexity.clHeaderTail, Complexity.clLastRound, Complexity.clPackedRecords, Complexity.clPrepHeader, Complexity.clProducerHorizon, Turing.FinTM, Turing.pairDecode_pairEncode}
\difficulty{1}
For every source machine, parameters, instance $x$, and cursor $i$, the retained instance of the persistent body word is $x$ itself, and the stored last-round index equals the last round of the schedule at input length $|x|$, tape count, and producer horizon.
\end{lemma}

\begin{lemma}[The stored round index is native]
\lean{Complexity.clA5StoredRound_native}
\label{Complexity.clA5StoredRound_native}
\leanok
\uses{Complexity.PolyTimeComputable, Complexity.PolyTimeComputable.comp, Complexity.clA5Instance, Complexity.clA5StoredHeader, Complexity.clA5StoredRound, Complexity.clA5Times_native, Complexity.clA5_pt_const, Complexity.clA5_pt_unaryLength, Complexity.clHeaderField_native, Complexity.clHeaderTail_native, Complexity.clLastRound, Complexity.clNative_append, Turing.FinTM}
\difficulty{3}
For every source machine, the function sending a persistent body word to the unary word of its stored last-round index is polynomial-time computable; every projection and fixed repetition carries a native ledger.
\end{lemma}

\begin{definition}[Saturating cursor update]
\lean{Complexity.clA5Next}
\label{Complexity.clA5Next}
\leanok
\uses{Complexity.clA5StoredRound, Complexity.clHeaderField, Complexity.clHeaderTail, Turing.FinTM, Turing.pairEncode}
One cursor update on the persistent body word extends the unary cursor by one token while the cursor is at most the stored last-round index, and leaves it unchanged beyond that saturation point; the complete protected producer answer is untouched.
\end{definition}

\begin{lemma}[The cursor update is native]
\lean{Complexity.clA5Next_native}
\label{Complexity.clA5Next_native}
\leanok
\uses{Complexity.PolyTimeComputable, Complexity.PolyTimeComputable.comp, Complexity.clA5Le_native, Complexity.clA5Next, Complexity.clA5StoredRound_native, Complexity.clA5_pt_cond, Complexity.clA5_pt_const, Complexity.clA5_pt_unaryLength, Complexity.clHeaderField_native, Complexity.clHeaderTail_native, Complexity.clNative_append, Complexity.clNative_pair, Turing.FinTM}
\difficulty{3}
For every source machine, the saturating cursor update on persistent body words is polynomial-time computable.
\end{lemma}

\begin{lemma}[The update tracks the bounded orbit]
\lean{Complexity.clA5Next_pack}
\label{Complexity.clA5Next_pack}
\leanok
\uses{Complexity.clA5Next, Complexity.clA5Pack, Complexity.clA5Stored_exact, Complexity.clHeaderField, Complexity.clHeaderTail, Complexity.clLastRound, Complexity.clProducerHorizon, Turing.FinTM, Turing.pairDecode_pairEncode}
\difficulty{3}
For every source machine, parameters, instance $x$, and cursor $i$ at most the last round plus one, the native cursor update of the body word at $i$ is the body word at $\min(i+1, R+1)$, where $R$ is the last round; the saturated cursor is stationary.
\end{lemma}

\begin{lemma}[Exact fuel and its polynomial bound]
\lean{Complexity.clA5Fuel}
\label{Complexity.clA5Fuel}
\leanok
\uses{Complexity.PolyBound, Complexity.PolyTimeComputable, Complexity.PolyTimeComputable.comp, Complexity.clA5Pack, Complexity.clA5StoredRound_native, Complexity.clA5Stored_exact, Complexity.clLastRound, Complexity.clNative_linear, Complexity.clPackedRecords_native, Complexity.clProducerHorizon, Turing.FinTM, Turing.FinTM.computesFunInTime_lengthBits, Turing.FinTM.computesFunInTime_pairEncodeFixed, Turing.FinTM.computesInTime_iff, Turing.MultiTapeTM.output_length_le}
\difficulty{4}
For every source machine and parameters, the last-round function of the input length is polynomially bounded, and the function sending an input to the binary representation of its last-round value is polynomial-time computable; the fuel is the last round index, read from an actual unary producer followed by the native binary length machine.
\end{lemma}

\begin{lemma}[One positive emission round of the host]
\lean{Complexity.clA5Round}
\label{Complexity.clA5Round}
\leanok
\uses{Complexity.CLA5Clean, Complexity.clA5HostEntry, Complexity.clA5HostTM, Complexity.clA5Host_clean, Complexity.clA5Modules, Complexity.clA5Modules_bound, Complexity.clA5_guard_add, Turing.Action.apply, Turing.Cfg.ofWords, Turing.FinTM, Turing.FinTM.controlAction, Turing.MultiTapeTM.runFrom, Turing.MultiTapeTM.runFrom_add, Turing.MultiTapeTM.runFrom_succ_eq_step, Turing.MultiTapeTM.step, Turing.stateWord}
\difficulty{5}
Let the emission module satisfy the clean contract preserving its argument and emitting a chunk, and the update module the clean contract replacing the argument and emitting nothing. Then from the anchor over any body word $w$, the host first returns to the anchor at a positive time at most $1$ plus both budgets at $|w|$, having appended exactly the emitted chunk and replaced $w$ by its update; empty chunks and a stationary saturated cursor are covered.
\end{lemma}

\begin{lemma}[Sums of polynomial bounds]
\lean{Complexity.clA5Poly_add}
\label{Complexity.clA5Poly_add}
\leanok
\uses{Complexity.PolyBound}
\difficulty{3}
If two numerical functions are polynomially bounded, then so is their pointwise sum; no computability claim about the bound is involved.
\end{lemma}

\begin{lemma}[Monomials of polynomial bounds]
\lean{Complexity.clA5Poly_call}
\label{Complexity.clA5Poly_call}
\leanok
\uses{Complexity.PolyBound}
\difficulty{3}
If a numerical function $f$ is polynomially bounded, then for all constants $a$ and $d$ the function $n \mapsto a(f(n)+1)^d$ is polynomially bounded; a size estimate is not thereby turned into a machine clock.
\end{lemma}

\begin{lemma}[Storage bound of the body word]
\lean{Complexity.clA5Pack_bound}
\label{Complexity.clA5Pack_bound}
\leanok
\uses{Complexity.PolyBound, Complexity.clA5Fuel, Complexity.clA5Pack, Complexity.clA5Poly_add, Complexity.clA5Poly_call, Complexity.clLastRound, Complexity.clPackedRecords, Complexity.clPackedRecords_machine, Complexity.clProducerHorizon, Turing.FinTM, Turing.pairEncode}
\difficulty{4}
For every source machine and parameters, there is a polynomially bounded function $W$ such that for every instance $x$ and cursor $i$ at most the last round plus one, the persistent body word at cursor $i$ has length at most $W(|x|)$.
\end{lemma}

\begin{lemma}[The emitter reduces to a chunk selector]
\lean{Complexity.clA5Output_of_nativeChunk}
\label{Complexity.clA5Output_of_nativeChunk}
\leanok
\uses{Complexity.CLFieldCode, Complexity.PolyBound, Complexity.PolyTimeComputable, Complexity.clA5Clean_emit, Complexity.clA5Clean_install, Complexity.clA5Fuel, Complexity.clA5HostTM, Complexity.clA5Modules, Complexity.clA5Modules_bound, Complexity.clA5Next, Complexity.clA5Next_native, Complexity.clA5Next_pack, Complexity.clA5Pack, Complexity.clA5Pack_bound, Complexity.clA5Poly_add, Complexity.clA5Poly_call, Complexity.clA5Round, Complexity.clA5Startup, Complexity.clChunk, Complexity.clChunks_serialize, Complexity.clEmitter_of_body, Complexity.clLastRound, Complexity.clProducerHorizon, Complexity.clTableau, Complexity.clTableauGroups, Complexity.clTableau_chunks, Std.Sat.CNF.serialize, Turing.Cfg.ofWords, Turing.FinTM, Turing.FinTM.ComputesFunInTime, Turing.FinTM.ComputesInTime.mono, Turing.MultiTapeTM.initCfg, Turing.MultiTapeTM.runFrom, Turing.stateWord}
\difficulty{6}
Let an emission function on body words be polynomial-time computable and agree, on every body word at a legal cursor $i$, with the chunk of round $i$ of the tableau group list, and with the empty word past the last round. Then the function sending an instance $x$ to the serialization of the pure Cook--Levin tableau of $x$ at the producer horizon is polynomial-time computable: startup, cursor updates, clean first returns, saturated silence, exact fuel, and one common polynomial budget are all discharged, and no satisfiability fact enters the word identity.
\end{lemma}

\begin{definition}[Comparator with replayed verdict]
\lean{Complexity.clA5CompareTM}
\label{Complexity.clA5CompareTM}
\leanok
\uses{Complexity.clCmpTM, Complexity.clPreparedTM, Turing.FinTM}
The extended comparator machine behaves as the natively loaded four-word comparator, except that on reaching a completed verdict control it emits that single verdict bit and physically halts; the source comparator and its signed arithmetic are unchanged.
\end{definition}

\begin{lemma}[One arithmetic bit in linear time]
\lean{Complexity.clA5Compare_compute}
\label{Complexity.clA5Compare_compute}
\leanok
\uses{Complexity.clA5CompareTM, Complexity.clCmpCfg, Complexity.clCmpOrbit, Complexity.clCmpSize, Complexity.clCmpTM, Complexity.clCmpVerdict, Complexity.clCmpVerdict_spec, Complexity.clCmp_run, Complexity.clFields, Complexity.clFields_width, Complexity.clFirst, Complexity.clMap_run, Complexity.clNum, Complexity.clPreparedCfg, Complexity.clPreparedTM, Complexity.clPrepared_idle, Complexity.clPrepared_run, Complexity.clSlotCfg, Turing.Action, Turing.Action.apply, Turing.Action.mapState, Turing.Cfg, Turing.Cfg.mapState, Turing.Cfg.ofWords, Turing.FinTM.ComputesInTime, Turing.FinTM.ComputesInTime.mono, Turing.FinTM.computesInTime_iff, Turing.MultiTapeTM.initCfg, Turing.MultiTapeTM.runFrom, Turing.MultiTapeTM.runFrom_add, Turing.MultiTapeTM.runFrom_succ_eq_step', Turing.MultiTapeTM.step}
\difficulty{6}
For every four words, the verdict-replaying comparator computes, from the packed encoding of the four words, the single bit deciding whether the value of the first word plus the second equals the value of the third plus the fourth, within $35$ times the encoded argument length plus one.
\end{lemma}

\begin{lemma}[Native equality of binary values]
\lean{Complexity.clA5EqNum_native}
\label{Complexity.clA5EqNum_native}
\leanok
\uses{Complexity.PolyTimeComputable, Complexity.clA5CompareTM, Complexity.clA5Compare_compute, Complexity.clA5_pt_const, Complexity.clFields, Complexity.clNative_fields, Complexity.clNative_image, Complexity.clNum}
\difficulty{3}
If two word producers are polynomial-time computable, then the single-bit test deciding whether their outputs denote equal natural numbers is polynomial-time computable, using the banked four-word comparator with the two remaining words empty.
\end{lemma}

\begin{definition}[State of bounded binary decoding]
\lean{Complexity.clA5DecodeState}
\label{Complexity.clA5DecodeState}
\leanok
\uses{Complexity.clNum, Turing.pairEncode}
For a binary target word $w$ and a candidate count $j$, the decoding state retains $w$, the unary candidate of $j$ tokens, and the most recent successful unary answer: the unary value of $w$ when it is at most $j$, and the empty word otherwise.
\end{definition}

\begin{definition}[One bounded-decoding step]
\lean{Complexity.clA5DecodeStep}
\label{Complexity.clA5DecodeStep}
\leanok
\uses{Complexity.clHeaderField, Complexity.clHeaderTail, Complexity.clNum, Turing.pairEncode}
One decoding step yields the state whose unary candidate is extended by one token; the genuine binary encoding of the new candidate's length is compared with the retained target, and the candidate is recorded as the answer exactly on equality, the old answer being kept otherwise.
\end{definition}

\begin{lemma}[The decoding step is native]
\lean{Complexity.clA5DecodeStep_native}
\label{Complexity.clA5DecodeStep_native}
\leanok
\uses{Complexity.PolyTimeComputable, Complexity.PolyTimeComputable.comp, Complexity.clA5DecodeStep, Complexity.clA5EqNum_native, Complexity.clA5_pt_cond, Complexity.clA5_pt_const, Complexity.clHeaderField_native, Complexity.clHeaderTail_native, Complexity.clNative_append, Complexity.clNative_linear, Complexity.clNative_pair, Complexity.clNum_bits, Turing.FinTM.computesFunInTime_lengthBits}
\difficulty{3}
The bounded-decoding step on state words is polynomial-time computable; the equality test uses the banked arithmetic comparator and the native binary length routine.
\end{lemma}

\begin{lemma}[The decoding invariant advances]
\lean{Complexity.clA5DecodeStep_state}
\label{Complexity.clA5DecodeStep_state}
\leanok
\uses{Complexity.clA5DecodeState, Complexity.clA5DecodeStep, Complexity.clHeaderField, Complexity.clHeaderTail, Complexity.clNum, Turing.pairDecode_pairEncode}
\difficulty{3}
For every binary target $w$ and count $j$, one decoding step on the state at $(w, j)$ yields exactly the state at $(w, j+1)$; no decoded value is used as a loop clock, the supplied unary budget being the only repetition count.
\end{lemma}

\begin{lemma}[Initial decoding state]
\lean{Complexity.clA5DecodeState_zero}
\label{Complexity.clA5DecodeState_zero}
\leanok
\uses{Complexity.clA5DecodeState, Complexity.clNum, Turing.pairEncode}
\difficulty{2}
For every binary target $w$, the decoding state at count zero is the pair of $w$ with two empty words; the unary answer is empty also when the target denotes zero.
\end{lemma}

\begin{lemma}[The iteration follows the state word]
\lean{Complexity.clA5Decode_orbit}
\label{Complexity.clA5Decode_orbit}
\leanok
\uses{Complexity.clA5DecodeState, Complexity.clA5DecodeStep, Complexity.clA5DecodeStep_state}
\difficulty{3}
For every binary target $w$ and count $j$, iterating the decoding step $j$ times from the initial state yields exactly the decoding state at $(w, j)$.
\end{lemma}

\begin{lemma}[Linear growth of the decoding state]
\lean{Complexity.clA5Decode_size}
\label{Complexity.clA5Decode_size}
\leanok
\uses{Complexity.clA5DecodeState, Turing.pairEncode}
\difficulty{3}
For every binary target $w$ and count $j$, the decoding state has length at most $2|w| + 3j + 4$; candidate and answer storage grow only with the supplied unary clock, even for a malformed target denoting a huge value.
\end{lemma}

\begin{lemma}[Bounded binary-to-unary decoding is native]
\lean{Complexity.clA5Decode_native}
\label{Complexity.clA5Decode_native}
\leanok
\uses{Complexity.PolyTimeComputable, Complexity.PolyTimeComputable.comp, Complexity.clA5DecodeState, Complexity.clA5DecodeState_zero, Complexity.clA5DecodeStep, Complexity.clA5DecodeStep_native, Complexity.clA5Decode_orbit, Complexity.clA5Decode_size, Complexity.clA5Iter_native, Complexity.clA5_pt_const, Complexity.clHeaderTail, Complexity.clHeaderTail_native, Complexity.clNative_pair, Complexity.clNum, Turing.pairDecode_pairEncode, Turing.pairEncode}
\difficulty{4}
If a binary-word producer and a unary bound producer are polynomial-time computable, then so is the function returning the unary value of the produced binary word when that value is at most the produced bound, and the empty word otherwise; all malformed requests retain the proved polynomial cost.
\end{lemma}

\begin{lemma}[Tail passes compose additively]
\lean{Complexity.clA5Tail_add}
\label{Complexity.clA5Tail_add}
\leanok
\uses{Complexity.clHeaderTail, Turing.pairDecode}
\difficulty{3}
For all counts $a$ and $b$ and every word, following $a + b$ encoded header tails equals following $b$ tails of the result of following $a$ tails.
\end{lemma}

\begin{lemma}[Tails remove a whole encoded prefix]
\lean{Complexity.clA5Tail_fields}
\label{Complexity.clA5Tail_fields}
\leanok
\uses{Complexity.clA5Tail_add, Complexity.clFields, Complexity.clHeaderTail, Complexity.clPair_append, Turing.pairDecode_pairEncode}
\difficulty{3}
For every list of words and every trailer, following as many header tails as there are fields in the packed encoding of the list, applied to the packed encoding followed by the trailer, returns the literal trailer.
\end{lemma}

\begin{lemma}[Fields are recovered bit for bit]
\lean{Complexity.clA5Field_fields}
\label{Complexity.clA5Field_fields}
\leanok
\uses{Complexity.clA5Tail_add, Complexity.clFields, Complexity.clHeaderField, Complexity.clHeaderTail, Complexity.clPair_append, Turing.pairDecode, Turing.pairDecode_pairEncode}
\difficulty{4}
For every list of words, trailer, and legal index $i$ below the list length, the $i$-th header field of the packed encoding followed by the trailer is exactly the $i$-th word of the list, regardless of the suffix.
\end{lemma}

\begin{lemma}[Row streams contain exact field counts]
\lean{Complexity.clA5Tail_rows}
\label{Complexity.clA5Tail_rows}
\leanok
\uses{Complexity.clA5Tail_add, Complexity.clA5Tail_fields, Complexity.clHeaderTail, Complexity.clRows}
\difficulty{3}
For every row family with $l$ fields per row, count $n$, and trailer, following $n \cdot l$ header tails of the stream of the first $n$ rows followed by the trailer returns the literal trailer.
\end{lemma}

\begin{lemma}[Row-table lookup reads the stored value]
\lean{Complexity.clA5Field_rows}
\label{Complexity.clA5Field_rows}
\leanok
\uses{Complexity.clA5Field_fields, Complexity.clA5Tail_add, Complexity.clA5Tail_rows, Complexity.clHeaderField, Complexity.clRows, Complexity.clRows_split}
\difficulty{3}
For every row family with $l$ fields per row, bound $N$, time $t < N$, and coordinate $i$, the header field at index $t \cdot l + i$ of the stream of $N$ rows is exactly the field $i$ of row $t$.
\end{lemma}

\begin{lemma}[The visit table is a sequential row stream]
\lean{Complexity.clA5Visits_rows}
\label{Complexity.clA5Visits_rows}
\leanok
\uses{Complexity.clRows, Complexity.clVisitCode, Complexity.clVisitRow_exact, Complexity.clVisitRows, Complexity.prevVisit, Turing.FinTM}
\difficulty{3}
For every source machine, parameters $C$ and $e$, instance $x$, and count $N$, the ordered visit table of $N$ rows equals the generic sequential row stream whose row at time $t$ lists, per work tape, the optional-time code of the greatest strictly earlier visit over the reference length $|x| + C(|x|+1)^e$; time zero is included.
\end{lemma}

\begin{definition}[Stored unary input length]
\lean{Complexity.clA5InputSize}
\label{Complexity.clA5InputSize}
\leanok
\uses{Complexity.clA5StoredHeader, Complexity.clHeaderField}
From the persistent body word, the stored reference-input length is the length of header field four, the all-false reference input retained literally by the producer.
\end{definition}

\begin{definition}[Stored unary horizon]
\lean{Complexity.clA5Horizon}
\label{Complexity.clA5Horizon}
\leanok
\uses{Complexity.clA5StoredHeader, Complexity.clHeaderTail}
From the persistent body word, the stored source horizon is the length of the fifth header tail, the unary clock retained literally by the producer.
\end{definition}

\begin{definition}[Stored trajectory projection]
\lean{Complexity.clA5Trajectory}
\label{Complexity.clA5Trajectory}
\leanok
\uses{Complexity.clHeaderField, Complexity.clHeaderTail}
From the persistent body word, the stored chronological movement-count trajectory is reached by fixed native projections: the first field of the first tail's payload.
\end{definition}

\begin{definition}[Stored visit-table projection]
\lean{Complexity.clA5Visits}
\label{Complexity.clA5Visits}
\leanok
\uses{Complexity.clHeaderTail}
From the persistent body word, the stored inclusive visit table is the third header tail, reached by fixed native projections.
\end{definition}

\begin{lemma}[Stored projections on legal cursor states]
\lean{Complexity.clA5Data_exact}
\label{Complexity.clA5Data_exact}
\leanok
\uses{Complexity.clA5Horizon, Complexity.clA5InputSize, Complexity.clA5Pack, Complexity.clA5StoredHeader, Complexity.clA5Trajectory, Complexity.clA5Visits, Complexity.clHeaderField, Complexity.clHeaderTail, Complexity.clPackedRecords, Complexity.clPrepHeader, Complexity.clProducerHorizon, Complexity.clRecords, Complexity.clVisitRows, Turing.FinTM, Turing.pairDecode_pairEncode}
\difficulty{1}
For every source machine, parameters, instance $x$, and cursor $i$, the stored input length is $m = |x| + C(|x|+1)^e$, the stored horizon is the producer horizon $T$, the stored trajectory is the inclusive serialized trajectory of $T+1$ rows over the all-false input of length $m$, and the stored visit table is the ordered table of $T+1$ visit rows.
\end{lemma}

\begin{lemma}[Stored sizes are native unary values]
\lean{Complexity.clA5Sizes_native}
\label{Complexity.clA5Sizes_native}
\leanok
\uses{Complexity.PolyTimeComputable, Complexity.PolyTimeComputable.comp, Complexity.clA5Horizon, Complexity.clA5InputSize, Complexity.clA5StoredHeader, Complexity.clA5_pt_unaryLength, Complexity.clHeaderField_native, Complexity.clHeaderTail_native}
\difficulty{2}
The functions sending a persistent body word to the unary words of its stored input length and of its stored horizon are both polynomial-time computable.
\end{lemma}

\begin{definition}[Bounded numeric interpretation]
\lean{Complexity.clA5SmallNum}
\label{Complexity.clA5SmallNum}
\leanok
\uses{Complexity.clNum}
For a binary word $w$ and a bound $N$, the bounded value is the value of $w$ when it is at most $N$ and zero otherwise, so the interpretation is total also on malformed requests.
\end{definition}

\begin{lemma}[Bounded interpretation is native]
\lean{Complexity.clA5SmallNum_native}
\label{Complexity.clA5SmallNum_native}
\leanok
\uses{Complexity.PolyTimeComputable, Complexity.clA5Decode_native, Complexity.clA5SmallNum, Complexity.clNum}
\difficulty{2}
If a binary-word producer and a unary bound producer are polynomial-time computable, then the unary word of the bounded value of the produced word is polynomial-time computable.
\end{lemma}

\begin{definition}[Stored movement count at a time]
\lean{Complexity.clA5Count}
\label{Complexity.clA5Count}
\leanok
\uses{Complexity.clA5Horizon, Complexity.clA5SmallNum, Complexity.clA5Trajectory, Complexity.clHeaderField, Complexity.clRecFields, Turing.FinTM}
For a counter coordinate $i$, a persistent body word, and a time $t$, the stored count is the bounded value, under the stored horizon, of the trajectory field at index $t$ times the field count plus $i$, read from the actual chronological trajectory.
\end{definition}

\begin{lemma}[Stored counts are native]
\lean{Complexity.clA5Count_native}
\label{Complexity.clA5Count_native}
\leanok
\uses{Complexity.PolyTimeComputable, Complexity.PolyTimeComputable.comp, Complexity.clA5Count, Complexity.clA5Field_native, Complexity.clA5Sizes_native, Complexity.clA5SmallNum_native, Complexity.clA5Times_native, Complexity.clA5_pt_const, Complexity.clHeaderField_native, Complexity.clHeaderTail_native, Complexity.clNative_append, Complexity.clRecFields, Turing.FinTM}
\difficulty{3}
For every counter coordinate and polynomial-time unary time producer, the unary word of the stored movement count at the produced time is polynomial-time computable; projections, clock construction, repeated field scans, and decoding calls are all charged.
\end{lemma}

\begin{lemma}[Stored counts are the banked counts]
\lean{Complexity.clA5Count_exact}
\label{Complexity.clA5Count_exact}
\leanok
\uses{Complexity.clA5Count, Complexity.clA5Data_exact, Complexity.clA5Field_rows, Complexity.clA5Pack, Complexity.clA5SmallNum, Complexity.clCounts, Complexity.clCounts_bound, Complexity.clNum_bits, Complexity.clProducerHorizon, Complexity.clRecFields, Complexity.clRows_records, Turing.FinTM}
\difficulty{3}
For every source machine, parameters, instance $x$, cursor $j$, time $t$ at most the producer horizon, and coordinate $i$, the stored count read from the body word equals the canonical movement count at time $t$ over the all-false reference input; the bound never truncates valid data.
\end{lemma}

\begin{lemma}[Unary addition is concatenation]
\lean{Complexity.clA5Add_native}
\label{Complexity.clA5Add_native}
\leanok
\uses{Complexity.PolyTimeComputable, Complexity.clNative_append}
\difficulty{1}
If two unary values are polynomial-time computable, then the unary word of their sum is polynomial-time computable by literal concatenation.
\end{lemma}

\begin{lemma}[Unary subtraction is charged tailing]
\lean{Complexity.clA5Sub_native}
\label{Complexity.clA5Sub_native}
\leanok
\uses{Complexity.PolyTimeComputable, Complexity.clA5Drop_native}
\difficulty{1}
If two unary values are polynomial-time computable, then the unary word of their truncated difference is polynomial-time computable by the charged repeated-tail operation.
\end{lemma}

\begin{definition}[Stored shifted input position]
\lean{Complexity.clA5InputPos}
\label{Complexity.clA5InputPos}
\leanok
\uses{Complexity.clA5Count, Turing.FinTM}
For a persistent body word and a time $t$, the stored input position is the positive stored input count plus one, minus the negative stored input count, with natural subtraction keeping the left boundary at zero.
\end{definition}

\begin{lemma}[The stored input position is native]
\lean{Complexity.clA5InputPos_native}
\label{Complexity.clA5InputPos_native}
\leanok
\uses{Complexity.PolyTimeComputable, Complexity.clA5Add_native, Complexity.clA5Count_native, Complexity.clA5InputPos, Complexity.clA5Sub_native, Complexity.clA5_pt_const, Turing.FinTM}
\difficulty{1}
For every polynomial-time unary time producer, the unary word of the stored shifted input position at the produced time is polynomial-time computable.
\end{lemma}

\begin{lemma}[Stored positions match the schedule]
\lean{Complexity.clA5InputPos_exact}
\label{Complexity.clA5InputPos_exact}
\leanok
\uses{Complexity.clA5Count_exact, Complexity.clA5InputPos, Complexity.clA5Pack, Complexity.clCounts_schedule, Complexity.clProducerHorizon, Complexity.clSigned, Complexity.inputPosAt, Turing.FinTM}
\difficulty{3}
For every source machine, parameters, instance $x$, cursor $j$, and time $t$ at most the producer horizon, the stored shifted input position equals the public oblivious input position at time $t$ over the reference length $|x| + C(|x|+1)^e$.
\end{lemma}

\begin{definition}[Stored visit code at a time and tape]
\lean{Complexity.clA5Visit}
\label{Complexity.clA5Visit}
\leanok
\uses{Complexity.clA5Visits, Complexity.clHeaderField, Turing.FinTM}
For a work tape $\tau$, a persistent body word, and a time $t$, the stored visit code is the field at index $t$ times the tape count plus $\tau$ of the stored visit table, including its explicit presence marker.
\end{definition}

\begin{lemma}[Stored visit codes are native]
\lean{Complexity.clA5Visit_native}
\label{Complexity.clA5Visit_native}
\leanok
\uses{Complexity.PolyTimeComputable, Complexity.clA5Add_native, Complexity.clA5Field_native, Complexity.clA5Times_native, Complexity.clA5Visit, Complexity.clA5_pt_const, Complexity.clHeaderTail_native, Turing.FinTM}
\difficulty{2}
For every work tape and polynomial-time unary time producer, the stored visit code at the produced time is polynomial-time computable, read by the charged sequential field reader.
\end{lemma}

\begin{lemma}[Stored visit codes are the public schedule]
\lean{Complexity.clA5Visit_exact}
\label{Complexity.clA5Visit_exact}
\leanok
\uses{Complexity.clA5Data_exact, Complexity.clA5Field_rows, Complexity.clA5Pack, Complexity.clA5Visit, Complexity.clA5Visits_rows, Complexity.clProducerHorizon, Complexity.clVisitCode, Complexity.prevVisit, Turing.FinTM}
\difficulty{2}
For every source machine, parameters, instance $x$, cursor $j$, time $t$ at most the producer horizon, and tape $\tau$, the stored visit code equals the optional-time code of the public greatest strictly earlier visit; the reference machine is never rerun and the protected table is never replaced.
\end{lemma}

\begin{definition}[Stored-coordinate input fragment]
\lean{Complexity.clA5InputFragment}
\label{Complexity.clA5InputFragment}
\leanok
\uses{Complexity.CLFieldCode, Complexity.clA5InputPos, Complexity.clA5InputSize, Complexity.clFragment, Complexity.clGroup, Turing.FinTM}
For a persistent body word and a time $t$, the input fragment serializes the present-input table wired to input variable $p - 1$ when the stored input position $p$ is interior, with $0 < p$ at most the stored input length, and the blank-input table with a dummy index otherwise.
\end{definition}

\begin{lemma}[Input fragments are native]
\lean{Complexity.clA5InputFragment_native}
\label{Complexity.clA5InputFragment_native}
\leanok
\uses{Complexity.CLFieldCode, Complexity.PolyTimeComputable, Complexity.clA5Group_native, Complexity.clA5InputFragment, Complexity.clA5InputPos_native, Complexity.clA5Le_native, Complexity.clA5Sizes_native, Complexity.clA5Sub_native, Complexity.clA5_pt_and, Complexity.clA5_pt_cond, Complexity.clA5_pt_const, Turing.FinTM}
\difficulty{3}
For every source machine, state code, and polynomial-time unary time producer, the serialized input fragment at the produced time is polynomial-time computable; both template choices are finite protected tables, and the branch predicate and shifted variable address are charged.
\end{lemma}

\begin{lemma}[Input fragments match the pure groups]
\lean{Complexity.clA5InputFragment_exact}
\label{Complexity.clA5InputFragment_exact}
\leanok
\uses{Complexity.CLFieldCode, Complexity.clA5Data_exact, Complexity.clA5InputFragment, Complexity.clA5InputPos_exact, Complexity.clA5Pack, Complexity.clFragment, Complexity.clInputGroup, Complexity.clProducerHorizon, Turing.FinTM}
\difficulty{2}
For every source machine, state code, parameters, instance $x$, cursor $j$, and time $t$ at most the producer horizon, the stored-coordinate input fragment equals the serialized fragment of the pure input clause group at time $t$ over the reference length.
\end{lemma}

\begin{definition}[Stored-visit work fragment]
\lean{Complexity.clA5WorkFragment}
\label{Complexity.clA5WorkFragment}
\leanok
\uses{Complexity.CLFieldCode, Complexity.clA5Horizon, Complexity.clA5InputSize, Complexity.clA5SmallNum, Complexity.clA5Visit, Complexity.clFragment, Complexity.clGroup, Complexity.clHeaderTail, Turing.FinTM}
For a work tape $\tau$, a persistent body word, and a time $t$: an empty stored visit code selects the blank first-visit table; a nonempty code selects the earlier-visit table wired from the source time decoded, bounded by the stored horizon, from the code's payload. A present visit at time zero has its distinct nonempty marker.
\end{definition}

\begin{lemma}[Work fragments are native]
\lean{Complexity.clA5WorkFragment_native}
\label{Complexity.clA5WorkFragment_native}
\leanok
\uses{Complexity.CLFieldCode, Complexity.PolyTimeComputable, Complexity.PolyTimeComputable.comp, Complexity.clA5Group_native, Complexity.clA5Sizes_native, Complexity.clA5SmallNum_native, Complexity.clA5Visit_native, Complexity.clA5WorkFragment, Complexity.clA5_pt_cond, Complexity.clA5_pt_const, Complexity.clA5_pt_eq, Complexity.clHeaderTail_native, Turing.FinTM}
\difficulty{3}
For every source machine, state code, work tape, and polynomial-time unary time producer, the serialized work fragment at the produced time is polynomial-time computable, with the greatest-earlier source address read from the protected visit table.
\end{lemma}

\begin{lemma}[Work fragments match the pure groups]
\lean{Complexity.clA5WorkFragment_exact}
\label{Complexity.clA5WorkFragment_exact}
\leanok
\uses{Complexity.CLFieldCode, Complexity.clA5Data_exact, Complexity.clA5Pack, Complexity.clA5SmallNum, Complexity.clA5Visit_exact, Complexity.clA5WorkFragment, Complexity.clFragment, Complexity.clHeaderTail, Complexity.clNum_bits, Complexity.clPrev_spec, Complexity.clProducerHorizon, Complexity.clVisitCode, Complexity.clWorkGroup, Complexity.prevVisit, Turing.FinTM, Turing.pairDecode_pairEncode, Turing.pairEncode}
\difficulty{4}
For every source machine, state code, parameters, instance $x$, cursor $j$, time $t$ at most the producer horizon, and tape $\tau$, the stored-visit work fragment equals the serialized fragment of the pure work clause group at time $t$ and tape $\tau$ over the reference length; the stored earlier source is strictly below the target, so bounded decoding preserves every legitimate visit, including source zero.
\end{lemma}

\begin{lemma}[Unary division by a fixed constant]
\lean{Complexity.clA5Div_native}
\label{Complexity.clA5Div_native}
\leanok
\uses{Complexity.PolyTimeComputable, Complexity.PolyTimeComputable.comp, Complexity.clA5MapWord, Complexity.clA5Map_poly, Complexity.clA5_pt_const}
\difficulty{4}
If a unary value is polynomial-time computable, then for every fixed constant $k$ the unary word of its quotient by $k$ is polynomial-time computable: a finite-control counter emits one bit per completed block, and division by zero yields the empty word.
\end{lemma}

\begin{lemma}[Unary remainder by a fixed constant]
\lean{Complexity.clA5Mod_native}
\label{Complexity.clA5Mod_native}
\leanok
\uses{Complexity.PolyTimeComputable, Complexity.clA5Div_native, Complexity.clA5Sub_native, Complexity.clA5Times_native}
\difficulty{2}
If a unary value is polynomial-time computable, then for every fixed constant the unary word of its remainder is polynomial-time computable, by native subtraction after the finite-control quotient scan; the zero-constant case is included.
\end{lemma}

\begin{lemma}[Finite native selection]
\lean{Complexity.clA5Select_native}
\label{Complexity.clA5Select_native}
\leanok
\uses{Complexity.PolyTimeComputable, Complexity.clA5Sub_native, Complexity.clA5_pt_cond, Complexity.clA5_pt_const, Complexity.clA5_pt_eq}
\difficulty{4}
For a finite list of polynomial-time computable word functions and a polynomial-time unary selector, the function applying the selected list entry to the input, with the empty-word constant as out-of-range default, is polynomial-time computable; no unbounded transition table is consulted.
\end{lemma}

\begin{lemma}[Indexing the time-then-tape work family]
\lean{Complexity.clA5Work_get}
\label{Complexity.clA5Work_get}
\leanok
\difficulty{5}
For every element family on times and tapes, tape count $k$, bound $N$, and index $j < kN$, the $j$-th entry of the list enumerating, for each time below $N$, the $k$ tape entries in order, is the entry at time $j$ divided by $k$ and tape $j$ modulo $k$; the zero-tape case has no legal index.
\end{lemma}

\begin{definition}[Selector in the six-family order]
\lean{Complexity.clA5GroupAt}
\label{Complexity.clA5GroupAt}
\leanok
Given parameters $n$, $k$, $T$ and the six clause-family generators, the selector at ordinal $i$ subtracts the family offsets in order: the $n$ pin groups, the single initial group, $T$ state groups, $T+1$ input groups, $k(T+1)$ work groups in time-then-tape order via quotient and remainder, and $T$ accept groups, with the empty formula past the end.
\end{definition}

\begin{lemma}[The selector is list lookup]
\lean{Complexity.clA5GroupAt_exact}
\label{Complexity.clA5GroupAt_exact}
\leanok
\uses{Complexity.clA5GroupAt, Complexity.clA5Work_get, Complexity.clGroups}
\difficulty{5}
For all parameters and family generators and every ordinal $i$, the six-family selector at $i$ equals the lookup of the frozen ordered group list at $i$ with the empty default, including every family boundary.
\end{lemma}

\begin{lemma}[Branching on strict unary comparison]
\lean{Complexity.clA5IfLt_native}
\label{Complexity.clA5IfLt_native}
\leanok
\uses{Complexity.PolyTimeComputable, Complexity.clA5Add_native, Complexity.clA5Le_native, Complexity.clA5_pt_cond, Complexity.clA5_pt_const}
\difficulty{2}
If two unary values and two word functions are polynomial-time computable, then so is the function returning the first word when the first value is strictly smaller and the second word otherwise.
\end{lemma}

\begin{lemma}[Branching on a zero test]
\lean{Complexity.clA5IfZero_native}
\label{Complexity.clA5IfZero_native}
\leanok
\uses{Complexity.PolyTimeComputable, Complexity.clA5_pt_cond, Complexity.clA5_pt_eq}
\difficulty{2}
If a unary value and two word functions are polynomial-time computable, then so is the function returning the first word when the value is zero and the second otherwise.
\end{lemma}

\begin{lemma}[Native indexed bit read]
\lean{Complexity.clA5GetBit_native}
\label{Complexity.clA5GetBit_native}
\leanok
\uses{Complexity.PolyTimeComputable, Complexity.PolyTimeComputable.comp, Complexity.clA5Drop_native, Complexity.clA5_pt_eq, Complexity.clA5_pt_head}
\difficulty{3}
If a word producer and a unary index producer are polynomial-time computable, then the single-bit read of the produced word at the produced index, with false as out-of-range default, is polynomial-time computable.
\end{lemma}

\begin{lemma}[Native pinning units]
\lean{Complexity.clA5Pin_native}
\label{Complexity.clA5Pin_native}
\leanok
\uses{Complexity.PolyTimeComputable, Complexity.clA5GetBit_native, Complexity.clA5Template_native, Complexity.clA5_pt_cond, Complexity.clFragment}
\difficulty{3}
If a word producer and a unary index producer are polynomial-time computable, then the serialized fragment of the unit clause pinning variable $i$ to the word's bit at $i$, with the false default out of range, is polynomial-time computable.
\end{lemma}

\begin{definition}[Current emission cursor]
\lean{Complexity.clA5Cursor}
\label{Complexity.clA5Cursor}
\leanok
\uses{Complexity.clHeaderField}
From the persistent body word, the current bounded emission cursor is the length of the unary word stored before the immutable producer answer.
\end{definition}

\begin{lemma}[Cursor and instance are native]
\lean{Complexity.clA5Cursor_native}
\label{Complexity.clA5Cursor_native}
\leanok
\uses{Complexity.PolyTimeComputable, Complexity.PolyTimeComputable.comp, Complexity.clA5Cursor, Complexity.clA5Instance, Complexity.clA5_pt_unaryLength, Complexity.clHeaderField_native, Complexity.clHeaderTail_native}
\difficulty{1}
The unary current cursor and the retained original instance are both polynomial-time computable projections of the persistent body word.
\end{lemma}

\begin{definition}[Residual ordinal after the pin family]
\lean{Complexity.clA5InitialIndex}
\label{Complexity.clA5InitialIndex}
\leanok
\uses{Complexity.clA5Cursor, Complexity.clA5Instance}
From the persistent body word, the residual ordinal after the pin family is the current cursor minus the retained instance's length.
\end{definition}

\begin{definition}[Residual ordinal for the state family]
\lean{Complexity.clA5StateIndex}
\label{Complexity.clA5StateIndex}
\leanok
\uses{Complexity.clA5InitialIndex}
From the persistent body word, the residual ordinal for the state family is the post-pin residual minus one, skipping the singleton initial family.
\end{definition}

\begin{definition}[Residual ordinal for the input family]
\lean{Complexity.clA5InputIndex}
\label{Complexity.clA5InputIndex}
\leanok
\uses{Complexity.clA5Horizon, Complexity.clA5StateIndex}
From the persistent body word, the residual ordinal for the inclusive input family is the state residual minus the stored horizon.
\end{definition}

\begin{definition}[Residual ordinal for the work family]
\lean{Complexity.clA5WorkIndex}
\label{Complexity.clA5WorkIndex}
\leanok
\uses{Complexity.clA5Horizon, Complexity.clA5InputIndex}
From the persistent body word, the residual ordinal for the time-then-tape work family is the input residual minus the stored horizon plus one.
\end{definition}

\begin{definition}[Residual ordinal for the accept family]
\lean{Complexity.clA5AcceptIndex}
\label{Complexity.clA5AcceptIndex}
\leanok
\uses{Complexity.clA5Horizon, Complexity.clA5WorkIndex, Turing.FinTM}
From the persistent body word, the residual ordinal for the final accept family is the work residual minus the tape count times the stored horizon plus one.
\end{definition}

\begin{lemma}[Family offsets are native]
\lean{Complexity.clA5Indices_native}
\label{Complexity.clA5Indices_native}
\leanok
\uses{Complexity.PolyTimeComputable, Complexity.PolyTimeComputable.comp, Complexity.clA5AcceptIndex, Complexity.clA5Add_native, Complexity.clA5Cursor_native, Complexity.clA5InitialIndex, Complexity.clA5InputIndex, Complexity.clA5Sizes_native, Complexity.clA5StateIndex, Complexity.clA5Sub_native, Complexity.clA5Times_native, Complexity.clA5WorkIndex, Complexity.clA5_pt_const, Complexity.clA5_pt_unaryLength, Turing.FinTM}
\difficulty{3}
For every source machine, the unary residual ordinals of all five families, after the pins, for the state family, the input family, the work family, and the accept family, are polynomial-time computable by native unary arithmetic on stored words.
\end{lemma}

\begin{definition}[Tape choice at the work ordinal]
\lean{Complexity.clA5WorkChoice}
\label{Complexity.clA5WorkChoice}
\leanok
\uses{Complexity.CLFieldCode, Complexity.clA5WorkFragment, Complexity.clA5WorkIndex, Turing.FinTM}
At the current work-family ordinal, the choice selects the work fragment of the tape given by the ordinal modulo the tape count, at the time given by the ordinal divided by the tape count, and the empty word when the tape count is zero.
\end{definition}

\begin{lemma}[The tape choice is native]
\lean{Complexity.clA5WorkChoice_native}
\label{Complexity.clA5WorkChoice_native}
\leanok
\uses{Complexity.CLFieldCode, Complexity.PolyTimeComputable, Complexity.clA5Div_native, Complexity.clA5Indices_native, Complexity.clA5Mod_native, Complexity.clA5Select_native, Complexity.clA5WorkChoice, Complexity.clA5WorkFragment, Complexity.clA5WorkFragment_native, Complexity.clA5WorkIndex, Turing.FinTM}
\difficulty{4}
For every source machine and state code, the work-family tape choice is polynomial-time computable: finite selection over the per-tape fragment functions, driven by the native unary quotient and remainder of the work ordinal, including the empty zero-tape family.
\end{lemma}

\begin{definition}[Ordered fragment selector]
\lean{Complexity.clA5Fragment}
\label{Complexity.clA5Fragment}
\leanok
\uses{Complexity.CLFieldCode, Complexity.clA5AcceptIndex, Complexity.clA5Cursor, Complexity.clA5Horizon, Complexity.clA5InitialIndex, Complexity.clA5InputFragment, Complexity.clA5InputIndex, Complexity.clA5InputSize, Complexity.clA5Instance, Complexity.clA5StateIndex, Complexity.clA5WorkChoice, Complexity.clA5WorkIndex, Complexity.clFragment, Complexity.clGroup, Turing.FinTM}
On a persistent body word, the selector emits, by successive guarded offsets of the current cursor: a pinning unit for cursors inside the retained instance; the initial template at the post-pin ordinal zero; a state template; a stored-coordinate input fragment; a stored-visit work fragment; an acceptance template; and the empty word past all families. Only the current group's fragment is emitted; the formula terminator belongs to the separate last-round check.
\end{definition}

\begin{lemma}[The fragment selector is native]
\lean{Complexity.clA5Fragment_native}
\label{Complexity.clA5Fragment_native}
\leanok
\uses{Complexity.CLFieldCode, Complexity.PolyTimeComputable, Complexity.PolyTimeComputable.comp, Complexity.clA5Add_native, Complexity.clA5Cursor_native, Complexity.clA5Fragment, Complexity.clA5Group_native, Complexity.clA5IfLt_native, Complexity.clA5IfZero_native, Complexity.clA5Indices_native, Complexity.clA5InputFragment_native, Complexity.clA5InputSize, Complexity.clA5Pin_native, Complexity.clA5Sizes_native, Complexity.clA5Times_native, Complexity.clA5WorkChoice_native, Complexity.clA5_pt_const, Complexity.clA5_pt_unaryLength, Complexity.clFragment, Complexity.clGroup, Turing.FinTM}
\difficulty{4}
For every source machine and state code, the complete ordered fragment selector is one total polynomial-time transduction: every numeric operation is unary bounded, stored binary counters are decoded only within the explicit unary horizon, and templates and tape choices are fixed finite data of the machine.
\end{lemma}

\begin{lemma}[The selector matches the pure tableau]
\lean{Complexity.clA5Fragment_exact}
\label{Complexity.clA5Fragment_exact}
\leanok
\uses{Complexity.CLFieldCode, Complexity.clA5AcceptIndex, Complexity.clA5Cursor, Complexity.clA5Data_exact, Complexity.clA5Fragment, Complexity.clA5GroupAt, Complexity.clA5GroupAt_exact, Complexity.clA5Horizon, Complexity.clA5InitialIndex, Complexity.clA5InputFragment_exact, Complexity.clA5InputIndex, Complexity.clA5InputSize, Complexity.clA5Instance, Complexity.clA5Pack, Complexity.clA5StateIndex, Complexity.clA5Stored_exact, Complexity.clA5WorkChoice, Complexity.clA5WorkFragment_exact, Complexity.clA5WorkIndex, Complexity.clFragment, Complexity.clGroup, Complexity.clHeaderField, Complexity.clHeaderTail, Complexity.clInputGroup, Complexity.clProducerHorizon, Complexity.clTableauGroups, Complexity.clWorkGroup, Turing.FinTM, Turing.pairDecode_pairEncode}
\difficulty{5}
For every source machine, state code, parameters, instance $x$, and cursor $j$, the native fragment selector on the body word at $j$ agrees word for word with the serialized fragment of the $j$-th member of the unchanged pure tableau group list at the producer horizon; the guarded offsets are precisely list lookup, and a work ordinal's quotient and remainder name its time and tape in the original order.
\end{lemma}

\begin{lemma}[Native equality of unary values]
\lean{Complexity.clA5Eq_native}
\label{Complexity.clA5Eq_native}
\leanok
\uses{Complexity.PolyTimeComputable, Complexity.clA5Le_native, Complexity.clA5_pt_and}
\difficulty{2}
If two unary values are polynomial-time computable, then the single-bit test deciding their equality is polynomial-time computable, as the conjunction of the two unary comparisons.
\end{lemma}

\begin{definition}[Exact chunk selector]
\lean{Complexity.clA5Emit}
\label{Complexity.clA5Emit}
\leanok
\uses{Complexity.CLFieldCode, Complexity.clA5Cursor, Complexity.clA5Fragment, Complexity.clA5StoredRound, Turing.FinTM}
On a persistent body word, the emitter outputs the current fragment while the cursor is at most the stored last round, appending the single formula-terminator bit exactly at the last round, and outputs nothing in the saturated state; empty group fragments still have their own positive round.
\end{definition}

\begin{lemma}[The chunk selector is native]
\lean{Complexity.clA5Emit_native}
\label{Complexity.clA5Emit_native}
\leanok
\uses{Complexity.CLFieldCode, Complexity.PolyTimeComputable, Complexity.clA5Cursor_native, Complexity.clA5Emit, Complexity.clA5Eq_native, Complexity.clA5Fragment_native, Complexity.clA5Le_native, Complexity.clA5StoredRound_native, Complexity.clA5_pt_cond, Complexity.clA5_pt_const, Complexity.clNative_append, Turing.FinTM}
\difficulty{2}
For every source machine and state code, the complete chunk selector on body words is polynomial-time computable.
\end{lemma}

\begin{lemma}[The chunk selector is exact]
\lean{Complexity.clA5Emit_exact}
\label{Complexity.clA5Emit_exact}
\leanok
\uses{Complexity.CLFieldCode, Complexity.clA5Cursor, Complexity.clA5Emit, Complexity.clA5Fragment_exact, Complexity.clA5Pack, Complexity.clA5Stored_exact, Complexity.clChunk, Complexity.clHeaderField, Complexity.clHeaderTail, Complexity.clLastRound, Complexity.clProducerHorizon, Complexity.clTableauGroups, Turing.FinTM, Turing.pairDecode_pairEncode}
\difficulty{2}
For every source machine, state code, parameters, instance $x$, and cursor $j$, the chunk selector on the body word at $j$ equals the chunk of round $j$ of the pure tableau group list when $j$ is at most the last round, and the empty word otherwise; member counts and the sole-last-terminator convention are preserved.
\end{lemma}

\begin{lemma}[Physical output identity]
\lean{Complexity.clA5OutputIdentity}
\label{Complexity.clA5OutputIdentity}
\leanok
\uses{Complexity.CLFieldCode, Complexity.PolyTimeComputable, Complexity.clA5Emit, Complexity.clA5Emit_exact, Complexity.clA5Emit_native, Complexity.clA5Output_of_nativeChunk, Complexity.clProducerHorizon, Complexity.clTableau, Std.Sat.CNF.serialize, Turing.FinTM}
\difficulty{1}
For every source machine, state code, and parameters, the function sending an instance $x$ to the serialization of the pure Cook--Levin tableau of $x$ at the producer horizon is computable by one finite machine in polynomial time from genuine native input. No certificate or satisfiability assumption is involved.
\end{lemma}

\begin{definition}[Raw block of an assignment]
\lean{Complexity.clA5Block}
\label{Complexity.clA5Block}
\leanok
\uses{Complexity.CLFieldCode, Complexity.clPack, Complexity.clWidth, Turing.FinTM}
For an input-variable count $m$, a total assignment $a$, and a time $t$, the raw snapshot block at $t$ reads the assignment bits at the packed variables of the $t$-th block of width equal to the snapshot block width.
\end{definition}

\begin{definition}[Bit-level meaning of one template]
\lean{Complexity.clA5Meaning}
\label{Complexity.clA5Meaning}
\leanok
\uses{Complexity.CLFieldCode, Complexity.CLTemplateKind, Complexity.clA5Block, Complexity.clBlockDecode, Complexity.clBlockEncode, Complexity.clInputSlice, Complexity.clStateSlice, Complexity.clSymbolCode, Complexity.clWorkSlice, Complexity.emitted, Complexity.stepState, Complexity.writtenOrKept, Turing.FinTM}
For a template kind, wiring data $m$, source time $s$, target time $t$, selected index $p$, and an assignment $a$, the meaning states: for the initial kind, that the raw target block equals the encoding of the start snapshot reading bit $a(p)$ or a blank; for the state kind, that the target state slice encodes the stepped state of the decoded raw source block; for the input kind, that the target input slice encodes the selected bit or a blank; for a work kind, that the target work slice encodes the symbol written or kept by the decoded source at its last earlier visit, or a blank; and for the accept kind, that the decoded source does not emit false.
\end{definition}

\begin{lemma}[Template satisfaction is bitwise meaning]
\lean{Complexity.clA5Group_eval}
\label{Complexity.clA5Group_eval}
\leanok
\uses{Complexity.CLFieldCode, Complexity.CLTemplateKind, Complexity.clA5Block, Complexity.clA5Meaning, Complexity.clBlockDecode, Complexity.clBlockEncode, Complexity.clGroup, Complexity.clInputSlice, Complexity.clPredicate, Complexity.clSourceBits, Complexity.clStateSlice, Complexity.clSymbolCode, Complexity.clTargetBits, Complexity.clTemplate_eval, Complexity.clWidth, Complexity.clWindowBit, Complexity.clWire, Complexity.clWorkSlice, Complexity.emitted, Complexity.stepState, Complexity.writtenOrKept, Turing.FinTM}
\difficulty{5}
For every source machine, state code, template kind, wiring data, and assignment $a$, the wired clause group evaluates to true under $a$ exactly when the bit-level meaning of the template holds: the wiring exposes the raw source and target blocks and the selected input bit exactly, so satisfaction means bitwise equations, never merely equality of decoded snapshots.
\end{lemma}

\begin{lemma}[Flattening conjoins group evaluations]
\lean{Complexity.clA5Flatten_eval}
\label{Complexity.clA5Flatten_eval}
\leanok
\difficulty{3}
For every list of clause groups and every assignment, the evaluation of the flattened CNF formula is true exactly when every group in the list evaluates to true.
\end{lemma}

\begin{lemma}[The six constraints of the tableau]
\lean{Complexity.clA5Tableau_eval}
\label{Complexity.clA5Tableau_eval}
\leanok
\uses{Complexity.CLFieldCode, Complexity.clA5Flatten_eval, Complexity.clA5Group_eval, Complexity.clA5Meaning, Complexity.clGroups, Complexity.clInputGroup, Complexity.clTableau, Complexity.clTableauGroups, Complexity.clWorkGroup, Turing.FinTM}
\difficulty{4}
For every source machine, state code, parameters $C$ and $e$, instance $x$, horizon $T$, and assignment $a$, with $m = |x| + C(|x|+1)^e$: the tableau evaluates to true under $a$ exactly when the first $|x|$ variables are pinned to the bits of $x$; the initial meaning holds; the state meaning holds at every $t < T$; every inclusive input group with $t \le T$ evaluates to true; every inclusive work group with $t \le T$ evaluates to true on every tape; and the acceptance meaning holds at every $t < T$.
\end{lemma}

\begin{lemma}[Boundary-aware input bits]
\lean{Complexity.clA5InputBit}
\label{Complexity.clA5InputBit}
\leanok
\uses{Complexity.inputBitAt}
\difficulty{3}
For every word $y$ and position $p$, the input symbol at $p$ is the bit of $y$ at $p - 1$ when $0 < p \le |y|$ and a blank otherwise.
\end{lemma}

\begin{lemma}[Input clauses specify the input slice]
\lean{Complexity.clA5Input_meaning}
\label{Complexity.clA5Input_meaning}
\leanok
\uses{Complexity.CLFieldCode, Complexity.clA5Block, Complexity.clA5Group_eval, Complexity.clA5InputBit, Complexity.clInputGroup, Complexity.clInputSlice, Complexity.clSymbolCode, Complexity.inputBitAt, Complexity.inputPosAt, Turing.FinTM}
\difficulty{3}
Let $y$ be a word of length $m$ and let the assignment $a$ agree with $y$ on the first $m$ variables. Then the input clause group at time $t$ evaluates to true under $a$ exactly when the input slice of the raw block at $t$ encodes the input symbol of $y$ at the oblivious input position of time $t$.
\end{lemma}

\begin{lemma}[Initial clauses pin the initial snapshot]
\lean{Complexity.clA5Initial_meaning}
\label{Complexity.clA5Initial_meaning}
\leanok
\uses{Complexity.CLFieldCode, Complexity.clA5Block, Complexity.clA5InputBit, Complexity.clA5Meaning, Complexity.clBlockEncode, Complexity.snapshotAt, Complexity.snapshotAt_zero, Turing.FinTM}
\difficulty{3}
Let $y$ be a word of length $m$ and let the assignment agree with $y$ on the first $m$ variables. Then the initial meaning holds exactly when the raw block at time zero equals the encoding of the machine's genuine initial snapshot on $y$, with the empty-input case included.
\end{lemma}

\begin{lemma}[Work clauses specify the work slice]
\lean{Complexity.clA5Work_meaning}
\label{Complexity.clA5Work_meaning}
\leanok
\uses{Complexity.CLFieldCode, Complexity.clA5Block, Complexity.clA5Group_eval, Complexity.clBlockDecode, Complexity.clSymbolCode, Complexity.clWorkGroup, Complexity.clWorkSlice, Complexity.prevVisit, Complexity.writtenOrKept, Turing.FinTM}
\difficulty{2}
For every input count $m$, time $t$, tape $\tau$, and assignment $a$, the work clause group evaluates to true under $a$ exactly when the work slice of the raw block at $t$ is the encoding of a blank on a first visit, and otherwise of the symbol written or kept, at the greatest strictly earlier visit $s$, by the snapshot decoded from the raw block at $s$.
\end{lemma}

\begin{lemma}[Bitwise reconstruction of all blocks]
\lean{Complexity.clA5Reconstruct}
\label{Complexity.clA5Reconstruct}
\leanok
\uses{Complexity.CLFieldCode, Complexity.clA5Block, Complexity.clA5Meaning, Complexity.clBlockDecode, Complexity.clBlockDecode_encode, Complexity.clBlockEncode, Complexity.clBlock_ext, Complexity.clBlock_input, Complexity.clBlock_state, Complexity.clBlock_work, Complexity.clInputSlice, Complexity.clPrev_spec, Complexity.clStateSlice, Complexity.clSymbolCode, Complexity.clWorkSlice, Complexity.inputBitAt, Complexity.inputPosAt, Complexity.prevVisit, Complexity.snapshotAt, Complexity.snapshotAt_inputSymbol, Complexity.snapshotAt_state_succ, Complexity.snapshotAt_workSymbol, Complexity.stepState, Complexity.writtenOrKept, Turing.FinTM, Turing.FinTM.Oblivious}
\difficulty{5}
Let the machine be oblivious, let $y$ have length $m$, and suppose an assignment's raw blocks satisfy: the block at time zero is the encoded initial snapshot; the state meaning at every $t < T$; the input-slice equations at every $t \le T$; and the work-slice equations at every $t \le T$ on all tapes. Then for every $t \le T$ the raw block at $t$ equals the encoding of the genuine snapshot of the run on $y$ at time $t$; junk blocks are excluded bit for bit.
\end{lemma}

\begin{definition}[Reading the assigned input]
\lean{Complexity.clA5ReadInput}
\label{Complexity.clA5ReadInput}
\leanok
For an input-variable count $m$ and an assignment $a$, the recovered verifier input is the word of the first $m$ assigned bits in order.
\end{definition}

\begin{lemma}[Recovered bits are assigned bits]
\lean{Complexity.clA5ReadInput_get}
\label{Complexity.clA5ReadInput_get}
\leanok
\uses{Complexity.clA5ReadInput}
\difficulty{1}
For every count $m$, assignment $a$, and position $j < m$, the $j$-th bit of the recovered input word is exactly the assigned bit at variable $j$.
\end{lemma}

\begin{lemma}[Pinning recovers an exact-length certificate]
\lean{Complexity.clA5Certificate}
\label{Complexity.clA5Certificate}
\leanok
\uses{Complexity.clA5ReadInput, Complexity.clA5ReadInput_get}
\difficulty{4}
For all parameters $C$ and $e$, every instance $x$, and every assignment agreeing with $x$ on the first $|x|$ variables, there is a certificate $u$ of length exactly $C(|x|+1)^e$ such that the recovered input word of the first $|x| + C(|x|+1)^e$ variables is the concatenation of $x$ and $u$; the certificate polynomial is not enlarged and no bounded-length witness is substituted.
\end{lemma}

\begin{lemma}[Run output is the emission concatenation]
\lean{Complexity.clA5Run_output}
\label{Complexity.clA5Run_output}
\leanok
\uses{Complexity.emitted, Complexity.snapshotAt, Turing.FinTM, Turing.MultiTapeTM.initCfg, Turing.MultiTapeTM.outputSymbol, Turing.MultiTapeTM.runFrom, Turing.MultiTapeTM.runFrom_succ_eq_step', Turing.MultiTapeTM.step_output}
\difficulty{3}
For every machine, input $y$, and horizon $T$, the physical output after $T$ steps is the chronological concatenation of the optional emissions of the snapshots at times $0, \dots, T-1$, including the write on a halting transition.
\end{lemma}

\begin{lemma}[No false bit means no false emission]
\lean{Complexity.clA5NoFalse}
\label{Complexity.clA5NoFalse}
\leanok
\uses{Complexity.clA5Run_output, Complexity.emitted, Complexity.snapshotAt, Turing.FinTM, Turing.MultiTapeTM.initCfg, Turing.MultiTapeTM.runFrom}
\difficulty{4}
For every machine, input $y$, and horizon $T$, the append-only output after $T$ steps contains no false bit exactly when no snapshot before time $T$ emits the bit false; the equivalence is an exact word statement with no halting assumption.
\end{lemma}

\begin{lemma}[No false emission means acceptance]
\lean{Complexity.clA5Decider_accept}
\label{Complexity.clA5Decider_accept}
\leanok
\uses{Complexity.clA5NoFalse, Complexity.emitted, Complexity.snapshotAt, Turing.FinTM, Turing.FinTM.ComputesInTime, Turing.FinTM.computesInTime_iff, Turing.MultiTapeTM.indicator}
\difficulty{3}
Let a machine compute, on input $y$ within time $T$, the single indicator bit of membership of $y$ in a language $V$. Then no snapshot before time $T$ emits the bit false exactly when $y \in V$; silent or junk traces cannot satisfy this equivalence, because the total decider supplies its actual verdict bit.
\end{lemma}

\begin{lemma}[Snapshot blocks are nonempty]
\lean{Complexity.clA5Width_pos}
\label{Complexity.clA5Width_pos}
\leanok
\uses{Complexity.CLFieldCode, Complexity.clWidth, Turing.FinTM}
\difficulty{1}
For every machine and optional-state binary code, the snapshot block width is positive, since every product code contains at least the two input-symbol bits.
\end{lemma}

\begin{definition}[Assignment of a genuine run]
\lean{Complexity.clA5TraceAssignment}
\label{Complexity.clA5TraceAssignment}
\leanok
\uses{Complexity.CLFieldCode, Complexity.clA5Width_pos, Complexity.clBlockEncode, Complexity.clWidth, Complexity.snapshotAt, Turing.FinTM}
For a machine with a state code and an input word $y$, the trace assignment maps the first $|y|$ variables to the literal bits of $y$, and each later variable to the corresponding bit, located by quotient and remainder modulo the block width, of the raw encoding of the chronological snapshot of the run on $y$.
\end{definition}

\begin{lemma}[Trace assignments retain the input bits]
\lean{Complexity.clA5Trace_input}
\label{Complexity.clA5Trace_input}
\leanok
\uses{Complexity.CLFieldCode, Complexity.clA5TraceAssignment, Turing.FinTM}
\difficulty{1}
For every machine, state code, input word $y$, and position $j < |y|$, the trace assignment at variable $j$ is the $j$-th bit of $y$.
\end{lemma}

\begin{lemma}[Trace assignments encode every snapshot]
\lean{Complexity.clA5Trace_block}
\label{Complexity.clA5Trace_block}
\leanok
\uses{Complexity.CLFieldCode, Complexity.clA5Block, Complexity.clA5TraceAssignment, Complexity.clA5Width_pos, Complexity.clBlockEncode, Complexity.clPack, Complexity.clWidth, Complexity.snapshotAt, Turing.FinTM}
\difficulty{4}
For every machine, state code, input word $y$, and time $t$, the raw block of the trace assignment at time $t$ equals the encoding of the genuine snapshot of the run on $y$ at time $t$; quotient and remainder recover each complete block.
\end{lemma}

\begin{lemma}[Soundness of the tableau]
\lean{Complexity.clA5Sound}
\label{Complexity.clA5Sound}
\leanok
\uses{Complexity.CLFieldCode, Complexity.clA5Block, Complexity.clA5Certificate, Complexity.clA5Decider_accept, Complexity.clA5Initial_meaning, Complexity.clA5Input_meaning, Complexity.clA5ReadInput, Complexity.clA5ReadInput_get, Complexity.clA5Reconstruct, Complexity.clA5Tableau_eval, Complexity.clA5Work_meaning, Complexity.clBlockDecode, Complexity.clBlockDecode_encode, Complexity.clTableau, Complexity.emitted, Complexity.snapshotAt, Std.Sat.CNF.Satisfiable, Turing.FinTM, Turing.FinTM.ComputesInTime, Turing.FinTM.Oblivious, Turing.MultiTapeTM.indicator}
\difficulty{5}
Let an oblivious machine decide, within time $T$ and for every word $y$ of length $|x| + C(|x|+1)^e$, the indicator of membership of $y$ in a verifier language $V$. If the Cook--Levin tableau of the instance $x$ at horizon $T$ is satisfiable, then there is a certificate $u$ of length exactly $C(|x|+1)^e$ with the concatenation of $x$ and $u$ in $V$. The raw-block reconstruction excludes junk assignments before any decoder use, and totality plus the no-false family forces the accepting verdict.
\end{lemma}

\begin{lemma}[Completeness of the tableau]
\lean{Complexity.clA5Complete}
\label{Complexity.clA5Complete}
\leanok
\uses{Complexity.CLFieldCode, Complexity.clA5Block, Complexity.clA5Decider_accept, Complexity.clA5Initial_meaning, Complexity.clA5Input_meaning, Complexity.clA5Tableau_eval, Complexity.clA5TraceAssignment, Complexity.clA5Trace_block, Complexity.clA5Trace_input, Complexity.clA5Work_meaning, Complexity.clBlockDecode, Complexity.clBlockDecode_encode, Complexity.clBlockEncode, Complexity.clBlock_input, Complexity.clBlock_state, Complexity.clBlock_work, Complexity.clStateSlice, Complexity.clTableau, Complexity.emitted, Complexity.prevVisit, Complexity.snapshotAt, Complexity.snapshotAt_inputSymbol, Complexity.snapshotAt_state_succ, Complexity.snapshotAt_workSymbol, Complexity.stepState, Std.Sat.CNF.Satisfiable, Turing.FinTM, Turing.FinTM.ComputesInTime, Turing.FinTM.Oblivious, Turing.MultiTapeTM.indicator}
\difficulty{5}
Let an oblivious machine decide, within time $T$ and for every word $y$ of length $|x| + C(|x|+1)^e$, the indicator of membership of $y$ in a verifier language $V$. If some certificate $u$ of length exactly $C(|x|+1)^e$ makes the concatenation of $x$ and $u$ belong to $V$, then the Cook--Levin tableau of $x$ at horizon $T$ is satisfiable: the literal input assignment together with the raw encodings of the genuine snapshots satisfies every family, and the singleton true output supplies the no-false clauses.
\end{lemma}

\begin{lemma}[Equisatisfiability with the certificate relation]
\lean{Complexity.clA5Equisat}
\label{Complexity.clA5Equisat}
\leanok
\uses{Complexity.CLFieldCode, Complexity.clA5Complete, Complexity.clA5Sound, Complexity.clProducerHorizon, Complexity.clTableau, Std.Sat.CNF.Satisfiable, Turing.FinTM, Turing.FinTM.ComputesInTime, Turing.FinTM.DecidesInTime, Turing.FinTM.Oblivious, Turing.MultiTapeTM.indicator}
\difficulty{2}
Let an oblivious machine decide a verifier language $V$ within time $c(A(n+1)^d+1)^2$. For every instance $x$, the Cook--Levin tableau of $x$ at the producer horizon is satisfiable exactly when some certificate $u$ of length $C(|x|+1)^e$ makes the concatenation of $x$ and $u$ belong to $V$.
\end{lemma}

\begin{lemma}[The tableau serialization is a Karp reduction]
\lean{Complexity.clA5Reduction}
\label{Complexity.clA5Reduction}
\leanok
\uses{Complexity.NP, Complexity.SAT, Complexity.clA5Equisat, Complexity.clA5OutputIdentity, Complexity.clFieldCode, Complexity.clNPVerifier, Complexity.clProducerHorizon, Complexity.clTableau, Std.Sat.CNF.Satisfiable, Std.Sat.CNF.decode, Std.Sat.CNF.decode_serialize, Std.Sat.CNF.serialize}
\difficulty{3}
For every language $L$ in $\mathrm{NP}$, the map sending an instance $x$ to the serialization of its Cook--Levin tableau at the producer horizon is a polynomial-time many-one reduction from $L$ to $\mathrm{SAT}$; neither the certificate parameters nor any public statement is modified.
\end{lemma}

\begin{theorem}[Cook--Levin hardness of SAT]
\lean{Complexity.SAT_NPHard}
\label{Complexity.SAT_NPHard}
\leanok
\uses{Complexity.NPHard, Complexity.SAT, Complexity.clA5Reduction}
\difficulty{1}
The language $\mathrm{SAT}$ of serialized Boolean CNF formulas admitting a satisfying assignment is $\mathrm{NP}$-hard: every language in $\mathrm{NP}$ polynomial-time many-one reduces to $\mathrm{SAT}$.
\end{theorem}

\begin{theorem}[SAT is NP-complete]
\lean{Complexity.SAT_NPComplete}
\label{Complexity.SAT_NPComplete}
\leanok
\uses{Complexity.NPComplete, Complexity.SAT, Complexity.SAT_NPHard, Complexity.SAT_mem_NP}
\difficulty{1}
The language $\mathrm{SAT}$ of serialized satisfiable CNF formulas is $\mathrm{NP}$-complete: it belongs to $\mathrm{NP}$ and every language in $\mathrm{NP}$ polynomial-time many-one reduces to it.
\end{theorem}

\begin{theorem}[Hardness of 3SAT]
\lean{Complexity.SAT3_NPHard}
\label{Complexity.SAT3_NPHard}
\leanok
\uses{Complexity.NPHard, Complexity.NPHard.polyTimeReducible, Complexity.SAT3, Complexity.SAT_NPHard, Complexity.SAT_reducible_SAT3}
\difficulty{1}
The language $\mathrm{3SAT}$ of serialized satisfiable CNF formulas with at most three literals per clause is $\mathrm{NP}$-hard: every language in $\mathrm{NP}$ polynomial-time many-one reduces to it.
\end{theorem}

\begin{theorem}[3SAT is NP-complete]
\lean{Complexity.SAT3_NPComplete}
\label{Complexity.SAT3_NPComplete}
\leanok
\uses{Complexity.NPComplete, Complexity.SAT3, Complexity.SAT3_NPHard, Complexity.SAT3_mem_NP}
\difficulty{1}
The language $\mathrm{3SAT}$ of serialized satisfiable CNF formulas with at most three literals per clause is $\mathrm{NP}$-complete: it belongs to $\mathrm{NP}$ and every language in $\mathrm{NP}$ polynomial-time many-one reduces to it.
\end{theorem}
